% ============================================================
%  Harald Høffding, "Danske Filosofer"
%  (Copenhagen: Gyldendalske Boghandel, Nordisk Forlag, 1909; 206 pp.)
%  ENGLISH TRANSLATION — SKELETON
%
%  Translation-only project: a clean modern e-book edition exists
%  (SAGA Egmont / Lindhardt & Ringhof, ISBN 9788726363128), so no
%  Danish transcription is produced. Base texts:
%    * e-book (17 chapters, clean digital text)
%    * 1909 first-edition scan (Royal Danish Library) for anything the
%      e-book omits — NOTABLY the untitled opening Introduction on
%      printed pp. 1–2, which the modern e-book drops entirely.
%
%  Chapter numbering is reproduced exactly as printed in 1909: the first
%  chapter is numbered "I." (roman), the rest "2."–"17." (arabic). This
%  quirk is in both the original and the e-book; it is preserved, not
%  normalised.
% ============================================================
\documentclass[12pt,a4paper]{book}
\usepackage[utf8]{inputenc}
\usepackage[T1]{fontenc}
\usepackage{libertinus}
\usepackage{libertinust1math}
\usepackage{textalpha}   % occasional Greek carried verbatim from the Danish
\usepackage[english]{babel}
\usepackage[protrusion=true,expansion=true]{microtype}
\usepackage{setspace}
\usepackage[margin=1.2in]{geometry}
\usepackage{fancyhdr}
\setlength{\headheight}{15pt}
\usepackage{hyperref}

\hypersetup{
  pdftitle={Danish Philosophers},
  pdfauthor={Harald Høffding},
  hidelinks
}

\pagestyle{fancy}
\fancyhf{}
\fancyhead[LE,RO]{\thepage}
\fancyhead[RE]{\textit{Danish Philosophers}}
\fancyhead[LO]{\textit{\leftmark}}

\setstretch{1.3}

\title{\textbf{Danish Philosophers}}
\author{Harald Høffding\\[6pt]
  \normalsize translated from \textit{Danske Filosofer} (1909)}
\date{København: Gyldendalske Boghandel, Nordisk Forlag, 1909}

% --- helper: unnumbered chapter that still enters the ToC and running head
\newcommand{\dfchapter}[1]{%
  \chapter*{#1}%
  \addcontentsline{toc}{chapter}{#1}%
  \markboth{#1}{#1}%
}

\begin{document}
\maketitle
\thispagestyle{empty}
\newpage

\tableofcontents
\newpage

% ============================================================
%  INTRODUCTION  (printed pp. 1–2, untitled in 1909; OMITTED by the e-book —
%  translate from the Royal Library scan)
%  Opening words: "Et lille Land som vort ligger aabent for Indflydelser
%  udefra …" ending "… At dansk Tænkning har sin Ejendommelighed, kan ses
%  ved Sammenligning med svensk Tænkning."
% ============================================================
\dfchapter{Introduction}
\label{ch:df-intro}

\noindent A small country such as ours lies open to influences from without.
The spiritual currents from the great centres of culture reach decisively
into its intellectual life, and this life cannot attain that firm continuity
through the ages which becomes possible when a country is able to bring forth
a great multiplicity of thinking minds --- minds that never let development
come to a standstill and that, through the various types they present, can
break the influences coming from abroad, so that the continuity in a people's
intellectual life may be the better preserved.

Nevertheless, serious and independent thought has been devoted here at home to
the questions that it is philosophy's calling to discuss --- the
presuppositions of science, the significance of its results for one's view of
life, the foundation of ethical evaluation, the nature and laws of the human
spirit, the relation between religion and humanity. The influences from
without have been assimilated in a distinctive way, and our national character
has shown itself in the manner in which the problems have been treated. In the
treatment of such questions as those mentioned, which lie on the borders
between strict science and life, the personalities of the thinkers must
necessarily contribute in a higher degree than in other inquiries. And where
personality operates according to its whole distinctiveness, the national
traits too assert themselves. For these are not layers or pieces joined from
without to the other qualities of individuals, but determine those qualities
from the very first and sound along with them, every time the strings of the
mind are struck.

That Danish thought has its own distinctive character can be seen by
comparison with Swedish thought. Sweden's intellectual life has received the
same currents from abroad as Denmark's, but in the domain of thought these
have been assimilated and developed in another way and in another direction
than with us. Speculation and mysticism have played a far greater role in
Sweden than in Denmark, and, bound up with this, an urge toward systematic
completeness has asserted itself there which has never found firm ground with
us. Danish thought has been essentially borne by a psychological and ethical
interest, and again and again a level-headed criticism has made itself felt
against the construction of systems.\footnote{See my essay on Swedish and
Danish philosophy in the collection \emph{Norden} (1902), as well as my essay
\emph{Filosofien i Sverige} [Philosophy in Sweden] (\emph{Det Letterstedtske
Tidsskrift}, 1879). --- On the philosophical movements abroad that come to
influence Danish philosophy, fuller information will be found in my work
\emph{Den nyere Filosofis Historie} [The History of Modern Philosophy] by
means of the index. --- Biographies I do not as a rule give, chiefly because I
can refer to the excellent articles by my colleague Professor Kroman in the
\emph{Biografisk Leksikon}.} Among our foremost thinking minds (Holberg, the
Ørsteds, Sibbern, Kierkegaard) this distinctive character makes itself felt,
and it gives rise to thoughts and works which, though they have not entered
into the great international discussion of the questions of knowledge and of
life, nonetheless have their value --- and that not for us alone. A
characterization of the Danish philosophers can teach us a great deal by way
of orientation toward those questions, and it will often appear that our
thinkers might very well have made their mark on a larger stage, had
circumstances assigned them one. The heaven of thought is, like the physical
heaven, accessible to all, not only to those born among a great people. This
book will therefore be able to cast light on the philosophical problems
themselves, and not merely serve to characterize Danish intellectual life as
it expresses itself in the domain of thought.

\begin{center}---\end{center}

% ============================================================
%  I. LUDVIG HOLBERG   (Da.: "I. LUDVIG HOLBERG")
% ============================================================
\dfchapter{I. Ludvig Holberg}
\label{ch:df-holberg}

\noindent Danish philosophy begins with the very founder of our literature. He
is the first here at home to take an independent stand on the problems of
thought, at the same time as he brings the ideas and attempted solutions that
had arisen abroad into Danish intellectual life.

Holberg no doubt stands first and foremost for us as a poet, but it would be
wrong to see him solely from this point of view and to interpret his moods amid
the vicissitudes of his life merely in the light of the fate his poetic work
met with. Inquiry held a central interest for him and filled an essential part
of his life, from his young days to his advanced old age. He held, as does his
Erasmus Montanus, that studies are an ornament in the days of good fortune and
a consolation in those of misfortune. Nothing befell him so heavily that it was
not made lighter by study --- as he himself said. It is characteristic that a
similar remark is attributed to Leibniz, who likewise had a many-sided interest
in reading and inquiry. But Leibniz had a brighter and lighter view of life and
of men than Holberg, whose heavy nature did not let him rest content with the
crystal-clear prism into which Leibniz had ground his view of life. He could not
be led by speculative considerations to overlook the disharmonies of life or to
be reconciled with them. For Holberg the evil in the world was a hardship not so
easily broken. Human life suffered in particular under a great opposition
between appearance and reality. ``The shadow is taken for the body'' --- this
expression recurs again and again in Holberg, and he wrote his comedies, his
\emph{Niels Klim}, and his \emph{Moralske Tanker} [Moral Reflections] in order
to counteract this perpetual confusion. He does not wish to preach morality, but
to tear illusions apart and awaken men to self-knowledge.

When he thus ascribes to himself an express moral purpose in all his works,
this is to be understood only as an interpretation he afterwards placed upon his
authorship. There was of course at work in him an involuntary urge to unfold his
intellectual power in the fields of poetry, historical research, and thought,
and only when in his old age he looked back on his rich and manifold activity
did he find, in a moral aim, a somewhat artificial and philistine
explanation.\footnote{Compare my essay on Holberg in \emph{Mindre Arbejder}
[Minor Works], Second Series, pp.~150--153.} But his work of thought was a side
of his activity that he looked back upon with particular pleasure. Just as he
thought he could ``say without scruple that he is almost the first Dane who, in
place of meagre and confused chronicles, has given us tolerably well-ordered
histories,'' so too he prided himself on having ``set the study of morals on its
feet again, which right up to this time had lain as if buried here in the
North.''

But Holberg was not merely a moral philosopher. He occupied himself with the new
science and the new view of the world that had developed in the age of the
Renaissance, and whose most conspicuous and at the same time most far-reaching
results were the changed position of the Earth in the picture of the world and
the unboundedness of the world. And besides this, his thought was taken up with
the great oppositions of the world and of life, which culminated in the problem
of evil --- a problem that in Holberg's very youth was being discussed by such
great thinkers as Bayle and Leibniz.

His thought loved to set oppositions sharply against one another, so that the
problem might spring forth all the more from the strong collision --- the
opposition, for example, between the infinite horizon that science had opened up
and human life with its striving and its wishes; the oppositions within human
nature; the conflict of good and evil. In treating questions of this kind he was
fond of ``paradoxical opinions,'' by which he understood such as ``with good
reasons overturn the opinions that have gained citizenship.'' It was not a bent
for hair-splitting or for playing with ideas that guided him here: the paradoxes
had to have ``probability'' --- that is, to be really well-founded --- and this
they were when it was widened experience that led to them. He saw that when one
places oneself on the broad ground of experience --- when one bases one's
judgments on experience --- one comes to doubt much that had hitherto stood as
self-evident. It is no wonder that one arrives at fixed opinions when one holds
to the given traditions and lets them determine what one experiences and how one
interprets what one has experienced.

It is on this point above all that Holberg's significance as a thinker lies ---
in the weight he places on as wide an experience as possible and on the problems
that experience presents at any given time. Both these tendencies make a
conclusive systematization difficult, if not impossible. When Holberg, in the
preface to \emph{Moralske Tanker}, speaks of completing his system,\footnote{``In
view of my age and infirmity, I regard myself as a man ready to set out on a
journey, who must draw up his \emph{Systema} before he takes his leave and sets
off.''} he surely meant thereby only a rounding-off or conclusion so far as his
own person was concerned. Even the most open and acute thinker halts at a
boundary, but he need not regard this boundary as the boundary of all thought
and experience.

With this it agrees that when Holberg is to say to which philosophers he owes the
most, he names John Locke, the chief representative of empiricist philosophy, and
Pierre Bayle, the great critic and raiser of problems. On the whole he owes most
to English and French intellectual life. Nowhere on his travels abroad did he
feel so much at home as in Oxford and Paris. He rejoiced that he had come to
study at the French and English academies, and he called England in particular
``the true school of philosophy.''

When Holberg, after his travels, after the experience he had gathered, and after
his thought had been sharpened by his persistent occupation with the problems of
the age, attempted to philosophize in the Danish academic world, great
difficulties arose for him from various quarters.

The Danish language was but little fashioned for more advanced thinking; it had
narrow limits to its diffusion --- and within those limits it was not respected
by its own people. ``My name,'' says Holberg, ``might perhaps have been tolerably
well known abroad, had I not written in the Danish language, which is so confined
that often in Denmark itself one must search for someone who understands it.''
The greater gratitude must we feel toward the founder of our literature that he
did not despise this language, but did what he could to cultivate it into an
instrument for thought. Had he, for the sake of his fame, written in a foreign
tongue, a Danish life of thought would not have developed. What he began was
carried on by his successors, especially by Eilschow, Sneedorff, and Kraft.

But in this still little-developed language one could not express oneself so
freely as in the other countries, especially in England. ``Here in our North,''
says Holberg, ``we are plagued with sharp censorships, by which all gaiety and
spirit are damped in a writer.'' When a book is furnished with an ``Imprimatur,''
it is usually a sign that there is nothing to it. It was above all from the
theological side that the censorship was prompted; Holberg noticed this
especially with \emph{Niels Klim}. ``I took particular care that the clergy
should not get occasion to array themselves in armour against me, since I know
from experience that these people are the most implacable.'' ``It is only those
adversaries I fear who clothe themselves in the armour of religion.'' And in his
last years he grew even more cautious: ``I put my words and writings on screws
more than before, since I have noticed that at one time they were being examined
more closely.'' --- Holberg's religious opinions were not orthodox. Of that there
is no doubt, though it may be difficult to decide where earnestness leaves off in
his utterances and where caution begins. People have exaggerated both the
earnestness and the caution --- have found both more orthodoxy and more tactics
in him than there are. He had his theology, which he held to be reasonable and
consonant with his experience. It goes not much further than Rousseau's and
Kant's, but he thought that, by right interpretation, it could be said to be
contained in the biblical and dogmatic conceptions.

Most of the students (two-thirds) at the University of Copenhagen in Holberg's
time studied theology, and the scholastic philosophy still cultivated here at
home was regarded as an important means of guarding the true doctrine. By formal
logic thought was to be trained in distinctions and syllogisms so as to be able
to set dogmatics in system. In \emph{Erasmus Montanus} Holberg presents a young
dialectical fencer of the kind that the constant disputation exercises were bound
to produce. His formal virtuosity drives him to undertake to prove whatsoever it
may be. ``That is why,'' says Erasmus Montanus, ``I have studied, and that is why
I have racked my brains, so that I can say what I will and defend it before the
whole world.'' In the scenes where Erasmus disputes with the folk on the Hill,
however, Holberg does not let him use the logic that was taught at the
university. Erasmus is fond of proceeding by affirmative syllogisms in the second
figure,\footnote{A is B and C is B, therefore A is C! Nille cannot fly, a stone
cannot fly, therefore Nille is a stone!} which scholasticism knew very well to be
invalid. But Holberg employs such sophisms as Aristophanes and Plato had already
ascribed to their opponents. Instead of showing the uselessness of formal
quibbling, he has achieved a comic effect through the way in which the respect
that the folk on the Hill nonetheless had for Copenhagen learning caused sheer
nonsense to stand as proven truth.

Among the subjects to which logical acuteness is applied, Erasmus names first and
foremost the theological ones --- for example, the question whether the angels
were created before men. But such questions are named too as whether the Earth is
round or oval, and concerning the distance and size of the sun and stars; and it
is on a question of this latter kind that Erasmus is brought down by the
home-grown Inquisition. Now it is characteristic that questions of this latter
kind interested Holberg himself in a high degree and led him to some of his
boldest thoughts, as will be shown in what follows. Here he had ``the Hill'' and
``the learned men of the Hill'' against himself, and he would have come into a
conflict similar to that of Erasmus Montanus, had he here openly drawn the
consequences. But in the comedy his main concern was to set useful studies over
against mere sword-play, and so he was ruthless enough to let Erasmus fall
precisely where he was in the right. The Lieutenant in the fifth act, the play's
moralist, has, like Holberg, chiefly studied ``Latin authors, natural law, and
moral matters,'' which Montanus regards as trifling. The moral of the play is
expressed in the Lieutenant's words: ``The first commandment in philosophy is to
know oneself\ldots\ The more one does, the more there seems to remain to learn.
But you make philosophy into an art of fencing.''

Instead of the subtleties or the broodings over the nature of God and the angels
that were then called metaphysics, Holberg wished to turn to subjects that could
yield a positive return for philosophizing. Instead of countless uncertain things
he wanted a few certain ones. --- Let us then bring forward the chief questions
that especially occupied Holberg's thought.

\begin{center}---\end{center}

Often in his later years Holberg declares that moral philosophy is his real
study. The chief work here is \emph{Moralske Tanker} (1744). It is not a
systematic book, but a series of little essays that occasionally stray into less
scientific subjects, for there is talk not only of poetry and of studies, but
also of the reasons why the author has not married, and of the interesting
question whether a bad stomach is as bad as a bad wife. But the main thought, to
which he returns again and again, is that there exists a morality founded on the
light of nature. In his \emph{Epistler} [Epistles] too, the other chief source of
his philosophy, he maintains this proposition. ``Natural morals are one study,
theological morals another. When one undertakes to write of the former, it is
required that one found the doctrine on the light of nature.'' But natural
morality is not, as was claimed, ``a heathen morality.'' ``Christ's morality''
too is founded on the light of nature; the difference is merely that the same
``morality'' is carried out in greater perfection.

Holberg's ethics is religious in so far as he connects it with ``natural
religion.'' But he will have morality placed before the special dogmas and
mysteries. ``One must, through Christ's morality and a sound philosophy, lay the
foundation before the dogmas and mysteries of Christianity are taught.''
``Children must be made into human beings before they become Christians.'' A human
confirmation must precede the Christian one; for if a person learns theology
before he learns to become a human being, he never becomes a human being. ``One
begins by crying: Believe! before one shows what one ought to believe and not to
believe. That is to go to work like that judge who began with the execution and
afterwards had the case examined.'' One ought to impart to men a critical
faculty, that they may not take the shadow for the body.

For Holberg the foundation of the Reformation was ``the Christian freedom to
search before one believes.'' ``It is in a person's power to search, but not to
believe.'' Error after careful examination is more excusable than orthodoxy
without preceding inquiry. Strife is better than brutal concord. Holberg himself
would rather read heretics than the orthodox, because the latter answer more with
abuse than with reasons. And it was his conviction that one had to give up the
outworks of the doctrine of faith in order to defend the main fortress the
better. For his own part he was willing to subscribe to all the fundamental
articles of faith found in the Confession; yet he would rather reduce than
increase their number, for there is need of nowhere near so many articles of
faith as are commonly alleged to be quite necessary. --- Unfortunately Holberg
does not say which ones he would strike out.

He will not have inquiry brought to a halt by one's taking refuge in ``God's
motives and reserved causes'' or ``the imperfection of reason and God's
inscrutable ways.'' That is a sign of laziness. Such an appeal is permissible only
when one has striven with all one's might to find out the cause and has sought to
reconcile the difficulties that arise with reason.

Nevertheless, as already indicated, there is in Holberg's moral catechism
\emph{(Epistle 46)} both a religious and a rational element. It is founded on
natural religion. The first three articles treat of God and immortality. The
fourth article demands that no doctrine be accepted before it is made
demonstrable; for otherwise one may fall into false belief as readily as into
true, and reason, God's most excellent gift, would be of no use. The fifth
article demands the rejection of all doctrines that conflict with God's holy
attributes, and the sixth runs thus: Hate and persecute no one who against his
will goes astray in his faith!

Of the two elements the rational lies deepest, and logically the fourth article
ought to have been placed before the three preceding ones. Very clear is a
statement from \emph{Moralske Tanker} that the right method would be to explain to
the young the chief commandments of nature, to set forth what fairness, justice,
and goodness are --- that they are virtues in God's creatures and \emph{therefore}
must consistently be attributes of the Godhead itself. --- Thus the attributes
ascribed to God must first be known from the human world. It is then a circular
inference when one derives morality from theology, grounding human justice by the
divine. The ethical, according to Holberg, takes priority over the religious; the
latter is understood by means of the former, before the reverse can become the
case.

The ``light of nature'' Holberg finds revealed in certain principles that for him
are immediately evident. ``The great commandment of nature'': Do not to others
what you would not have them do to you! --- will, he thinks, be granted by
everyone who hears it explained --- even though Locke is right that it cannot be
said to be innate. It is as clear as that a thing cannot at one and the same time
be and not be. --- Morality thus has its own source, even though it stands in
close connection with natural religion. Whether that principle is really
self-evident is another question. The problem of how we come to set up the
principles on which all ethical demonstration rests, Holberg does not enter into.

As regards the individual qualities and actions, Holberg impresses upon us that
all perfection depends on a definite limitation: ``All virtue consists in
mediocrity, and as soon as it passes beyond the bounds, it is metamorphosed into
vice.'' ``The Danish proverb, Fear God and follow the highway! is quite well
founded; for those who are always looking for byways and wish to shorten the
roads most often arrive last. And it is of such crafty wayfarers that people here
in this country say that they go by Per Gante's shortcut.'' --- Unfortunately
there is perhaps something Danish in these rules. But one would gravely
misunderstand Holberg if one supposed that he recommends a middle-of-the-road
morality. He does so as little as Aristotle did in his doctrine of virtue as the
right mean.\footnote{See on this my \emph{Etik} [Ethics] XI, 9 (cf.\ IV, 2).
(Third edition, p.~221 f.; p.~89).} What is too much or too little is determined
by the individual's nature and may differ from individual to individual. What for
one is a byway may be a highway for another. It was not Holberg's meaning to
exclude new roads. His comparisons and examples show clearly what he means. In
case of sickness one patient is to have one pill, another two or three. A choleric
man must be judged otherwise than one born with a ``temperate complexion.'' One
who has no gall cannot be praised for coolness of temper. To ride on a tame ass is
no merit, but to govern a wild and mettlesome horse deserves applause and clapping
of hands. In the chastisement of misdeeds the punishment ought to be fitted to the
malefactor's constitution; a misdeed may be equally great and yet not equally
punishable in different persons.

It is thus an individualization of morality that Holberg demands. In ethical
evaluations and demands, account is to be taken of the inner personal conditions
of individuals, and all cannot be cut by one measure. Holberg does not, it is
true, go deeply into the problem that here comes to light, namely how far regard
for the individual dispositions of particular persons can be reconciled with the
demands that the existence of society makes. But the problem of the individual
here steps forth for the first time in Danish thought. --- On closer examination
it will prove to be bound up with the question of the demonstrability and
practicability of ethical principles.\footnote{Compare my essay on the Law of
Relation in Ethics (\emph{Etiske Undersøgelser} [Ethical Investigations], 1891,
pp.~42--83).}

\begin{center}---\end{center}

Holberg not only felt an urge to think about ethical questions. He is at the same
time the first here at home to have brooded over the influence that the new
conception of the world --- founded by Copernicus, Galileo, Descartes, and Newton
--- was bound to have on one's very view of life. He emphasizes here the
enlargement that the picture of the world has received through the new discoveries
and ways of seeing. The new astronomy has widened our conceptions and given us a
greater idea of God. More than ever, speculations about God now stand as
unreasonable. He declares himself in agreement with the Greek Simonides that the
more one thinks about such questions, the more difficult they become. It is of no
use to say, with ``common theology,'' that Simonides speaks out of heathen
blindness, but that revelation has brought enlightenment; for the figurative
expressions that the Bible predominantly uses give no proper concept, indeed lead
to ``coarse and carnal conceptions.'' Holberg finds that those versed in nature
have higher thoughts of the divine Being than the common crowd, indeed than many
of the clergy themselves, in that they do not, like the latter, make attempts to
express the incomprehensible. The deeper we enter into the investigation of what
eternity, infinity, omniscience, omnipresence (``all-being''), and omnipotence
are, the more we drown in them and see ever less.

According to the new picture of the world, moreover, the Earth stands as a grain
of sand over against the infinite universe; it is no longer the centre of the
world, and even the sun about which it revolves is only one star among countless
other stars. Men are like worms that creep about on a little hillock and yet
imagine that everything exists for their sake. From this side too it appears how
hopeless metaphysical and theological speculations are.

In the world of natural science and of philosophy there still stood, in Holberg's
time, the dispute between Descartes's and Newton's teachings concerning the
motions of the heavenly bodies and the laws underlying them. Descartes assumed
that the universe consisted of systems, each of which was in whirling motion. The
globes belonging to the same vortex were driven round by the motion of the
surrounding mass. That the globes did not, by virtue of the law of inertia, fly
off into the universe away from the centre of their system was explained by the
fact that the surrounding mass, in which they as it were swam, continually carried
them round the centre. Newton rejected this vortex theory and explained the
globes' remaining in their orbits by gravity or attraction, which constantly
counteracted the tendency to fly off in a straight line. Holberg seems to have
been much taken up with this discussion. ``According to the cursory knowledge he
has of it,'' he prefers to hold to the Cartesian conception, because it appears to
him simple and intuitable, whereas Newton appears to him to appeal to a mysterious
force when he speaks of attraction.

But if the universe consists of many worlds in constant whirling motion, how does
this square with the biblical account of the creation of the world? This point had
in Holberg's younger days caused him much doubt and unease. He settled the matter
for himself, however, in this way: that the Mosaic creation-story tells only of
the coming-into-being of the vortex-system to which the Earth belongs. There is,
after all, nothing said about its having been the first or the only world that was
created! Perhaps there takes place a perpetual arising and passing-away of worlds!

It is the boldest thought Holberg hints at. And with this thought there is given a
great background for Holberg's view of the world, which his own caution and the
pettiness of the age prevented him from bringing forth more definitely, but which
is nonetheless felt in the freedom of spirit that his thought on the whole bears
witness to. His conceptions of the decisive questions unfolded under the influence
of a wide horizon, which would have come forth more clearly had he not himself,
too much in the less happy sense of the words, ``feared God and followed the
highway.''

Yet one point still remains as that in which Holberg deviated in the most decisive
way from the orthodox way of thinking --- namely his position on the problem of
evil, a problem that shortly before his appearance had been discussed by Bayle and
Leibniz. Pierre Bayle had maintained the impossibility of explaining the origin of
evil if one started from the conception of an all-good and almighty God as the
author of all things --- and had declared that the only consistent assumption,
when one took one's stand on the standpoint of natural reason with revelation left
aside, was the presupposition of two original principles, a good and an evil, in
constant strife and struggle, as the Manichaeans had taught. Holberg, however,
could as little accept this dualistic conception as he could follow Leibniz in his
optimistic speculations. But one of his friends, whom he designates by B. S. (it
is supposed to be Rector Bernhard Schnabel in Roskilde), had shown him a way out:
if God from the beginning stands over against a chaotic matter that is not
perfectly receptive to the forms he wishes to impress upon it, but offers a
certain resistance, then it is understood that not only the perfect exists, but
also the imperfect. One then avoids the Manichaean assumption of two active
principles, a good and an evil God, and holds fast to the belief in a good God
whose carrying-out of his intentions is hampered by a recalcitrant matter. Such a
view (of a kind similar to that in which Rousseau and Voltaire, and later Stuart
Mill and William James, ended up) would --- Holberg thinks --- have disarmed
Bayle, had revelation not taught a creation out of nothing. On this point his
accommodation and apparent submission to Church doctrine is quite plain. With his
teaching about the origin of our world-system as the successor of an earlier one,
the solution of the problem of evil here indicated would have fitted excellently.

For the rest, Holberg has yet another argument against Bayle. The different beings
must be measured by different measures. A human being's perfection is other than a
fly's and than an angel's; every being is perfect in its kind. This line of
thought accords well with Holberg's individualizing of ethics, and it presupposes
that the measure is taken from the very nature of the individual beings, each for
itself. It raises a problem similar to the individualization of ethics, and in
sharper form.

Here, as at several other points, Holberg has developed his thoughts neither fully
nor consistently. But he has done a foundational work and made free movement of
thought possible here at home. For the first time, great problems were discussed
in his writings in Danish and in a Danish spirit.

\begin{center}---\end{center}

% ============================================================
%  2. SNEEDORFF AND THE WOLFFIANS   (Da.: "2. SNEEDORFF OG WOLFIANERNE")
% ============================================================
\dfchapter{2. Sneedorff and the Wolffians}
\label{ch:df-sneedorff}

\noindent Holberg's figure fills so much of Danish intellectual life in the
eighteenth century that there seems to be an empty space after him. In
independent thinking, no one at the close of the century came near him. But he
had awakened a desire to think and had shown the possibility of using the Danish
language for it. Among the best of the younger men there stirred a sense of shame
that there did not yet exist a Danish scientific literature, while other European
peoples had long since begun to use their mother tongue in science.

The one who, as a thinker, worked most nearly in Holberg's spirit was Jens
Schielderup Sneedorff. Like Holberg he was influenced by English and French
literature. Shaftesbury and Montesquieu in particular exerted an influence on
him. He was born in 1724 and studied at Göttingen, whose university was visited
by many foreigners and where philosophy, political science, and natural science
flourished. Later he served (1751--1761) as professor of law and political
science at the newly founded Academy at Sorø, which in its brief flowering became
the cradle of Danish scientific literature. He died as early as 1764.

Sneedorff praises Holberg's great services to our language and our literature.
``But we should not be worthy of having such a fellow citizen if, out of laziness
or baseness, we were to believe that there was nothing for us to do.'' The
struggle for the mother tongue and for true, natural thinking goes hand in hand
in Sneedorff. Away with barbarous technical terms and academic sword-play ---
with dead memory and dead abstraction!

According to the historians of language, the foundation of the more modern Danish
language was laid by Sneedorff and his contemporaries. ---

Sneedorff's thought concerns partly matters of law and the state, partly general
philosophical points of view.

In the first respect the chief work is \emph{Den borgerlige Regering} [Civil
Government] (1757). The weight is here laid on the principle that a good
government must follow principles that build on experience and accord with the
nature of the people; only such a government can win the people's trust. Written
laws, not least constitutions, are unnecessary --- indeed they may be harmful on
account of their unalterability, which hinders the constant interplay with
experience. The Danish absolute monarchy was in Sneedorff's eyes the only
absolute monarchy that was grounded on the people's trust, and he protests
therefore against its being called despotism. --- But in order that experiences
may be made, and that the nature of the people may reveal itself, public opinion
must unfold freely, and to this again belong enlightenment as well as freedom in
all civil relations for all estates, not least for the hitherto so neglected
peasantry. Freedom makes men active and reasonable, and without freedom there is
no honour! ---

Absolute monarchy is here recognized because it is enlightened and popular, no
longer because it is by the grace of God. A turn in thought and in justification
that was bound to have far-reaching consequences!

Sneedorff was one of those who, in the domain of general ideas, worked to prepare
the great agrarian reforms in Denmark at the close of the eighteenth century.
Whereas Holberg had still held that the peasant must be kept to work by
compulsory laws, Sneedorff and those of like mind maintained that both regard for
personal freedom and for the common good --- considerations that at bottom
coincide --- demanded changes in the peasants' condition.\footnote{Compare E.\
Holm: \emph{Om det Syn paa Kongemagt, Folk og borgerlig Frihed, der udviklede sig
i den dansk-norske Stat i Midten af det attende Aarhundrede} [On the View of
Royal Power, People, and Civil Liberty that Developed in the Dano-Norwegian State
in the Middle of the Eighteenth Century] (University Programme, 1883),
pp.~107--114.}

Sneedorff's general philosophical standpoint already reveals itself in his
doctrine of law and the state, but he develops it more closely, partly in
polemical forms, in his periodical \emph{Den patriotiske Tilskuer} [The Patriotic
Spectator] (especially in the first volume, 1761).

Where it is a matter of experience and the description of experience, or where it
is a matter of finding expression for the life of feeling, there, according to
Sneedorff, one cannot apply strict definitions and methods of proof. What
essentially matters is to procure as full a content as possible, and a logical or
mathematical formulation would here only have a narrowing effect --- indeed, at
bottom, rest on an illusion. On the ground of experience we reach only
probabilities (``likelihoods''), not necessities. And we must be cautious about
finding contradictions, for these may fall away with widened experience, which
can teach us that there are far more possibilities than we imagined.

It is characteristic that Sneedorff, in his discussion of the relation of
experience to logical contradictions, emphasizes the opposite side from the one
on which Holberg dwelt. Holberg laid weight on the point that the greater the
experience we gain, the more contradictions and difficulties we shall find. It
was his business to find and pose problems. Sneedorff dwells especially on the
working of widened experience to make the provisionally found contradictions fall
away. This is bound up with the brighter tone that lies over his philosophizing,
and with his hope in progress, but also with a weakening in the sense for the
sharpness of problems --- of the ever-recurring problems.

But Sneedorff is not dogmatic. Just as he has a clear eye for the hypothetical
character of the propositions we draw from experience --- the chief source of our
knowledge --- so too he has an eye for the role that analogy plays in the
border-region of thought, indeed already when we wish to infer what goes on in
the inner life of other beings. And he maintains that only individual things are
real, while species and generic designations are due to our thinking. There is
expressed in this last proposition a sense for the individual, which we have
already met with in Holberg and shall find again in later Danish thinkers
(Treschow, A.\ S.\ Ørsted, Kierkegaard). Sneedorff protests against philosophy's
being metaphysics and logic alone. The study of the mind's operations and
movements demands fineness and loftiness of thought no less than pure
metaphysics, and jurisprudence and aesthetics now claim their right, after
theology, philology, mathematics, and metaphysics have so long had the say.

\begin{center}---\end{center}

Holberg and Sneedorff belong to the kind of thinkers who strive to build on
experience, with full awareness that by this path no absolute conclusion can be
reached. They take a stand against the thinking that starts from certain
dogmatically established principles, from which one then proceeds by strict
inferences. Experience here receives a subordinate significance. For a large
part, Sneedorff's expositions are directed against Jens Kraft (1720--1765),
professor of mathematics at Sorø and author of a whole series of philosophical
textbooks. Kraft had, in his \emph{Kritiske Breve til Videnskabernes Fremvæxt og
Smagens Forbedring} [Critical Letters on the Growth of the Sciences and the
Improvement of Taste] (1761), defended the mathematical (deductive) method in
science. He emphasized the indispensability of abstract thinking and of
mathematics if knowledge is to be thorough and if it is to be applied in
practice, and he fears that the preponderance of ``the fine sciences'' will lead
to indolence and superficial cleverness. On Sneedorff's emphasizing of the method
of experience, however, he does not enter.

Some years before (1751--1753), Kraft had published a systematic presentation of
metaphysics after the method of the German philosopher Wolff. By metaphysics he
understands ``the doctrine of the first and most necessary fundamental teachings
needed for human knowledge.'' The first part is ontology, the doctrine of the
most universal truths, which can be reduced to two: that the self-contradictory
is impossible, and that everything has its cause. The second part is cosmology,
the doctrine ``of the forces and primary physical properties of natural things,''
the key to physics. Here he parts from Wolff by adopting Newton's hypothesis and
rejecting Descartes's; the book thereby marks a turning point in Danish science.
The third part is psychology. Here Wolff distinguished between rational
(speculative) and experimental (empirical) psychology. Kraft fuses the two
together, remarking: ``Among all the theoretical discoveries about the soul I
have found but few of importance, and I have therefore held that one might with
propriety and right introduce them together with the experimental in a single
work.'' Here, then, Kraft himself wished to let experience come into its own. The
fourth and last part is ``the natural doctrine of God, or natural theology.''
Here God's existence is proved partly by the argument that there must be a first
beginning to the changes in the world, since otherwise these would not have their
sufficient ground, and the question of the origin would otherwise ever recur
without being answered; partly by the argument that it lies in the concept of God
as the infinite and perfect Being that he must exist, since he has his ground in
himself and is not dependent on anything else that could exclude his existence.

Kraft's work is the only completed system written in Danish. Apart from the
individual points singled out, it rests entirely on Wolff's philosophy, the first
German system. Chr.\ Wolff was a disciple of Leibniz, and his philosophy is at
once a popularization and a systematization of the master's thoughts. He
supplanted the old scholasticism at the German universities, just as Descartes
had supplanted it in France and Holland, and Locke in England. But his
theological colleagues at Halle feared the consequences of the new philosophy,
partly because it so zealously maintained ``the principle of sufficient reason,''
which seemed to them to lead to fatalism, partly because it taught the
possibility of a purely human morality. In Denmark too it was ill regarded while
Pietism held sway. One of the first Danes to be influenced by it, Henrik Stampe,
is said for that reason to have lodged ``half incognito, at least as regards
court chaplain Bluhme,'' in Wolff's house. Wolffianism was first introduced at
the University of Copenhagen by jurists (Stampe, Kofod Anker) and theologians
(Gunnerus). But the first purely philosophical writer here at home of this
tendency was Frederik Christian Eilschow (1725--1750).

Just as Wolff in Germany created a philosophical terminology to replace the Latin
one, and for the most part did so, so too it was one of Eilschow's endeavours to
replace the many foreign scientific words with Danish ones, and these had then
first to be coined. Eilschow has often been felicitous in his proposals in this
respect. Words such as \emph{Enkelthed} (singularity), \emph{Sandsynlighed}
(probability), \emph{Sædelighed} (morality), \emph{Personlighed} (personality),
\emph{Selvbevidsthed} (self-consciousness), \emph{Sindsbevægelse} (emotion), and
\emph{Overflade} (surface) are due to him. He maintains that the sciences ought
to be presented in the people's own language. Against those who held that there
was no reason to abandon Latin, since practical people had no time for reading of
a scientific kind anyway, he appeals to his own experience. There are many who,
without belonging to the learned, buy and read scientific books whose language
they understand; many hours, indeed whole nights, they spend absorbed in
them.\footnote{Eilschow: \emph{Cogitationes de scientiis vernacula lingva
docendis}, Hafniæ 1747, p.~7.}

Eilschow praises the English because they lay such great weight on experience,
and on account of the many profound thoughts that come from England. But he
attaches himself predominantly to Wolff, whose philosophy he wishes to make more
accessible, also to ladies,\footnote{\emph{Forsøg til en Fruentimmerphilosophie}
[Attempt at a Philosophy for Ladies]. Copenhagen 1749.} and whose work he also
wishes to continue.\footnote{\emph{Philosophiske Breve over adskillige nyttige og
vigtige Ting} [Philosophical Letters on Various Useful and Important Things].
Copenhagen 1748.} Wolff had, indeed, set metaphysics in system, dividing it into
fundamental doctrine (ontology), the doctrine of the world (cosmology), the
doctrine of the soul (psychology), and natural theology. Eilschow did not himself
treat all these subjects; such a treatment was, as already mentioned, given by
Kraft some years later. It is natural theology that Eilschow especially occupies
himself with. It is from the concept of possibility that Eilschow will conduct his
proof of God's existence. No finite thing's possibility has its ground in itself;
but all possibilities have their ultimate ground in God's possibility --- and
this has its ground in itself, just as all numbers presuppose the number one.
This is, however, not to be understood as though the possibilities of things were
due to God's will. They stand independently of his will in his eternal thought;
only their realization is due to the will. Against the eternal possibilities the
will can do nothing --- the divine will as little as the human.

Now the source of evil --- of pain and of sin --- lies in the limitation to which
everything finite as such is subject, and which follows from the fact that finite
beings stand in connection with one another. This connection, which in itself is
a good, brings with it barriers and therewith the possibility of imperfection.
Evil cannot then be charged to God's account, since he has as little created the
possibilities as he has created his own understanding.\footnote{\emph{Philosophiske
Breve}, p.~150. --- A similar line of thought is found in Jens Kraft's
\emph{Naturlige Theologie}, pp.~13--16. --- Eilschow states, however, that we
cannot infer from the imperfections found in one part of the world to the
imperfection of the whole.}

As little as Leibniz and Wolff has Eilschow answered the objection that Bayle
raised and that Holberg repeated: why, then, has God created a world of which he
knew that it contained the possibility of pain and sin, a possibility that neither
divine nor human will could annul? Holberg is more consistent than Eilschow.
Understanding and will, possibility and actuality, are here in such a degree at
war with one another that they must really belong to two different beings. The God
who is willing but not able is not almighty; he has a barrier --- whether one
now, with Holberg, seeks it in an eternal matter, or, with the Manichaeans, in an
eternal evil being, an Ahriman. --- Eilschow lacks Holberg's sense for the
seriousness of problems.\footnote{A purely philosophical standpoint on the problem
of evil is given in my \emph{Religionsfilosofi} [Philosophy of Religion],
§§~88--89, and my \emph{Etik} [Ethics], chap.~VI.}

On the question of the relation between reason and revelation, Eilschow likewise
makes use of the distinction between possibility and actuality. God can reveal
something that one cannot derive from reason and from the possibilities it
establishes; but he cannot reveal anything that overthrows natural knowledge,
contradicts the clearest reports of the senses, and upsets all that on which the
certainty of all truths rests.

Just as Sneedorff recognized absolute monarchy because it worked in the service of
enlightenment and of the popular cause, and not because it sprang from the grace
of God, so Eilschow and the other Wolffians recognized revelation in the
confidence that it does not contradict reason and the eternal possibilities. The
question was bound to become --- and did become --- on which side the wrong lies
when a contradiction actually shows itself. There came the time, toward the close
of the century, when men were in no doubt that a contradiction was present, nor in
doubt on which side the wrong lay.

For the time being, Eilschow was sure that he had in hand the means to refute the
freethinkers, that is, those who deny the divine providence.\footnote{The word
freethinker comes from England (free-thinker) and first occurs around the year
1700, used of and by men such as Toland and Collins. Earlier, from the sixteenth
century on, designations such as deists and naturalists were used in France.} They
sin against metaphysics and are refuted by means of the aforementioned doctrine of
the possibilities and their ultimate ground. Without metaphysics one easily goes
astray, precisely in such questions as the origin of evil and providence. Yet
Eilschow thinks that the freethinkers really merely wish for freedom to sin, and
that they deny the truth in order to escape obeying it. There are, to be sure,
some --- such as Spinoza --- who have led a fine life; but that is inconsistent,
since they expect no reward or punishment after death. ---

More kindly disposed than toward the freethinkers is Eilschow toward the animals.
He is one of the first philosophers of the modern age to maintain that we have
duties toward animals. In order for one to have a duty toward a being, it is not
necessary that it have reason, provided only that it has the capacity to suffer.
There exists a community between them and us, in that we take their service and in
return have duties toward them. ---

Wolff's systematic philosophy held sway at the university until into the new
century. Its last representative was Børge Riisbrigh (born 1731, died 1809), who
was professor of philosophy from 1767 to 1803. His teaching is praised for its
formal clarity, but he was not very productive. The most considerable work from
his hand is a treatise (from 1802--03) on Kant's philosophy, a work that does him
honour and that will be discussed later on.

\begin{center}---\end{center}

% ============================================================
%  3. DANISH KANTIANS AND THEIR OPPONENTS
%     (Da.: "3. DANSKE KANTIANERE OG DERES MODSTANDERE")
% ============================================================
\dfchapter{3. Danish Kantians and Their Opponents}
\label{ch:df-kantianere}

\noindent The Wolffian system was in Germany supplanted by one of the most
profound works of thought that the history of philosophy knows --- by Kant's
critical philosophy. Kant had been aroused by Hume's critique of the causal
principle, the principle on which all natural science and all speculative
philosophy and theology built. At the same time Newton's mathematically developed
natural science stood before him as the ideal of scientific knowledge. When he
now asked how such science was possible despite Hume's demonstration that the
causal principle was neither self-evident nor demonstrable, he found the solution
in this: that there are certain forms with which the human mind, according to its
nature, constantly works, and of which it must therefore make use when it is to
establish what has objective validity --- what is truth and reality. Thereby it
became at once possible to acknowledge science in its independence and to see its
limits; these must lie where the human mind's forms of intuition and thought no
longer find any given material to be applied to. Science cannot therefore go
further than the possibility of experience goes. When this possibility fails, a
person can --- and must --- have a personal faith to support his holding fast to
the command of conscience; but a speculative theology is impossible.

The law of conscience was, according to Kant, a pure law of reason: just as the
human mind in its science strives to bring observations under the forms and laws
of intuition and understanding --- under the forms of space and time, of the
concept of magnitude and the concept of cause --- so it can convince itself of the
morality of an action only when it can show that the action can be acknowledged on
the basis of a universal law valid for all; the maxim I follow in my action must
be such that everyone in my place could follow it.

As Kant, from developed science, seeks back to the innermost nature of thought
itself at work, so on the ethical-religious domain he sought back to the
self-legislation of conscience. From the principles he thus gained, he judged the
speculative systems and the theological dogmas. His great sense for the original
forces led him further, on the aesthetic domain, to point to genius as the
involuntarily welling source of art. In thought, in conscience, and in genius
there expresses itself, for Kant, one and the same fundamental force in different
forms. By this great idea Kant at once forms the conclusion of the eighteenth
century's intellectual labour and points far beyond its horizon.

It was a young, promising scholar who especially aroused attention for Kant's
philosophy in Denmark. In the spring of 1793 Christian Hornemann (born 1769)
delivered lectures on Kant's critique of reason. They were cut short by the
author's death the same year, and the great hopes that had been placed in him fell
to the ground. With great enthusiasm he had taken Kant's principal ideas into
himself, and a period of study with Reinhold at Jena had developed his
philosophical interest. In his lectures (which are printed in his \emph{Efterladte
Skrifter} [Posthumous Writings], 1795) he gave a clear account of the relation of
the critical philosophy to earlier conceptions and then began to go through Kant's
principal work; but he got no further than the end of the doctrine of space and
time. --- Hornemann's main task seems to have been to show how the religious
problem would present itself under the presupposition of the new conception.

The purely theoretical side of Kant's philosophy was but little brought forward
here at home, and just as little the aesthetic. It was the moral- and
religious-philosophical side that was especially attended to. The most
considerable champion of Kant's ethics and doctrine of right was the man who later
became Denmark's greatest jurist, Anders Sandøe Ørsted (born 1778); in 1798 he won
the university's gold medal for a prize essay \emph{Om Sammenhængen mellem Dyds- og
Retslærens Princip} [On the Connection between the Principle of the Doctrine of
Virtue and that of the Doctrine of Right], an excellent presentation for its time
of Kant's ethics and philosophy of right, and an account of the discussions to
which this gave rise among contemporaries, both abroad and here at home. In the
standing dispute between the ``Bible party'' (the Palaeologists) and the ``Reason
party'' (the Neologists) --- the first represented by Bishop Balle, the second by
Horrebow (editor of the periodical \emph{Jesus og Fornuften} [Jesus and Reason])
--- Ørsted took part with a longer essay (in \emph{Philosophisk Repertorium},
1797--1798), directed against Horrebow's popular half-philosophy. From a standpoint
similar to Lessing's in \emph{Erziehung des Menschengeschlechts} and Fichte's in
\emph{Kritik aller Offenbarung}, he points beyond the two contending parties. ---
Ørsted did not --- as we shall see in a later connection --- remain at the strictly
Kantian standpoint he had taken in these essays. But admiration and gratitude
toward his great teacher followed him all his life, and in his advanced old age he
expressed it in a beautiful way.

It was, as Ørsted remarks in his autobiography,\footnote{\emph{Af mit Livs og min
Tids Historie} [From the History of My Life and My Times], I (1851), p.~22. --- A
rather remarkable and independent position toward Kant's ethics was taken by the
then professor of moral philosophy Anders Gamborg, who, following Hutcheson and
Adam Smith, maintained a morality independent of religion and springing from love
of humanity. In his first work (\emph{Forskel mellem Dyd og gode Handlinger} [The
Difference between Virtue and Good Actions], 1783) he does not yet know Kant's
writings. In a later work (\emph{Jesu Moral} [The Morality of Jesus], 1799) he
declares that he is neither a eudaemonist nor a Kantian; yet he is glad that in his
insistence on the unconditional character of the duty of truthfulness he agrees
with Kant, though their grounds differ.} not so much the young theologians as the
young jurists who occupied themselves with philosophy. And in the questions that
occupied the last decade of the eighteenth century, the vigorous and free spirit
that characterizes Kant's and Fichte's ethical and legal-philosophical views could
have worked extraordinarily fruitfully. But there did not arise here, as in
Germany, a series of independent and energetic thinkers and reformers such as those
to whom the great reforms in Prussia after 1806 were due. The ``Königsberg
idealism,'' to which, according to Gneist, Prussia's restoration was for so large a
part due, set no lasting effects in the life of the people here at home, and was
early superseded by the more passive and dreaming tendency of Romanticism. This
came partly from external political circumstances; partly, too, no doubt, from the
fact that twelve years before the close of the century we had carried through one
of the greatest reforms our country's history knows --- a reform in which,
moreover, as touched on in its place, philosophical ideas had also played a part.
Besides Ørsted, whose philosophical and juridical studies in youth belonged to the
preparation for his long activity as judge and legislator, one must nevertheless
name a Kantian whose writings, by their clarity and energy, stand far above most
of what was produced here at home in the public discussion that overflowed all its
banks in the last decade of the eighteenth century, before the strict press
regulation of 1799 stopped the free discussion of public affairs. The clergyman
Michael Gottlieb Birckner (born 1756, died 1798) maintained, in his work \emph{Om
Trykkefriheden og dens Love} [On Freedom of the Press and Its Laws] (1797), that
even in an absolute monarchy complete freedom of the press may very well prevail.
--- It was the ideal of enlightened absolutism --- free discussion of all thoughts,
and yet absolute obedience in all actions --- that was here carried through to its
uttermost consequence. Kant himself had, moreover, in his remarkable little essay
\emph{Was ist Aufklärung?} (1784), stated that only an enlightened autocrat, not a
free state, dared say: Reason as much as you will --- but obey! But it was then
also Kant's opinion that such an autocrat the world had had only in Friedrich II of
Prussia. Birckner's Kantianism had led him to continue a line of thought in which
Sneedorff and his generation had already moved. In other domains too Birckner
appears as a pronounced Kantian. Even more strongly than Kant, and without the
sense for the sublime that marks the master's presentation, he asserts the sole
right of reason in the ethical life. All feeling --- with the exception of the pure
feeling of respect that the moral law inspires --- was for him really self-love.
Even love of fatherland and love of humanity he would therefore not acknowledge as
the proper ethical motives; they cannot be required and commanded, but must develop
involuntarily, and it is not the immediate respect for the right that makes itself
known in them, but the individual's personal captivation.\footnote{See his essays
\emph{Om Kierlighed til Fædrelandet} [On Love of the Fatherland] and \emph{Om
Christendommens højeste Moralprincip} [On the Highest Moral Principle of
Christianity] (\emph{Efterladte Skrifter} [Posthumous Writings], ed.\ A.\ S.\
Ørsted, Copenhagen 1800). --- On the ethical and psychological misunderstandings in
Kant's and Birckner's conception, see my \emph{Etik} [Ethics] XII, 1.} ---

Among the Danish Kantians, Baggesen must also be named. In him it was not the
strict, rational side of Kant's philosophy that aroused interest, but precisely
that character of sublimity which Kant's thoughts can present, especially when they
move on the borderlines of experience. Kant satisfied for a time the young poet's
rapturous searching, though he admitted that the purely theoretical enthusiasm was
foreign to him. He philosophized, as he himself says in a letter to
Reinhold,\footnote{\emph{Aus Jens Baggesens Briefwechsel mit Reinhold und Jacobi},
I, p.~6.} out of ethical and religious interest, and his work of thought really
always aimed at procuring symbols for what moved in his feeling and his
imagination. Philosophy was for him ``Gedankenspiel'' [a play of thoughts]:
``Since I do not work in philosophy and on philosophy, my philosophizing ---
whether it consist in the reading of philosophical works or in conversation with
philosophical friends --- is only play, but the dearest of all plays.''\footnote{\emph{Op.\
cit.}, II, p.~56. --- It is of interest to hear Reinhold's judgment of Baggesen,
made to a third party: ``Sein darstellender Genius ist sein Plagegeist. In den
Augenblicken, wo er ihn in Ruhe lässt, ist Baggesen im Denken Philosoph und im
Handeln verständiger Mann von der ersten Grösse --- aber das sind auch nur
Augenblicke --- sonst ist er Engel oder Teufel, und befindet sich im Himmel oder in
der Hölle --- leider! öfters da als dort. Als Korrespondent (wenn er schreibt) und
als auf einige Tage besuchender Gast hat er wohl Niemanden über sich.'' [His
representative genius is his tormenting spirit. In the moments when it leaves him in
peace, Baggesen is in thought a philosopher and in action a sensible man of the
first rank --- but those too are only moments --- otherwise he is angel or devil,
and finds himself in heaven or in hell --- alas, more often there than here. As a
correspondent (when he writes) and as a guest visiting for a few days he has, I
think, no one above him.] Reinhold to Erhard, 2 Aug.\ 1796.
(\emph{Denkwürdigkeiten des Philosophen und Arztes Erhard}, ed.\ Varnhagen von
Ense, 1830, p.~426).} Poetry and philosophy strove over this strange spirit. And he
stood not only on the border between two domains of the spirit, but also on the
border between two centuries. We shall meet him later in the remarkable collision
with Grundtvig, which promised so much but came to so little, because in Baggesen
the play with symbols gained the upper hand over the seriousness of thought.

\begin{center}---\end{center}

Kantianism was met with no small resistance here at home from various sides, not
least from the side of self-satisfied ``Enlightenment.'' And literary spokesmen
such as Rahbek were vexed at the young Kantians' arrogance and the new technical
terms they threw about. ``Enlightenment'' was, moreover, agreed with Orthodoxy
that it could not be so difficult as Kant thought to get the riddles of existence
solved. Why all the many preparations, when one had one's Bible and one's catechism
or one's ``sound common sense'' to resort to?

But the new tendency of thought also found really philosophical opponents, and some
of these have even with clarity pointed out defects and unsolved questions in the
critical philosophy. Tyge Rothe and Johannes Boye collided with criticism in the
course of their work at developing their own ideas, while Riisbrigh, Niels Schow,
and Treschow directly examined its foundation and range.

Tyge Rothe's (born 1731, died 1795) authorship falls into two periods, a
historical-political and a philosophical. To the first period belong his writings
on the influence of Christianity on the condition of the peoples in Europe and on
the history of the North, which gave him occasion to work for the emancipation of
the peasantry. To the second period belongs his chief philosophical work
\emph{Philosophies Ideer til Kundskab om vor Art og til Glæde over den}
[Philosophy's Ideas toward Knowledge of Our Kind and Joy in It] (1788--1789). It is
here not, as has been said and repeated, one of the ordinary spokesmen for
Enlightenment and the doctrine of felicity who has the word. Rothe is influenced by
men such as Herder and Jacobi. He was thinking of laying down his pen when Herder's
\emph{Ideen} appeared, but yet took courage to walk in his footsteps. He is clear
that neither man's nobility nor his happiness can be judged merely by his
enlightenment. Free humanity, working its way toward inner harmony and toward
sympathy with all that is human, is for him man's vocation. Social life develops
out of the free and undisturbed urge of hearts, not out of need and misery, nor out
of pure reason. And like his two German masters he finds in immediate feeling an
organ for the highest: by a revelation of nature, not through handed-down articles
of faith or by the way of pure reason, we discover our own being as caused by an
eternal being, and therewith the causal law itself is given in its inevitability.
All philosophy is therefore theology. The doctrines of positive religion are
interpreted according to the doctrines of nature's revelation, which experience too
confirms; Newton's demonstration of the lawfulness of nature leads back to nature's
springing from the will of a single being; he has given us the physico-theology
that so gloriously characterizes our enlightened age.

Kant's teaching belongs (like Spinoza's, Berkeley's, and Hume's) to ``the dark
philosophical systems.'' It shakes the empiricist philosophy on which alone we can
build with security. --- It is evidently partly Kant's discussion of the causal
principle as not immediately given or evident, partly his critique of the proofs of
God's existence, that repels Rothe. He is sure that when Kant has closely
determined how far strict knowledge goes, then men will feel that they cannot be
content with it, and will work their way out of the dark labyrinth into the clear
daylight of immediate assurance.

Rothe represents a standpoint that otherwise does not come forth here at home at the
close of the eighteenth century. He lacks Herder's originality and Jacobi's
remarkable union of criticism and pathos, and his sentimentality and artificiality
contrast strongly with their assertion of the right of the involuntary and the
immediate in the life of the spirit. But it was not without significance that a
tone of this kind too, weak and sometimes false though it was, came to sound among
the many voices of the age here at home. ---

Rothe had received influence not only from German writers, but also from French
literature. It was especially the Genevan Charles Bonnet, the naturalist and
philosopher,\footnote{Bonnet was among the foreigners supported by the Danish
government under Frederik V. His still instructive psychological work \emph{Essai
analytique sur les facultés de l'âme} (of particular interest for the psychology of
attention and of the will) was printed in Copenhagen in 1760 and dedicated to
Frederik V. Though a republican, he praised the Danes as fortunate in having such a
monarch: ``Ce republicain envieroit le sort de l'heureux Danois, si un citoyen de
Genève pouvoit envier quelque chose!'' [This republican would envy the lot of the
happy Dane, if a citizen of Geneva could envy anything!]} who inspired him and was
his guide in the contemplation of nature. He embraces Bonnet's doctrine of
preformation, according to which each species was created for itself and in such a
way that the germs of all later individuals were contained in the first
representatives of the species. All beings form one great chain, but the links in
this chain of nature remain through the ages specifically (``according to kind'')
distinct, and no transformation of the species or transition from one species to
another takes place. The original dispositions in nature cannot be annulled by any
mechanical force. Man is perfect in his first disposition, but external
circumstances, great upheavals of the Earth, have intervened inhibitingly and
altered the disposition, so that a long and laborious development became necessary.
At every stage of development, however, it holds that a being attains the
satisfaction it is possible for it to reach. Through this consideration the harmony
of existence asserts itself for Rothe's thought. He expressly rejects the way out
that Eilschow and Kraft (following Leibniz and Wolff) had seized upon: to
distinguish between the possibilities for the divine understanding and the actuality
as it springs from the divine will. ``Only one could God's plan be; that plan was
the only one that could be put into effect, for away with the foolish talk that God
pondered and weighed and chose!'' ``What could be in God's realm of creation, that
is there. What could subsist, that subsists.'' Why influences (or, as Tyge Rothe
says, ``impacts'') hinder the unfolding of the dispositions, of that, according to
Rothe, one can give no explanation. For our sake the eternal laws are not altered;
nor are they given for our sake. But this thought does not darken Rothe's bright
view of things. --- Optimism is rising --- from Holberg through Eilschow and Kraft
to Tyge Rothe. ---

\medskip

Whereas Rothe conceives development as merely the unfolding of original germs,
Johannes Boye (born 1756, died 1830) stands as a philosopher of development in the
more modern sense of the word. In his remarkable work \emph{Statens Ven} [The Friend
of the State] (1792 ff.) he lays down the thought that everything endowed with
force, man included, must struggle with resistance and violence to reach its goal.
No one sleeps himself good, but every step forward toward perfection must be won by
struggle. Only through strife and labour is man helped forward to what he is. A
struggle of all against all is waged. Bravery was a virtue before piety became one,
and physical superiority prevailed before spiritual superiority could be exercised.
Nature brings up its children strictly --- but thereby it brings up heroes. The lazy
and the weak are pushed aside in that strife of forces through which order and
harmony are won. This blessed strife is the means that Providence uses for the
realization of the world-plan.

Right and morality depend on the development of culture, and it is the annals of
history, not speculative vapouring, that can teach us something about their origin.
Man has developed from an animal state and from brutality, through all degrees of
culture, to the height he now occupies. And the state is the most important aid to
culture. The life of the state is grounded on the understanding of the necessity of
renouncing some of one's rights in order to preserve the others. Through social life
and by the aid of language man has ceased to be an animal; without them we should
still be living among the animals in the forest and eating acorns with them. But
development does not halt at a self-limitation in expectation of reciprocity. The
fundamental urge in us is first satisfied in sympathy and noble self-feeling. Man
then feels himself as part of a whole; his thought gazes out over wider circles, and
his interest is extended far beyond his own self. At this standpoint virtue is its
own reward. Even self-sacrifice becomes possible only through the urge to use the
great force led by sympathy. ---

From this standpoint Boye combats the critical philosophy's attempt to build
morality and right on pure reason. This philosophy halts, he says, at an
unreconciled opposition between pure reason and the psychological motives. ``A good
will'' must, after all, be a will that is good toward someone, and should love of
humanity be reprehensible because in it man follows a deep urge? The ideas of the
good and the right develop out of human interests and their conflict. ---

Boye's line of thought, here given in condensed form, does not come into its own in
his books and is therefore easily misunderstood by superficial readers. It is a
manly and free view that stirs in him, and it had not deserved the oblivion that has
become its lot, but which is explained by the spirit of the age, which was first
facilely optimistic and soon after rapturously Romantic. Yet Boye lacked sufficient
historical sense and sufficient material to give his ideas a thorough foundation.

A.\ S.\ Ørsted, who in his prize essay answered Boye's attack on Kant's ethics,
acknowledges his work as ``spirited and tasteful.'' He agrees with Boye that the
most probable hypothesis is that humanity in its childhood did not stand high above
the animals --- an assumption that Kant too had put forward (in his essay
\emph{Muthmaaslicher Anfang des Menschengeschlechts}, 1786). But he impresses the
difference between the historical development of knowledge or of conscience and the
philosophical demonstration of the presuppositions on which all knowledge and all
conscience build --- or between the beginning of knowledge and its foundation. ---
Ørsted is right that Kant's opponents constantly overlooked (and in part still
overlook) this difference; but Kant has undeniably himself contributed much to the
misunderstanding, especially by not giving the psychological and historical inquiry
its due. ---

\medskip

Børge Riisbrigh found in his advanced old age the time and the inclination to
discuss the philosophy that supplanted the doctrine he had proclaimed throughout
his long university activity. And he does it with an acuteness and an impartiality
that do him honour. His inquiry appeared in the Transactions of the Society of
Sciences in 1803,\footnote{Its title is: ``Undersøgelse om og hvorvidt deres Mening
er grundet, som holde for, at sand Filosofi og et fuldstændigt Begreb om sand
Filosofi først i vore Tider ere blevne til'' [Inquiry whether, and how far, the
opinion is well founded of those who hold that true philosophy and a complete
concept of true philosophy have first arisen in our times]. --- Philosophy was not
among the subjects represented in the Danish Royal Society of Sciences from its
founding (1742). Only under A.\ P.\ Bernstorff's presidency were (1795) ``the
properly philosophical sciences, or the so-called speculative philosophy, which in
recent times has acquired so much importance,'' added to the mathematical, physical,
and historical subjects. C.\ Molbech: \emph{Det kongelige danske Videnskabernes
Selskabs Historie i dets første Aarhundrede}, p.~283.} the same year in which he
withdrew from his teaching activity, and it is still of interest. He shows first
that already before Kant men had been aware that what our knowledge presupposes and
applies in its investigation of objects cannot, at the very moment of application,
have come from the objects themselves, but must have been in the mind beforehand,
whether in its first origin it stems from experience or not. There is, then, a
``pre-experiential knowledge'' (knowledge \emph{a priori}). But Kant first sought
the source of these principles in the human mind itself and showed their purely
formal character, and gave a catalogue of them. He thereby laid the foundation of a
new science within philosophy, the scientific critique of reason. But he did not,
like some of his successors, assert that before him men had neither true philosophy
nor knew what philosophy was. And he by no means spurned experience; he assumed both
an experiential and a pre-experiential philosophy.

Riisbrigh wishes, as a counterpart to Kant's inquiry into pure reason, an inquiry
into empirical reason, and he thinks that one would then come to see its great
significance. One would then again appreciate what Socrates, Aristotle, Bacon, and
Locke had taught. On behalf of science, indeed of all human culture, one must warn
against undervaluing the science of experience. Our inquiry must, after all, occupy
itself with objects other than its own forms. And on the other hand these forms are
themselves something given: the mind has as little given itself its essential forms
as it has produced the forms in the world of sense! ---

\medskip

In the same volume of the Transactions of the Society of Sciences that brought
Riisbrigh's essay, there is an essay by the philologist and archaeologist Niels
Schow (born 1754, died 1830) that likewise presents no small
interest.\footnote{Its title is: ``Den nuherskende Skepticisme eller Kritik i
Filosofien kan, naar den grunder sig paa den menneskelige Fornuft og rigtigt
anvendes, ej være farlig, hverken for videnskabelig Cultur eller Moralitet'' [The
now-prevailing skepticism or criticism in philosophy cannot, when it is grounded on
human reason and rightly applied, be dangerous, either to scientific culture or to
morality].} The author finds Kant's significance especially in the movement in
thought that he had aroused; whereas he finds the critique that G.\ E.\ Schulze had
directed at Kant's philosophy (especially in \emph{Kritik der theoretischen
Philosophie}, 1801) altogether decisive. It is especially the critique of Kant's
strange concept ``Ding an sich'' [thing in itself] that is at issue. Schow maintains
that whether the absolute reality whose existence we know is contained in the
subject or in the object, or in what relation it stands to both, we do not know; we
only trace its effects. It exceeds the limit of human reason to give a positive
determination. Even the most sublime philosophical presentations --- like Spinoza's
--- are here anthropomorphisms. The task of philosophy is to determine the limits to
which our knowing is ever subject --- to discuss the results of the science of
experience in their significance for the ultimate questions --- and to work toward a
harmony between acting and knowing. Beyond the hypothetical we do not come.

\medskip

Already several years before Riisbrigh's and Schow's essays appeared, Niels Treschow
(born 1751, died 1833), in Kristiania, where he was then headmaster, had delivered
lectures on Kant's philosophy\footnote{\emph{Forelæsninger over den Kant'ske
Philosophie} [Lectures on the Kantian Philosophy]. Copenhagen 1798.} and given a
critical assessment of it that went deeper than any of the hitherto mentioned works
on this subject. In a clear and lively presentation Treschow gives a survey of
Kant's philosophy and attaches to it his critical remarks. He finds the derivation
of the pure forms of the understanding unsatisfactory in Kant, and shows in
particular that Kant could not rightly appeal to the current logic's division of
judgments when he himself alters this. And even if Kant had succeeded in finding all
the forms that knowledge makes use of, it would still have to be shown what it is in
the given, in the ``material,'' that makes possible and occasions the application of
a definite particular form in the individual particular case. What makes it, for
example, that I see here a quadrangle, there a triangle? It cannot be derived from
space as a general form of intuition. What makes it that I apply here the concept of
magnitude, there the concept of cause? It cannot be derived from the understanding's
general power to combine observations in certain forms (categories). There must in
the things themselves lie something at the ground of the particular modes of
combination. Already our sensations must teach us what nature itself has combined
and what it has separated.

There is lacking, according to Treschow, in Kant's presentation of the forms of
thought (the categories), such important concepts as likeness (identity) and
difference, agreement and conflict, which are nevertheless applied in all thinking,
and which are just as much applied to the given as the concepts of magnitude and of
cause. And the said concepts are even of quite special significance: the
understanding's work really consists merely in the comparison of what is given in
experience. It is according to the likeness or unlikeness of things, their agreement
or conflict, that the understanding decides which combinations among them have
validity and which not. --- Treschow has here pointed to an important point in the
theory of knowledge. Kant had not seen that the formal logical principles, just as
much as the mathematical principles and the causal principle, are applied to the
given material, so that their validity too must be examined by the critique of
reason.

It is the standpoint of empiricist philosophy that Treschow maintains against Kant.
He praises the Britons, who here not only have themselves come nearest to the truth
and made considerable progress in knowledge of human nature, but have also become
teachers for others. He upholds John Locke's standpoint against Leibniz and Kant,
and reproaches the Germans for not having been willing to learn from the English.
Locke was right that every form presupposes a material, and that no concepts can be
formed when certain representations are not given beforehand.

But Treschow misunderstands Kant when he says that, according to Kant, the
understanding forms its concepts ``before any experience comes to it.'' It is,
indeed, as A.\ S.\ Ørsted showed against Boye, Kant's express teaching that
knowledge \emph{begins with} experience, even if not all knowledge \emph{grounds
itself on} experience. Experience itself was, according to Kant, composed of given
material (the sensations) and involuntarily applied forms (forms of intuition and
categories).

With the limitation that Kant establishes for our knowledge, Treschow is likewise
not satisfied, as little as with the way in which Kant will have that urge satisfied
which science cannot satisfy. ``Kant himself cannot deny that there lies in our
nature an indescribable longing to know more than he holds it possible to attain ---
indeed, he confesses, what is still more, that the aim of our existence necessarily
demands a knowledge that we are able to attain by none of the indicated ways [those
of sense-observation and of understanding]. Does he then intend to bring us back
under the yoke of blind faith and the hierarchy, or is it deliberately that he, as
in a heroic poem, entangles the knot in order to drive attention to the highest and
earn the greater admiration by loosing it?'' --- Treschow despairs of Kant's art and
will not build on the new foundation that Kant wished to lay for one's view of life.
It is questionable to give practical reason and faith the right to overstep the
limits of knowledge. How can it be proved that certain articles of faith indicate
the only way in which the ethical ideas can be maintained? And if one gives up all
demonstration and appeals merely to personal interest, then one opens the door to
disputes of faith and persecutions. About that which is only to have practical
significance, everyone --- and not least those in power --- will think they have a
voice. --- Treschow's misgivings are not unfounded. Another question is whether Kant
is not right that what is decisive for every view of life is in the end a personal
urge arising from the experience of life. With this every faith, whether it supports
itself on science or not, must be able to accord. The great thing in Kant is
precisely the art with which he tightens the opposition between the power of thought
and the urge of life. He only did not see that the urge of life cannot once for all
be shown to have found its satisfaction in the dogmas of ``natural religion,'' any
more than in those of any other religion. ---

Treschow, whom we have here met only as a critic, we shall later meet as a
distinctive and independent thinker who in several directions exercised a
considerable influence on Danish intellectual life.

\begin{center}---\end{center}

% ============================================================
%  4. HENRIK STEFFENS   (Da.: "4. HENRIK STEFFENS")
% ============================================================
\dfchapter{4. Henrik Steffens}
\label{ch:df-steffens}

\noindent It seems to be a strife that recurs at the close of every century, when
the new century begins. In Weimar there was no doubt as to when the nineteenth
century began; it was celebrated by Goethe, Schiller, Schelling, and Steffens when
the clock struck twelve on the evening of 31 December 1800. This little circle
characterizes the transition between the two centuries. The eighteenth century had
been the age of criticism and rationalism. But men such as Rousseau, Lessing, and
Herder, Kant, Goethe, and Schiller pointed beyond their time by their sense for the
original, the natural, the genial --- for the deep-lying forces whose essence is
never wholly exhausted in the forms that an age or a century is able to give it.
These men represent the work and the struggle of the new humanistic Renaissance.
And just as high as they stand above the eighteenth century's common rationalism,
so high do they stand above the so-called Romanticism that marks the nineteenth
century's first reaction against its predecessor. It was the ideas of the said men
that were now, in the philosophy and poetry of Romanticism, carried out into the
forms of speculative thought and rapturous imagination. Romanticism builds on their
intellectual work, at the same time as it marks a certain reaction against it. Thus
the two youngest participants in that New Year's celebration, Schelling and
Steffens, mark a reaction against Kant's manly assertion of the independence and
self-limitation of thought, at the same time as they wish to build on the
foundation he had laid.

While France and England contended for world domination, in the North men immersed
themselves in the world of ideas and imagination. Schiller had struck the note in
his poem at the beginning of the new century:

\begin{verse}
Freiheit ist nur in dem Reich der Träume,\\
und das Schöne blüht nur im Gesang.
\end{verse}

\noindent [``Freedom is only in the realm of dreams, and the beautiful blooms only
in song.'']

As regards Denmark, there had in the last decades of the eighteenth century been a
lively discussion of public affairs, and a keen sense of political freedom had
developed. Sneedorff, Tyge Rothe, Birckner, and Johannes Boye bear witness to it.
But the hope that Birckner especially had entertained, that absolute monarchy and
absolute freedom of the press could be united, was dashed, and the general European
reaction penetrated also across our country's borders. It was a pressure from
Russia's side that gave the decisive impulse to the appearance of the press
ordinance of 1799, which stopped the free discussion.\footnote{See on this: Edvard
Holm: \emph{Den offentlige Mening og Statsmagten i den dansk-norske Stat i
Slutningen af det attende Aarhundrede} [Public Opinion and the Power of the State in
the Dano-Norwegian State at the Close of the Eighteenth Century] (University
Programme, 1888), p.~175 f.} Remarkably quickly minds settled to rest\footnote{The
Norwegians took the matter up in a more energetic way than the Danes, and they also
soon got a great political task. ``For Norway the seed that had been strewn at the
close of the eighteenth century acquired power to germinate, when the thought of the
necessity of a free constitution came to coincide with the question of leading an
independent political life. Most of the more considerable men of Eidsvoll had, as
young people in Copenhagen in the [seventeen-]nineties, received their first
awakening impressions in a political respect.'' E.\ Holm, \emph{op.\ cit.}, p.~183
f. --- Grundtvig once stated, in a conversation with Sibbern, that the movement in
the last decade of the eighteenth century was not allowed ``to run the life out of
itself.'' It rose again a generation later. (Sibbern: \emph{Dikaiosyne}, 1843,
p.~9.)} --- partly in the world of philistinism, partly in that of Romanticism.
Those who sought beyond the torment of the day, or sought consolation for the
misfortunes that soon broke in over the country, found what they sought in the world
of speculation and imagination.

And yet --- for the Danes, who had had neither a Kant nor a Goethe, it was a great
spiritual treasure that was opened with the new spiritual current called
Romanticism. Our literature, our life of thought, and our poetry were reborn a
second time. The time after Holberg had, as regards originality, been a trough
between waves. Now it went upward again, and here as in Holberg's time the new was
taken up and worked over in an independent way. The people set its own stamp on what
it appropriated.

It was Henrik Steffens who first proclaimed the gospel of Romanticism in Denmark.
Born in Stavanger (1773), and thus a Norwegian like Holberg and Treschow, he
nonetheless lived his childhood in Denmark and studied in Copenhagen and in Kiel.
The great spiritual movement in Germany drew him to itself, and after a rich
intercourse with the greatest minds of the age in the fields of poetry, natural
research, and philosophy, he returned in 1802 to Copenhagen, where in the winter of
1802--1803 he delivered his famous lectures, which are in part available in print
(\emph{Indledning til filosofiske Forelæsninger} [Introduction to Philosophical
Lectures], Copenhagen 1803).

\begin{center}---\end{center}

The word ``romantic'' (or, as was earlier said, ``romanic'') seems to have come
from England to France and Germany. It comes from ``roman,'' a fanciful tale of the
kind that first developed among the Romance peoples. At first it was used in a
disparaging sense: wild, fantastic, lawless --- later, especially from Rousseau's
time, in a laudatory sense. The word is found used of the Alps --- first in the
first, then in the second sense.\footnote{Compare on this Ludwig Friedländer's
excursus in the second part of his \emph{Sittengeschichte Roms}, 5th ed., II,
pp.~245--247.} It denotes the opposite of the clear, the definite, the surveyable,
and the intelligible, and is used first of landscapes and of human characters. It
has become the standing expression for the fundamental mood and the way of thinking
in Germany and in the North in the first generation of the nineteenth century.

Romantic is the infinite, that which leads out beyond all limitation, so that the
mind's longing can wholly unfold itself. The romantic forms the opposite of the
classical in art and of the rational in thought.

Romantic is the far-off, which awakens melancholy and longing, especially the
far-off in time, the old times that have long since vanished. The age of
Romanticism lived --- or thought it lived --- in the memory of the great and
valuable that was past, whereas the eighteenth century lived in joy over what had
come, or in hope of what was to come. Antiquity and the Middle Ages stood for
Steffens in mystic radiance as times in which the oppositions that now divide life
had not yet appeared --- ``in which science, art, and action, in unceasing union,
were able to bring forth the most divine and highest'' --- and he asked: ``Is not
the most beautiful and glorious thing that permeates our age the sense for that old,
vanished one?''

What one longed back to was especially, as is seen from the cited utterances, the
gathered, concentrated way in which life, the spiritual as the material, could be
led before the division of labour had intervened dividingly. Rousseau and Schiller
had already energetically brought forward the unfortunate consequences of the
division of labour, and not least on this point are they the forerunners of
Romanticism. Steffens states (in his seventh lecture) that what especially
distinguished that vanished time was the unity of what we must now separate as word
and deed. ``In living action everything was expressed that we now seek by words ---
all the revolutions and developments of nature, all the mysteries of religion. Life
itself, fresh as it sprang from the Godhead, had immediately the meaning that we now
only divine. Therefore all representations were living. No trace of the
disintegrating death that in more recent times has made the sciences and poetry,
and through them religion, so monstrous.''

This disintegration is now to cease, the division of labour is to be abolished, and
religion, poetry, and science are again to become one! --- Such was the gospel of
Romanticism.

\begin{center}---\end{center}

In order that philosophy may enter the service of the new gospel, the method must be
changed. Steffens demands of his hearers --- besides certain general knowledge and
an interest in solving the riddle of existence --- a capacity for the intuition of
wholes, which lets all particulars stand as members of an absolute whole. He wishes
to give ``a living system, not a dead classification'' --- wishes to ``open a more
significant view over life and existence than the one to which ordinary experience
in daily life, hemmed in by finite need, leads us. From the deeper view over nature,
history, and our innermost temper, the philosophical problem springs as from its
proper source.''

But the whole cannot be reached by the piling up of particulars, by an aggregate of
facts, such as naturalists and historians as a rule give us. If they bring about a
connection between the separated parts, it becomes only a quite external, mechanical
connection. For the particular is understood only in connection with the whole, and
therefore an intuition of the whole forms a necessary presupposition for any
occupation with particulars. And in all great, foundational researchers there
stirs, accordingly, such an intuition, even if only in the form of instinct or
presentiment. --- In Steffens' own studies he was animated --- as he says in his
autobiography --- by a dark presentiment of an all-embracing life, long only as free
imagination and of more poetic than scientific significance.

It was Romanticism's fateful error that it was so sure of the validity of its
presupposition that it did not apply the right means to work its way step by step
toward the great goal that was set by the presupposition of the unity of existence.
It revelled in the idea of this goal instead of being animated by it to an
absorption in the particulars. Since mechanical natural science, which for two
hundred years had shown how manifold changes in nature could be explained as motions
produced by pressure, impact, or attraction, did not seem to satisfy the romantic
ideal, it was regarded as an aberration. Science, Steffens complains, has closed
nature to us; the key to nature's secrets can be found only in the interior of our
own spirit. Nature must be regarded as the prehistory of the spirit and read from
above downward. --- Instead of the analogy on which natural science had hitherto
supported itself --- the analogy between change and motion --- Steffens will
therefore set another: the analogy between changes in nature and psychological
processes. Nature is to be seen as a series of stages by which the slumbering spirit
successively awakens to consciousness. Romanticism has its justification against a
materialistic dogmatism (or, as it has been called, ``a mechanical mythology''), but
the analogy it will set in place of the natural sciences' fruitful mechanical
analogy can lead only to poetry, not to science. ---

One might think that if nature is to be the prehistory of the spirit, there must be
talk of a real development, a development in time, as in the modern doctrine of
evolution. But Steffens expressly declares it a false hypothesis to assume a real,
actual transition from one life-form to another. He acknowledges that it is a
presentiment of nature's unity that stirs in the naturalists who embrace such an
assumption; but here this presentiment has led astray. The apparent becoming of new
forms is only a semblance, arising from the relativity of all oppositions. As
Steffens says in a later German work: ``Die Stufenfolge der Formen ist nur für den
Schein'' [The gradation of forms is only for appearance].\footnote{\emph{Grundzüge
der philosophischen Naturwissenschaft}. Berlin 1806, p.~187 (cf.\ 170).} The
difference of time was, for the romantic consciousness, too external to have
decisive significance. One felt oneself so deeply into nature that one was equally
at home in all forms and saw each of them as an eternal expression of the One that
reveals itself in all.

Nevertheless Steffens lays great weight on the difference between the various forms
of nature, and he finds it especially in the different degree of individuality that
the various beings possess. There expresses itself in nature an individualizing
tendency, which reaches its high point in man. In him and through him nature's
centre-seeking, inwardizing tendency reveals itself. Man has a whole world over
against him; but therefore a whole world also reveals itself in his interior. The
prehistory of the spirit is here concluded --- at the very point from which we must
set out in order to understand the other beings of nature. The last section in the
German work that underlies the Danish lectures\footnote{\emph{Beiträge zur inneren
Naturgeschichte der Erde}. Freiberg 1801, p.~275 ff. --- The work is dedicated to
Goethe and ``laid with holy awe in the Delphic temple of the higher poetry.''} has
for its heading: ``Durch die ganze Organisation sucht die Natur nichts als die
individuelleste Bildung'' [Through the whole organization nature seeks nothing but
the most individual formation]. Through reproduction, irritability, and
sensibility, organic nature approaches --- by means of increasing individualization
--- the realm of intelligences, and everything that shows itself there already lies
as slumbering disposition in unconscious nature. --- In his autobiography
(\emph{Was ich erlebte} [What I Experienced], Book Ten) he even says: ``What I call
philosophy of nature is nothing other than the conviction that in everything that is
an object of research there must be acknowledged something distinctive, something
that develops out of itself.''

\begin{center}---\end{center}

Steffens did not come to continue his activity here at home. People would not have
that visionary as a teacher for the young. He went to Germany and became professor
at Halle, later at Breslau, and finally at Berlin, where he died in 1845.

The romantic unity of poetry, religion, and philosophy later became not enough for
him. He saw that this ideal of youth could be realized only slowly and by scattered
ways, and the urge to immerse himself in great subjects with his whole soul, which
in his youth had led him to speculation about nature, later led him to embrace the
Lutheran doctrine. His soul found rest in religious faith, and religion now became
for him the only point of light in the chaos of history. (\emph{Was ich erlebte},
Book Eight.)\footnote{Compare the detailed characterization I have given of Henrik
Steffens in my \emph{Mindre Arbejder} [Minor Works], Second Series, pp.~162--180.
--- To his fine edition of Henrik Steffens' \emph{Forelæsninger} (1905), Mr.\ B.\
T.\ Dahl has attached an interesting essay: \emph{Nogle Sider af Henrik Steffens'
ydre og indre Liv} [Some Aspects of Henrik Steffens' Outer and Inner Life]. ---}

By his thought he exercised a considerable influence, especially through Sibbern,
who came into a close relationship with him at Halle and at Breslau. It is especially
Steffens' strong assertion of the individual\footnote{It was also on this point that
Steffens influenced Schleiermacher: compare \emph{Den nyere Filosofis Historie} [The
History of Modern Philosophy], 2nd ed., II, p.~193 f.} that is traceable in later
Danish philosophizing --- especially in the line Sibbern--Poul Møller--Kierkegaard.

To the history of poetry belongs his influence on Oehlenschläger, which he himself
has beautifully expressed thus: ``I gave him to himself.'' On Grundtvig too he
exercised a great influence, probably most by the second, unpublished series of
lectures, which treated poetry and entered upon historical-philosophical
considerations. The common opposition to rationalism was a point of union; and
Grundtvig stated (in a letter to Ingemann in 1824) that his own poetic-historical
tendency was a continuation of Steffens' from 1803. Yet he felt a great opposition
between them and declared at the same time that they were poles apart. What
especially caused the separation was that Steffens would have nothing at all to do
with Grundtvig's ideas about ``the North's, and especially Denmark's, spiritual
conditions and deep historical importance''; such thoughts were in Steffens' eyes
``Judaism'' --- but for Grundtvig, Denmark was precisely ``history's Palestine, the
nation-land where Christianity is to attain its highest earthly
transfiguration.''\footnote{In the letter to Ingemann referred to (see
\emph{Grundtvigs og Ingemanns Brevveksling 1821--1859}, p.~35), Grundtvig mentions
an ``ultimatum'' that he sent Steffens (after the latter's visit to him in 1824).
This is now printed in Fr.\ Nielsen: \emph{Grundtvigs religiøse Udvikling},
pp.~239--243.}

\begin{center}---\end{center}

There were two young Danish men who needed no awakening from the apostle of
Romanticism: the Ørsted brothers. In quiet they had themselves appropriated what the
new tendency could bring of value. Themselves having come from the best and deepest
in the eighteenth century's intellectual development, they stood, by their great
self-labour, more freely toward the newer tendency --- on the whole, two of the most
many-sided minds that Denmark has produced. We shall return to them later. Here
shall be cited only an utterance of H.\ C.\ Ørsted about Steffens in a letter to
Oehlenschläger (1 November 1807), who had become angry and unjust toward his former
friend:

``As for Steffens, I know very well that he wishes to squeeze nature into a system,
and that a very narrow one. God preserve me henceforth from becoming a
Steffensian\ldots\ With a rare manifoldness, though not depth, in his knowledge, he
unites manifold happy insights. When one does not take these for incontrovertible
results of the deepest philosophy, but for what they are --- flashes of lightning
for guidance --- then one learns much from him. Thus I have done, and can without
danger continue to do.''

\begin{center}---\end{center}

% ============================================================
%  5. HANS CHRISTIAN ØRSTED   (Da.: "5. HANS CHRISTIAN ØRSTED")
% ============================================================
\dfchapter{5. Hans Christian Ørsted}
\label{ch:df-hcoersted}

\noindent A ferment such as Romanticism was must prove its justification by the
development, and especially by the work, it leads to. No one was so well prepared to
receive the awakening influences of the new current, without yet being entangled in
it, as Hans Christian Ørsted.

Born in Rudkøbing (1777), where there was very little access to instruction, he had
to procure for himself, by self-study and by collaboration with his brother a year
younger, the knowledge he needed to become a student. There followed then some rich
years of youth in Copenhagen. His many-sided interest showed itself in that ---
while pursuing his pharmaceutical study --- he won the prize for an aesthetic and a
medical essay and eagerly studied Kant's philosophy, an interest he shared with his
brother. But while the latter interested himself most in Kant's ethics and
philosophy of religion, it was Kant's conception of the knowledge of nature and its
presuppositions that especially interested him. In \emph{Filosofisk Repertorium}
(1798) he published an essay on \emph{Naturfilosofiens første Grunde} [The First
Principles of the Philosophy of Nature], in which he first gives a survey of Kant's
work \emph{Metaphysische Anfangsgründe der Naturwissenschaft} and then develops his
own conception, especially through a critique of the atomic theory, which will still
have interest. Remarkable it is, however --- but characteristic of the age --- that
although the said work of Kant is already more dogmatic than Kant's own theory of
knowledge permitted, the young author nonetheless thought he had to go a step
further. Kant's method consisted in constantly asking on what presuppositions
scientific experience rests, and it was by this way that in his said work he thought
he could establish that there could be only two moving fundamental forces in the
material world, an attracting and a repelling one. A material point can, he thought,
act on another material point only in a straight line between the two points, and
here only two motions are possible: one by which the two points approach one
another, another by which they move apart. To this Ørsted fully subscribes, but he
thinks that Kant's line of thought leads only to a hypothesis, because it consists in
an inference back from experience. Ørsted will instead derive all natural laws
``from the nature of our cognitive faculty, which Kant has so excellently developed
in his \emph{Kritik der reinen Vernunft}.'' But he overlooks that Kant, in this his
principal work, had precisely impressed that it was only the fundamental
presuppositions of experience that found their explanation in the very nature of
human knowledge, while the particular natural laws had to be found out hand in hand
with experience itself: Kant's \emph{Metaphysische Anfangsgründe der
Naturwissenschaft} deviates in this respect from the principal work.

When Ørsted wrote the said essay, in which he wished to be more Kantian than Kant
himself, Schelling's first works on the philosophy of nature had already appeared.
But against them Ørsted takes a decided stand. He states: ``These books certainly
deserve attention for the beautiful and great ideas one finds in them, but they lose
much of their value on account of the not quite strict method by which the author
mixes in empirical propositions without sufficiently separating them from the a
priori propositions --- especially since the empirical propositions he adduces are
often fundamentally false.''

A few years after the appearance of this essay came Ørsted's years of travel as well
as his first experimental works. Already before the journey he had begun his teaching
activity at the university (1800), which was to become of such great significance for
Danish science and intellectual life, as well as for our material culture. We shall
not here follow his development in detail, but shall keep especially to the book in
which, in his old age, he gathered his conception of natural science and its
significance for one's view of life and of the world: \emph{Aanden i Naturen} [The
Spirit in Nature]. This book, unique in our literature, appeared in the year 1850,
but it consists of essays written at various times, and it bears witness to the
coherence in Ørsted's life and thought from the days of youth to old age's retrospect
on the work.\footnote{In the work \emph{Vort Folk i det nittende Aarhundrede} [Our
People in the Nineteenth Century], edited by V.\ Østergaard (1897), I have given a
characterization of the Ørsted brothers of a more popular kind than the one I give of
the two men in the present work.}

\begin{center}---\end{center}

While Steffens complained that science concealed the spirit of nature from us, Ørsted
wished precisely through science to penetrate into the spirit of nature. It was his
conviction that it was the same reason that expresses itself in outer nature and in
our own thinking and feeling. What is abiding in nature is the laws. In a waterfall
the abiding thing is not the individual bodily parts, but the law that binds them,
their motions and the effects of these --- the scattering of the drops, the formation
of foam, the sound caused by the plunging, the roaring, and the foaming --- together
into a single act of nature. We notice not merely a manifoldness, but a co-foldness,
which expresses a thought of nature, an idea. A thing's essence is the sum of the
co-operating natural laws by virtue of which it subsists. The natural laws are in
nature what thoughts are in our spirit.

The Kantian has here become a Romantic. Kant had found only an analogy between our
thoughts and their passing into one another through inference on the one side, and
the processes of nature on the other. Ørsted, like Schelling and Steffens, makes this
analogy into identity. But his ``dynamic'' conception of nature, according to which
matter is not a dead being subsisting for itself, but manifestations of activity
determined and limited by the laws of nature, he had already won through the study of
Kant.

The religious expression for ``the spirit in nature'' is God. For the fundamentally
active and the ordering in nature are not two different things, but one infinite,
living reason. God's will is religion's expression for the eternal laws of existence.
And, as Ørsted elsewhere expresses himself: ``the eternal, living reason we call God,
when we regard it in its self-consciousness, its personality.'' The opposition between
God and the world really holds only when by the world one understands that in
finitude which removes us from God. In the highest moments the opposition falls away,
in that what one calls the world stands as an effect of the Godhead.

\begin{center}---\end{center}

When Ørsted conceived matter as derived from force --- when in his work \emph{Den
almindelige Naturlæres Aand og Væsen} [The Spirit and Essence of General Natural
Science] (1811, reprinted 1847) he taught that matter is nothing other than space
filled by nature's fundamental forces, and that what gives things their unalterable
peculiarities is the laws according to which they are produced --- he has thereby
stated a view that in the more recent history of physics has received increasing
confirmation. Physics has more and more developed into a doctrine of natural forces.
With the cited statement of Ørsted from 1811 may be set alongside a statement of
Faraday from 1846: ``Neither in what is called empty space nor in bodies can I
discover anything but the forces and the lines according to which they act.''
Professor Christiansen sees in Ørsted's thoughts on this point ``a foreboding of what
a distant future was to bring.''\footnote{Chr.\ Christiansen: \emph{H.\ C.\ Ørsted som
Naturfilosof} [H. C. Ørsted as Philosopher of Nature] (Overviews of the Danish
Society of Sciences, 1903), p.~488.}

\begin{center}---\end{center}

Despite the sharp judgment that Ørsted in his day passed on Schelling's and Steffens'
philosophy of nature, he too nevertheless sometimes came upon fantastic analogies and
combinations. Thus when he sets the cardinal points of the compass alongside oxygen,
hydrogen, carbon, and nitrogen, and finds an agreement between the electrical figures
and the organic forms. But there was one thought that the romantic philosophy of
nature had uttered, and that guided Ørsted in a fruitful way in his more specialized
investigations --- it was the thought of the kinship and connection of the forces of
nature. This thought he pursued in detail. It led him to conceive the ray of light as
a series of small electric sparks. He has thereby anticipated the electromagnetic
theory of light later developed by Maxwell and Hertz.\footnote{Compare Chr.\
Christiansen: \emph{Den elektromagnetiske Lysteori} [The Electromagnetic Theory of
Light] (Transactions of the Danish Society of Sciences, 1889), pp.~189--193.} --- Of
special significance for his research, however, the idea of the inner connection of
the forces of nature became through this: that step by step it led him to undertake
the experiments by which he (July 1820) discovered that an electric current affects
the motion of the magnetic needle. With pride he could declare that this discovery was
not due to chance, but was the result of long reflection. He had on this point had to
correct the presuppositions he had in his youth taken over from Kant. According to the
way in which Kant had systematized Newton's teaching, there could be found only two
moving forces, attraction and repulsion, acting along the line connecting two material
points. But if an electric current was to have influence on a magnet, the motion had
precisely to go across the connecting line. That thought of the philosophy of nature
he could then carry through only in defiance of the authority of Newton and
Kant.\footnote{Compare Chr.\ Christiansen: \emph{Nogle Bemærkninger om
Naturvidenskabernes Betydning for vor Tid} [Some Remarks on the Significance of the
Natural Sciences for Our Time] (University Programme, 1905), p.~41. --- As is well
known, Ørsted did not pursue his discovery further. It was Ampère who first saw in the
fact demonstrated by Ørsted a single example of a general law, which he carried
through theoretically and experimentally. Philosophy is not to blame for Ørsted's not
going further; for Ampère too was a philosopher and has a fine place in the history of
philosophy. See \emph{Den nyere Filosofis Historie} [The History of Modern
Philosophy], 2nd ed., II, pp.~304--308.}

\begin{center}---\end{center}

As is seen from the foregoing, there was for Ørsted an intimate connection between his
scientific, philosophical, and religious circle of ideas. It was also one of his
life's chief thoughts that natural science should not stand isolated in human life,
but should reach into thought, faith, and action. With the publication of \emph{Aanden
i Naturen} he aimed, as he openly stated, at a change in the human view of the world.
He conceived the very cultivation of science as religion: it seeks, after all, that
which abides under all changes --- ``the bond of unity which brings it about that
things, in all their manifold distributions and separations, are yet not scattered
apart.''

In his biography of H.\ C.\ Ørsted, Carsten Hauch says of him: ``It was his conviction
that in a thorough conception of the real conditions of the world there lay a deeper
philosophy than that found in most philosophical systems, and a truer poetry than that
found in many poems; indeed, he even thought that our religious views too could
thereby be set right.''

In the great, splendid letter to Oehlenschläger of 1 and 6 November 1807, in which he
reads his friend the riot act because he wished to run away from school too early, he
expresses himself especially on the relation between science and poetry. The thinker,
he thinks, will only at the end of his wandering reach the point where the world of
art begins, while conversely the poet, when he has formed his presentiments with full
clarity, approaches science. He did not have Romanticism's shyness of science and
reality. And in his poem \emph{Luftskibet} [The Airship] he himself made an attempt to
show how natural science can enter the service of poetry --- an attempt that, however,
may precisely serve to show that it is human striving and human destiny that is the
soul in all poetry: the poem is at bottom about the fate of Montgolfier, the inventor
of the airship, not about natural phenomena. There is not much poetry in it, though.
The poetry was in Ørsted's own soul, in his thinking, his enthusiasm, and his work,
but he was not able to lay it into his poem.

He reproached the poets for making superstition poetic, instead of finding the poetic
in nature itself. And by superstition he understood the imagination that something can
act in nature otherwise than according to its laws. His aim with the publication of
\emph{Aanden i Naturen} was, as he writes in a letter (8 February 1850), ``to show the
often over-poetic and over-credulous people of our age that the understanding, which
they so frequently mistake, has the great calling to guard the sanctuary of conscience
and faith --- from which, to be sure, I would have a good deal of the old leaven of
superstition and prejudice swept out.'' Already earlier --- in his address at the
millennial festival of Christianity in Denmark (1826) --- he had also defended the
eighteenth century's ``boldly inquiring men,'' over whom ``certain of the proud and
hard men of our time break the staff in the name of Christianity.'' He did not share
Romanticism's judgment on the immediately preceding age. There was, after all, in his
own development too an inner connection between what he had brought with him from the
preceding century and what he had learned in the new one.

\begin{center}---\end{center}

The line of thought in the first part of \emph{Aanden i Naturen} led to a polemic with
Bishop Mynster, which especially concerned two points. --- Mynster maintained that
development must have begun with an original state of perfection. But Ørsted teaches
that everything in the world begins from something undeveloped and runs through a
series of stages. Why it is so, he cannot say, but experience shows it. The
theologians' doctrine that nature has become worse through man's sin cannot be
defended. The laws of nature have not changed, and there was suffering and death in the
world before man came into being. ``I submit it to the dogmatists to consider how far
their science's doctrine of sinfulness must in all respects be regarded as
incontrovertibly correct, or might gain by a new reworking.'' --- The other point was
how far our wishes and our need can determine what is truth. Against this Ørsted
expressed himself decidedly and in beautiful, heartfelt words: ``Should we not be
ashamed in our innermost being, if we caught ourselves wishing to have a truth other
than the real one?\ldots\ No, let us give the truth the honour! With it the good is
indissolubly bound up. The full truth itself brings its consolation with it.''

In this statement there makes itself known not only Ørsted's great love of truth, but
also his bright view of life. He was a great optimist. A guarantee that the truth
always brings consolation no one can provide, and besides, there may be a long way to
the full truth!\footnote{Compare my remarks on Ørsted's statement in my
\emph{Religionsfilosofi} [Philosophy of Religion] (§~115).} But the utterance is
characteristic of H.\ C.\ Ørsted's happy nature. Fredrika Bremer, the Swedish
authoress, took him to task for his optimism. He promises not to forget the dark side
of life. ``You forced me,'' he writes (1849) to her, ``to subject my thoughts on the
misfortunes, the pain, the hideous, to a new consideration, which has not been
unfruitful for me, and which I still continue. In my joy over the harmony of the
whole, I shall strive to grant pain and sorrow their lawful rights!''

Thus he wrote two years before his death. His considerations hardly altered his view
of life. His own life had, in the best sense, been happy. His youth's enthusiasm bore
fruit in a great discovery and in the presentation of a significant view of the world,
as well as in practical work for his science's influence on life (by the founding of
the Society for the Dissemination of Natural Science and by the founding of the
Polytechnic College). There was harmony in his personality as in his thought. The
enthusiasm of youth bore fruit in the work of manhood, and old age did not lose the
understanding of the thoughts of youth.

\begin{center}---\end{center}

% ============================================================
%  6. GRUNDTVIG AND H. C. ØRSTED   (Da.: "6. GRUNDTVIG OG H. C. ØRSTED.")
% ============================================================
\dfchapter{6. Grundtvig and H. C. Ørsted}
\label{ch:df-grundtvig-oersted}

\noindent In H. C. Ørsted, one of the elements that Romanticism wished to fuse into
unity had been separated out and asserted in its independence --- science. And the
relation between science and the other elements in Romanticism's unity was for him
rather this: that it should lie at the foundation, or that through it we should come
to poetry and religion.

In the opposite direction moved another of our people's great minds --- Grundtvig.
For him religion became the main thing, and poetry and science were to be
subordinate to it, without independence in relation to it. No wonder that a sharp
collision came between the two men. A significant dispute was carried on between them
in the years 1814 and 1815.

One might wonder that Grundtvig had anything at all to do with philosophy. And yet it
too intervened in his course of development and even furnished motifs that continued
to work after he had reached his definitive standpoint. Our most national tendency
received important impulses from German thought.

Steffens' lectures --- or, as Grundtvig says in \emph{Kirkespejlet} [The
Church-Mirror], ``the frightful noise that my half-German cousin, Henrik Steffens,
made in Ehlers' College, so that the Regensen nearly fell down'' --- awakened him
especially in a poetic-historical respect, as already mentioned above. The sense for
the great and original, Romanticism's best side, was strengthened in him, nourished
as it had already become by the great events of the time and by absorption in the old
sagas. Twenty-two years old, he went (1805) as a private tutor to Langeland, and it
was here that he brooded with Schiller, Fichte, and Schelling over the riddles of
life. Fichte especially had a great influence on him. He was gripped by the way in
which Fichte impressed the necessity of a conviction fixed in the will and proceeding
from man's interior, in opposition to barren reflection and fluttering notions
(\emph{Die Bestimmung des Menschen}, 1800). When Fichte, in \emph{Grundzüge des
gegenwärtigen Zeitalters} (1806), spoke out against the age's spiritlessness, its
dead learning and self-satisfied enlightenment, and not least when he emphasized the
significance of the spoken word in opposition to the written, these were thoughts
that bore fruit throughout Grundtvig's life and activity. And undoubtedly he was also
captivated by the way in which Fichte later (in \emph{Reden an die deutsche Nation},
1808) enthusiastically asserted the imperishable significance of nationality and
referred in particular to the German and Scandinavian peoples as those who from
primeval times had dwelt on the same soil --- as well as by Fichte's educational
ideas, which demanded a development for the young by the way of free appropriation and
self-activity. When Fichte died, Grundtvig --- who by then, however, had reached his
pronouncedly positive-religious standpoint --- honoured him in a beautiful memorial
poem.\footnote{Compare my essay: \emph{Fichtes Taler til den tyske Nation og de danske
Nordslesvigeres nationale Kamp} [Fichte's Addresses to the German Nation and the
National Struggle of the Danish North-Schleswigers] (\emph{Mindre Arbejder} [Minor
Works], Second Series, pp.~192--200). (On Fichte's influence on Danish intellectual
life see especially p.~195 f.)}

Poetry and philosophy were not enough for the young man, whose mind was in such strong
movement. He takes leave of the philosopher in a remarkable religious-philosophical
essay from the year 1807 (printed in \emph{Theologisk Maanedsskrift} [Theological
Monthly]): \emph{Om Religion og Liturgie} [On Religion and Liturgy]. It is a
religious-philosophical essay --- for one cannot, after all, take leave of philosophy
except through philosophy.

The first question is: Do we need religion? --- We live by an open sea. Poetry seeks
to fly across it, but vanishes in the clouds, which are lit by the radiance from the
distance. Philosophy investigates whence this radiance comes --- builds ships for
voyaging, but is stopped by the waves and must either become poetry or content itself
with understanding the Earth and life there in the finite world. Philosophy is the
age's clever daughter, but does not know her mother's birth. It can point out the
disharmonies of life, but not show how they arose from an original harmony, or how
they can be resolved into harmony. For all its speculation it does not, after all, get
beyond a hypothesis. So we must then either give up gaining knowledge of the relation
between time and eternity --- but that would lead back to animality --- or the
knowledge must come from above.

If it is next asked: What is religion? --- then history bears witness that it is a
divine force which raises man above all the surroundings of the Earth to the highest,
the supersensible, and by which a community is established between the finite and the
eternal. In Christianity these features appear clearly and with perfection. In it,
poetry and philosophy, which in the heathen religions always lay in strife, are united
into unity, and they can therefore both, each in its way, express religion in the
world of finitude.\footnote{The editor of the Monthly (P. E. Müller) declares in a
note that he ``does not quite understand this.'' I understand it thus: religion is the
fundamental thing, but poetry bears witness to it in its images, philosophy in its
questions. This interpretation accords with the essay's line of thought and with
Grundtvig's standpoint at that time.} Poetry gets its reality, and thought can work
with its approximations. But the unity itself of poetry and philosophy is a mystery
that one denies, both when one holds to poetry and when one merely holds to
philosophy. And it is in the liturgy --- especially in the sacraments --- that this
mystery finds its expression. --- And with that the young author has then made his
way over into theology and to his proposals for the ordering of the church service.

Although this essay lies prior to Grundtvig's religious breakthrough, it is of great
interest as testimony to a striving to pose the religious problem in sharper form than
rationalism and Romanticism made possible. Religion's difference from poetry and
science is sharply emphasized --- and yet the possibility of a harmony between all
three domains is still maintained, only in such a way that it is the religious element
that determines the timbre. It remains unclear how there can be said to be won
``knowledge'' of the relation between the eternal and the finite, when the result
becomes a mystery. Further: how can history show us divine forces at work? History
can, through its investigations, present what has happened and what forces have been at
work; but that these forces were divine, it cannot establish. And finally: can it
really stand so with us that we have only the recourse of sinking into animality if we
cannot win full knowledge of the relation between time and eternity? ---

These critical misgivings did not make themselves felt for the young theologian, who
now went to meet his inner spiritual struggles. He was outraged at a rationalist
historian's treatment of the Crusades and was led into a period of enthusiasm and
combativeness, until he was shaken through by the personal question of his own soul's
salvation and brought to the brink of madness. In a letter to Molbech (9 March 1811)
he declared shortly afterward: My whole condition would be the most inexplicable
riddle, if Satan were not at work in it! ---

After having gone through this crisis, he then appeared (especially in
\emph{Verdenskrøniken} [The Chronicle of the World], 1812) with the Bible in his hand
against the slack age, its literature, and its science. When Molbech had attacked his
book and its treatment of various personalities as unhistorical, and demanded that
history should be independent of theology, Grundtvig answered: ``If a Christian is to
write history, then he must write it biblically, call all learning that conflicts with
the Bible outright, without further proof, lies and error, let the Bible judge men and
ages, for it alone has the right to do so; it is Christ's viceroy on
Earth.''\footnote{\emph{Chr. Molbech og N. F. S. Grundtvig. En Brevveksling} [Chr.
Molbech and N. F. S. Grundtvig: A Correspondence], ed.\ L.\ Schrøder, 1888, p.~123
(cf.\ p.~101).}

There is now no more talk of religion as a higher, though mystical, unity of poetry
and philosophy. In the Bible religion speaks with authority, and all other spiritual
striving must subordinate itself to it.

While Molbech defended history against the prophet's hard judgments, H. C. Ørsted
entered the lists on behalf of philosophy and natural science.\footnote{\emph{Imod den
store Anklager} [Against the Great Accuser]. By H. C. Ørsted, Professor. 1814. (Motto:
Thou shalt not bear false witness against thy neighbour.) --- Grundtvig's reply bore
the title: \emph{Imod den lille Anklager, det er Prof. H. C. Ørsted, med Bevis for, at
Schellings Philosophie er ukristelig, ugudelig og løgnagtig} [Against the Little
Accuser, that is, Prof. H. C. Ørsted, with Proof that Schelling's Philosophy is
Un-Christian, Godless, and Mendacious]. By N. F. S. Grundtvig, Priest, 1815. (Motto:
Stopped shall be the mouths of those who speak lies.)}

In the collision between Ørsted and Grundtvig, two opposite worlds of thought met ---
that of science, with its constant striving for development, and that of biblical
faith, with its striving to bend everything spiritual under a single, once-and-for-all
given circle of ideas. And it was two opposite natures that met. The one --- one of
the healthiest souls that Danish intellectual life knows --- went his cheerful way
through the world of research, sure that nothing true and great, no real value, would
be lost by the constant searching, but that life would precisely become richer by it.
The other had, in a grievous crisis, gotten his gaze fixed on a single point in the
world of the spirit and of history, and that in which he had found his hold he now
held up before the age as the one thing and the highest --- as that which was to judge
all and everyone, but itself be judged by no one.

The dispute between them concerned the independence of science. Grundtvig declares
that it is not science ``itself'' that he combats --- only its independence. All
science is to be governed by a Christian idea. Ørsted answers to this that a
non-independent science --- a science that takes any consideration whatever other than
the one its subject demands --- is not science, and that we do, after all, come to
know nature better through nature itself than through the Bible! --- There was also
dispute about whether Schelling's philosophy is Christian or not; this has nowadays
less interest and would lead us too far. On the other hand it has its interest to see
somewhat more closely how Grundtvig at this time conceived the relation between
religion and science.

Philosophy as independent thinking can --- Grundtvig thinks --- have its significance
in an un-Christian age. Thus Kant and Fichte had a great work to perform by their
strict ideality and their demand that the truth about man, his condition, and his
vocation be found by the way of thought. When the age did not follow these
philosophers, it was because they demanded too much of the slack time, which had
abandoned Christianity only in order to avoid the struggle that goes with all higher
striving. From the standpoint of Christianity an independent philosophy is a
monstrosity; here reason is to submit to the written word, acknowledging its divine
power to overcome the world, and then a Christian logic is a quite glorious thing when
it is a matter of stopping the mouths of the deniers and mockers. --- In what this
``Christian'' logic differs from the usual logic --- whether it has other principles
than this --- one unfortunately learns as little here as later in Martensen, when he
proclaimed ``regenerate'' cognition as theology's organ.

It goes no better with the other sciences than with philosophy. Grundtvig asserts that
where chemistry, astronomy, and mathematics truly flourish, spiritual life is near its
ruin. It is attractive for unbelieving souls to divert themselves with them, since
they confirm the notion that man is only a body that dissolves at death, and since
occupation with the natural causes of the great phenomena of the heavens easily leads
to forgetting God. In historical research, whose independence Grundtvig had denied
against Molbech, he finds no such danger. Yet it holds of all sciences: when they wish
to be independent, Christianity is in their way, and they are in Christianity's way.

It is not, however, only from a religious standpoint that Grundtvig is doubtful of
natural science. ``I regard,'' he says, ``the passionate study and higher valuation of
physics, especially of experimental physics, which ends in mechanics, as a sign of the
near grave of scientific spirit.'' --- Here there makes itself known the romantic
opposition to the more recent natural science, as it had developed from Galileo's
time. A new concept of science was --- as we have seen in Steffens --- to take its
place, but was as little formulated as the principles of the ``Christian logic.''

The main thing, however, is for Grundtvig the religious standpoint. It is, after all,
religion that is to teach us to know man's condition and vocation and the way to
attain it, and that is to govern our desire, thought, and deed. It must then also be
consulted before we seek any other knowledge, and must stand livingly before our eyes
while we occupy ourselves with the sciences. For the believer, science is either
something abominable or else religion's servant. And by religion Grundtvig thought of
``the old-fashioned, simple faith,'' whose content is found in the Bible, understood
quite literally. ---

It was not only Grundtvig's personal need that thus led him to demand science
subordinated to religion. His great interest in the popular cause also contributed
here. The people's spiritual unity would, after all, be secured when the unlearned
found their spiritual nourishment in the Bible, taken according to its letter, and the
learned bent their science under quite the same authority. It was the continuation of
this line of thought that later (1825) led him to the discovery that not the Bible, but
the three articles of faith, were the real foundation. The Bible can, after all, be
interpreted in various ways, and it requires learning to be rightly understood. This
Grundtvig did not think was necessary as regards the articles of faith, and only
through their lying at the foundation can the unlearned become independent of the
learned's testimony.\footnote{\emph{Kirkespejlet}, p.~343.} ---

Grundtvig is right that the learned should not be authorities. The learned's right to
be heard rests on their reasons and their proofs, and whether these hold can likewise
be decided only by reasons and proofs. The world of thought stands open to everyone,
and the learned, who are just as much a part of the people as the unlearned, have the
duty to guard the world of thought as a part of the people's spiritual goods; but this
duty they can perform only when their work can be done in full freedom and
independence. It depends, of course, on the unlearned themselves how far they will make
use of the guidance the learned can provide.

The proposition that the unlearned should not be dependent on the learned's testimony
really means a retreat from the standpoint in the dispute with Ørsted, where he
conversely wished to make the learned dependent on the unlearned's testimony. And yet
Grundtvig also later had difficulty understanding that one can cultivate science for
its own sake, and that it can precisely thereby become of great human value. In the
winter of 1843--44 he delivered lectures on Greek and Norse myths and stated then that
he did not comprehend how one can take an interest in ``infinite space and the
Copernican system,'' when one did not start from the assumption that the zodiac and the
Milky Way, sun, moon, and all the stars exist for man's sake. He objected moreover to
astronomy that it had become too mathematical; it had become black on white instead of
gold on blue!\footnote{\emph{Brage-Snak om græske og nordiske Myter og Oldsagn}
[Brage-Talk on Greek and Norse Myths and Antiquities], 2nd ed., p.~1; 165 f. ---
Grundtvig is said constantly to have held fast to the belief that the sun moved around
the Earth. Compare Sofus Høgsbro: \emph{Mit Forhold til Grundtvig, Tscherning og
Monrad}, 1902, p.~13. --- Grundtvig still stands toward Copernicanism as Luther and
Melanchthon stood, when their young, enthusiastic astronomical colleague Rheticus,
Copernicus' first disciple, brought the new doctrine into the Wittenberg high school.
See \emph{Den nyere Filosofis Historie} [The History of Modern Philosophy], 2nd ed., I,
p.~103. --- For the rest, it is not Copernicus himself who brought about so decisive a
change in the view of the world; such a one came first through Giordano Bruno's
consistent extension of Copernicanism. \emph{Op.\ cit.}, pp.~117--124.} ---

From Romanticism's speculative ideas he moved away in his later years. In his both
delightful and brilliant way, on the other hand, he spoke out in the said lectures ---
in the very midst of Hegelianism's dominion here at home --- for an empiricist
philosophy after the English model. Odin, he says, means the spirit of the North, and
Odin was what the German philosophers with a very lofty air call an empiricist. He
wandered far and wide, and he looked far around from Hlidskjalf; he had besides Hugin
and Munin, and he consulted Mimer. It is much wiser, with old Odin and the Englishman,
to take life and the world as they are, and to make them as useful to us as possible,
than, with the young philosophy and the Germans, to waste one's time and one's powers
on recreating life and the whole world after one's own head, and then to fly into a
rage because it will not succeed, because the impossible will not let itself be
done.\footnote{\emph{Brage-Snak}, pp.~201--205.}

It is a sound realism that is here extolled; whether it is the ``Christian logic'' that
Grundtvig in his youth prophesied is another question. There lies in it no solution of
the question he discussed with H. C. Ørsted. And with all the significance that the
spiritual tendency proceeding from him has had for our people, that question
nonetheless stands ever open and will perhaps even take on sharper forms than ever
before: how the spiritual unity can be preserved in the people despite the opposition
between the ``unlearned's'' faith and the ``learned's'' independent science.

\begin{center}---\end{center}

% ============================================================
%  7. BAGGESEN AND GRUNDTVIG   (Da.: "7. BAGGESEN OG GRUNDTVIG")
% ============================================================
\dfchapter{7. Baggesen and Grundtvig}
\label{ch:df-baggesen-grundtvig}

\noindent A few years after the collision between H. C. Ørsted and Grundtvig came the
peculiar collision between Baggesen and Grundtvig. The two poets collided over the
great question of poetry's relation to life --- or rather over the question of the
measure of value for the life that poetry is to express. It was building up to a
battle in poetic form between great thoughts --- but unfortunately it blew over. Our
literature got only the heralds' signals. Perhaps poetry cannot give more, either.

We left Baggesen as a rapturous Kantian. The rapture was stronger than the criticism
--- and he soon found what he --- provisionally, as always --- was seeking, in
Jacobi's standpoint of feeling. In a long letter to Reinhold (of 20 November 1801) he
gave an account of his new view. ``Believe me, my dear Reinhold, faith is higher than
knowledge --- is the non-rational in reason, reason's proper power, without which its
faculty would be a blind faculty for empty forms.'' It marks a conclusion of all
reflection and analysis and cannot itself be analysed. It is a power that cannot be
divided, counted, measured, or comprehended --- but it is what bears life. In it the
Holy reveals itself. Religion is to sound not merely in our knowledge, but in our
whole soul. One must not take away the shudder it awakens --- not rob the starry
heaven of the night in which alone it is visible.

He knows very well that, and why, he is no good as a scientific philosopher: ``In the
first place I have more imagination than understanding, more philosophical fantasy
than philosophizing reason\ldots\ In the second place my language is too rich in
images. In the third place I am perhaps already by nature too thoroughly penetrated by
the true, for me to exert myself to seek the truth by an artificial way. In the fourth
place I do not believe in any completed system of human knowledge, because I believe
that the human spirit can never stand still, nor shall stand still.''

These last words he himself made truth. He needed to look about him more in the world
than Jacobi's faith invited to. His reflection worked at connecting nature, morality,
and religion in one great whole. Jacobi had admitted to Lessing that his standpoint of
faith could be reached only by a leap, a \emph{salto mortale}. In opposition to this
``blind faith's \emph{salto mortale}'' Baggesen now wished to set ``the seeing faith's
\emph{salto vitale}'' --- and his endeavours, which fill a very large part of his
later years, we know from the \emph{Philosophischer Nachlass} published after his
death by his sons. He wishes to give a speculative theology that develops the concept
of God as ``the truth of all truths, revealed at once in our conscience, in the
radiance of the starry heaven, and in the New Testament.''\footnote{\emph{Philosophischer
Nachlass} (1858 ff.), I, pp.~409, 465.} He now regarded himself as a philosopher and
would gladly have been a professor of philosophy in Copenhagen when a post there fell
vacant in the year 1813. He even writes to one of his sons in the year 1811: ``Not
poetry, but philosophy has from my childhood been my subject and the object of my
endeavours,'' and still more decidedly he expresses himself in this direction in a
letter to one of his friends ten years later.\footnote{These two letters are not found
in Baggesen's biography, but are cited in Kr.\ Arentzen's work \emph{Baggesen og
Oehlenschläger}, VI, p.~34. --- This work has been of very great use to me, especially
in the study of the relation between Baggesen and Grundtvig. It is the fruit of many
years' patient work, borne by heartfelt enthusiasm for the great men of our
literature. Clumsy and pieced-together in form though it is, it is nevertheless a
treasure-chest for those who occupy themselves with the literature of this period.
There is a remarkable contrast between its form and its lively and enthusiastic
author, who had a rare ability to awaken and inspire the young, whose teacher in
Danish he was, and who often became his friends. The use I have made of his book in the
present work has renewed my grateful memory of my teacher and friend.} --- Baggesen's
psychological-theological reflections always retained an individual character, which he
also acknowledged: ``every snail must have its own house!'' There are found in his
\emph{Philosophischer Nachlass} many spirited reflections, many testimonies of serious
reflection and attempts at religious speculation. He moves, as Sibbern has said of him,
in ``the domain of philosophizing reflection expressing itself with acuteness, wit, and
humour.''\footnote{In \emph{Zeitschrift für Philosophie}, 1859, there is an interesting
review of Baggesen's \emph{Philos.\ Nachlass} by Ulrici. --- A theological assessment
of Baggesen's attempt at a philosophy of Christianity is given by Martensen: \emph{Af
mit Levned} [From My Life], III, pp.~162--172. Martensen acknowledges him as a
``Christian thinker,'' but regrets that he has not rightly penetrated into the doctrine
of the Trinity, and that he denies that we can know God as he is in himself. So far,
namely, had ``Christian thinking'' reached in Martensen's time.}

In Baggesen's poetry this background of serious reflection now and then received
profound expression --- thus in \emph{Gengangeren} [The Revenant] (1807), where it is
demanded that the poet shall acknowledge his divine calling

\begin{verse}
in eternal yearning for the Infinite\\
and in infinite and eternal desire\\
for the attuning of the Finite\\
with eternity's ground-tone in his breast.
\end{verse}

But often he played with poetry, as he played with thought --- and it was this that
called forth the collision between him and Grundtvig. In the dispute against
Oehlenschläger, Grundtvig sympathized with Baggesen and was attacked by the latter's
opponents. He thought, like Baggesen, that the great poet had gone backward and no
longer conceived his calling with seriousness. And both took up the word for art's
connection with life and its riddles. But Grundtvig found that his ally carried on his
criticism and polemic too much as a game. In his \emph{Rimelige Strøtanker ved Jens
Baggesens Grav} [Rhymed Stray Thoughts at Jens Baggesen's Grave] (1815) Grundtvig says
of him:

\begin{verse}
That was for him, then, the whole misfortune:\\
too much he had upon his neck\\
and far too little at his back!
\end{verse}

Had he --- Grundtvig thought --- had a church wall (like the one that he, Grundtvig,
had himself found) at his back, he would have fared better against his opponents.

And in Baggesen's well-known riddle (``The eternal emblem, not only the origin of
literature, but the end of philosophy. A riddle as a prize-question for the spiritually
living'') Grundtvig could see nothing but a play of thought and word with the highest.
The emblem that Baggesen meant seems (according to a letter to one of his sons) to have
been the Cross, as that which is formed of two lines and therefore expresses the great
doubleness or opposition that reveals itself in everything and prevents every line of
thought running along a single line from exhausting existence. The Cross on Golgotha is
a single, great example: God himself is revealed in suffering. These ingenuities
Grundtvig did not care for. For him the riddle was once and for all solved in the
Bible, and Baggesen's tomfoolery was to him an offence. He called out to him:

\begin{verse}
Thou art to thyself a riddle,\\
that have I long known,\\
and rightly to read it\\
thou hast but short respite! ---

What is your riddle? Only a pastime?\\
You have, though, now not much time to waste!\\
A shadow of your whole life as a bard?\\
Perhaps --- but truly, then it were ill!\\
An image of life as a whole?\\
Does it, then, consist in its broken parts?
\end{verse}

In Baggesen's reply stand the proud lines:

\begin{verse}
Well do I know the Book\\
to which you point,\\
perhaps as well as any,\\
and the solution in it;\\
but if from it is to be drawn\\
what slakes the thirst of thought,\\
the old seal must be loosed,\\
and the well first opened.\\
In vain from it you await\\
the spring-water of all life,\\
if your spirit fetch it not\\
deep with reason's pail.
\end{verse}

Baggesen's ``seeing faith'' could not content itself with the reference to ``the
Book.'' He wished, after all, to work the word of the Bible together with the knowledge
of nature and with the utterance of conscience. His thought could not proceed along a
single line. He did not indeed think that the whole consisted only in broken parts, but
that it is the broken light that we know, and that reflection must gather the parts,
without being able to reach the full unity. That was the fundamental thought in his
philosophical reflections, and it was the thought of his riddle, here, however, hidden
and disturbed with all sorts of word-play.

The battle now came to be over the question: does that view of life stand highest which
supports itself on thought, while yet taking from the religious traditions what it can
unite with the knowledge of nature and with conscience --- or the one that holds solely
to the handed-down faith in full submission?

Baggesen thought he stood highest, because he knew and beheld what his opponent
believed, and because he could easily put himself in the latter's standpoint, while the
reverse was not the case:

\begin{verse}
From the place where I stand, I can well set myself down into yours;\\
but the steps up from yours to mine are not so easy!
\end{verse}

Grundtvig answers with a reference to Judgment Day:

\begin{verse}
Whether you really stand on the temple's pinnacle,\\
and whether I have many steps up to you,\\
that will clear itself up one day at the last.\\
But if I see rightly by the lamplight of the Word,\\
I make not a single step toward your height,\\
and never shall your ripeness tempt me.
\end{verse}

Quite apart from Baggesen's and Grundtvig's special standpoints, it is a significant
question that came up between them: what is high and what is low in the world of the
spirit? What measure have we here? --- It is, of course, no dispute of rank between
persons that is at stake, but the finding of the scale of values. Grundtvig refers to
the Bible and to the Last Judgment, and Baggesen builds on his religious speculation.
But should there not be a purely human measure for the values of life?

It looked for a moment as if the battle between the two great poets would enter upon
this question in earnest. Baggesen challenges to continued battle, for ``in the world
of spirit there ought to be war!'' --- and he wants no peace until all is cleared up.
And Grundtvig bids him welcome to the battle (``Welcome, you who know the law of
spirits!'') and indicates more closely what is being fought about and how it is to be
fought. It is not the poet's wreath that they fight about ---

\begin{verse}
No --- what is light and spirit, and what is darkness,\\
yes, what is truth, what is trumpery and lie;\\
what one shall believe, whom one shall faithfully worship ---\\
is knowledge holy, or is it profane? ---
\end{verse}

And he will lay his Bible aside. When it is a matter of what and why one is to believe,
one cannot, after all, begin by appealing to one's confession of faith, which the other
party does not share. Here he utters the great and beautiful words that are fit for a
motto in all strife over life's greatest questions.

\begin{verse}
We ask not greatly here about confessions,\\
but seek the ground of the broken tones\\
and scout for light with spiritual reason!
\end{verse}

This boded well. But it stopped at the boding. The battle-signal was given, but not
followed. Baggesen began again his spiritual and bodily wandering, and Grundtvig
immersed himself in historical research and in the shaping of his ecclesiastical views,
where soon the guarantees for what one believes came to play a greater role than the
causes that lead to faith; the confession stepped into the foreground, spiritual reason
into the background. Only later was there taken up in Danish intellectual life the
discussion to which the poetic heralds had called each other, without themselves making
earnest of fighting the battle to the end.

\begin{center}---\end{center}

% ============================================================
%  8. ANDERS SANDØE ØRSTED   (Da.: "8. ANDERS SANDØE ØRSTED")
% ============================================================
\dfchapter{8. Anders Sandøe Ørsted}
\label{ch:df-asoersted}

\noindent Of the two Dioscuri who here at home met the new century with the greatest
maturity, we have already treated the one, the natural scientist. The jurist, Anders
Sandøe (1778--1860), was a year younger. His life of thought presents the same unity in
its development as his brother's. Hans Christian was the happier, on the whole one of
the happiest figures in Danish intellectual life. Anders Sandøe was heavier; his
thought moved forward through thorough testing of all particulars and possibilities.
--- It is here not as the creator of Danish jurisprudence that he is to be treated, but
as the philosophical thinker.

He lived in the circle of the best here at home. Oehlenschläger owed it to the two
brothers that he was so well prepared to receive Steffens' influence. Oehlenschläger's
sister became Ørsted's wife, and in the Ørsted house the works that the young poet sent
home from his journey were read aloud before they appeared. With his brother Ørsted
lived in an intimate intercourse built on great common interests, and men such as
Mynster and Sibbern belonged to the circle. He was passed over by the university ---
one of its sins, which nevertheless turned to the good, since he was thereby, to the
benefit of his science and his development, led out into practical activity as judge and
legislator. Later he held high offices, as attorney-general, deputy of the Chancery, and
as minister.

As mentioned earlier, in his young days he was an eager Kantian. Soon Fichte gained
great influence on him, especially by his attempt to found a doctrine of right
independent of ethics. Ørsted was a personal friend of Fichte, who lodged in his house
during his visit to Copenhagen in the summer of 1807. Although he did not continue his
philosophical studies, taken up as he became with practical work, he nevertheless
preserved the fundamental thoughts he had appropriated throughout his whole life.
Speculative gifts he did not think he had, and he looked with mistrust on the intrusion
of German speculation (Schelling's and Hegel's) here in the country. On the other hand
he thought he possessed abilities in the direction of logical and ethical thinking. Of
Kant's influence on him and his contemporaries he has, in his old age, given a beautiful
testimony:\footnote{\emph{Af mit Livs og min Tids Historie} [From the History of My Life
and My Times]. 1851, p.~34.}

``Although I have long since come back from the idolatry with Kant's teaching with which
I entered upon my literary career, its pure and powerful morality has nevertheless had
an effect on my mind that has never been lost. And I am convinced that the same is the
case with a great part of the generation that at the same time received the formation of
its manhood. But even on the later generation, for whom the Kantian teaching may seem
gradually to have become foreign, it has certainly indirectly had an incalculable
effect; for although his morality, in the form in which he handed it down, has, and that
far more than is fitting, lost its repute, the fundamental tone that ran through it has
nevertheless long maintained itself under the altered form that morality has since
received\ldots\ It has been acknowledged by many that Kant and Fichte --- the latter of
whom, in enthusiasm for a pure moral doctrine that disdains all self-love, attached
himself in full measure to Kant --- contributed much to the upswing that Prussia took in
the years that followed immediately after the unfortunate war of 1806 and 1807. When one
of the statesmen who had had a part in leading Prussia's regeneration, and who had always
remained faithful to the spirit from which it proceeded (Minister of State v.\ Schön),
received on the occasion of his official jubilee a fine and warm homage from his fellow
citizens, he testified that if he had accomplished anything beneficial for his
fatherland, he owed it to the way of thinking that his teacher Kant had implanted in him,
and that he must therefore carry this thanks that had been brought to him back to the
source of which he was but a little brook. Without, for the rest, being able to make any
comparison, I am seized by a similar feeling when I reflect on what Kant has been to me
and to many others in my fatherland.'' ---

What gives A. S. Ørsted an interesting place in the history of Danish thought is the
significance that the relation between thought and experience --- between universal
propositions on the one side and reality's manifold and shifting content on the other
--- has had in his development, both as philosopher and as jurist. It was his practical
work as judge and legislator, in connection with the study of the legislation of various
countries and with the influence of philosophers such as Fichte and Treschow, that
brought him to lay weight on experience and to see that thought, neither in philosophical
systematization nor in legislation, can exhaust the richness of life. Its significance is
to orient one in life, not to create life. --- I bring forward some of the chief points
that are of interest in this respect.

\begin{center}---\end{center}

It was, according to Ørsted's own statement, in his work \emph{om Trykkefrihedslovgivningen}
[On the Legislation on Freedom of the Press], which appeared in 1801, that his altered
line of thought was first stated. It had from the first been his view that the foundation
of the legal order lies in the necessity that the individual's freedom be secured in such
a way that it does not merely depend on the good will of others. There is required, then,
a power that sets limits to encroachment. But the means by which this limitation is to be
reached must be determined by the people's whole moral condition at the given time, and
no law can formulate rules that can without further ado be applied to the individual case.
The judge must by a discretionary judgment decide whether this falls under the law: ``The
pure fact, as it shows itself to the senses, does not let itself be immediately brought
under a concept to which the law can attach a punishment or another legal consequence, but
the courts must first, by an act of judgment grounded on the consideration of all the
circumstances, recognize that the given fact has the legal quality that the statute
presupposes.''\footnote{\emph{Af mit Livs og min Tids Historie}, I, p.~91. (Ørsted here
himself gives a résumé of the said work on the legislation on freedom of the press.)} In
an essay from the same year \emph{om Regeringens Ret til at ophæve eller forandre
Stiftelser, som private Mænd have oprettet} [On the Government's Right to Abolish or Alter
Foundations Established by Private Men], Ørsted states that in examining any relation
whatever of a moral nature one will convince oneself that universal principles are not
enough to order it, when one is not supported by a sound practical tact and an honest
will.\footnote{\emph{Eunomia. Afhandlinger henhørende til Moralphilosophien,
Statsphilosophien og den dansk-norske Lovkyndighed} [Eunomia: Essays Belonging to Moral
Philosophy, Political Philosophy, and Dano-Norwegian Jurisprudence], I, p.~37.} And later
he returns, as we shall see, more fully to this subject.

An absolute right of property does not exist. The right of property is grounded on the
fact that the individual must have means for free development, for reaching his rational
goal. But by this grounding a limit is at once set. The owner cannot have the right to
make eternally valid provisions for the use of his property. That could lead to the dead
coming to rule over the living. Bequests to foundations can then be acknowledged only in
so far as, in the individual cases, they are found to accord with the common good and the
common will.\footnote{Compare \emph{Af mit Livs og min Tids Historie}, I, p.~118, where
Ørsted refers to the fact that an ordinance of 1781 declared invalid all declarations of
will that could hinder the country's welfare cause, the abolition of the common ownership
of land.}

As one sees, Ørsted seeks to set an experiential grounding of morality and right in place
of pure reason. It is the regard for what, in each individual case, ``the common good and
the common will'' demands that becomes the decisive thing. The expressions can be
understood only thus: that it is society's will, striving toward ``the common best,'' over
which --- by the government --- discretion is to be exercised in the individual cases. The
expressions recall both Bentham and Rousseau, whose ``volonté générale'' really makes
itself known in Kant's categorical imperative. Those expressions are characteristic of a
Kantian who has seen the necessity of considerations of experience in ethics and the
doctrine of right.

In an essay from 1807 \emph{om Religion og Stat} [On Religion and the State] (Eunomia I)
he maintains that all the conditions of a people's welfare that can be formulated as a
general ordering ought to be drawn into the legal order. Thus right serves not merely (as
Ørsted had assumed in earlier works) to secure the freedom of the one against the others,
but acquires a more positive significance. Excluded from the aim of the state is only what
is better effected by the free, self-directed care of individuals. But a completely
reassuring legal mechanism is not attainable. Conscientious consideration of the
circumstances at hand must seek to apply the general laws in an appropriate way.

By the changes that had thus entered into Ørsted's conception of morality and right, he
had broken both with the old natural law and with Kant's and Fichte's theories. But on the
other hand he did not give up the possibility of a philosophy of right. The task for such
a one then became to investigate how legal relations are to be ordered, when regard is
taken both to securing the individuals' freedom and personality and to society's
progressive well-being --- this word taken in a comprehensive sense, so that the spiritual
goods of life too are included under it. Comparative jurisprudence, as well as political
economy, acquire great significance for such a philosophy of right, which will be more
fruitful than the old natural law with its abstract principles and its artificial
deductions. Preliminary works toward such a science Ørsted finds in the works of
Montesquieu, Filangieri, and Bentham. He himself thinks that in his work on the history of
his life and his times he gives a contribution to it.\footnote{\emph{Af mit Livs og min
Tids Historie}, I, pp.~145--150.}

It was in the foregoing works especially the doctrine of right that Ørsted had treated. For
ethics, the remarkable essay \emph{Om Teori og Praxis i Sædelæren} [On Theory and Practice
in Moral Doctrine] (1811) (reprinted in Eunomia I) has a significance similar to that which
the essay on religion and the state had for the doctrine of right. Here too it is
maintained that there is an interval between the universal concept and the individual case
that is to fall under it --- an interval that can be filled only by a sound judgment.

Some years before, Niels Treschow, who was now professor at the university, had published
an interesting essay\footnote{\emph{Det danske Videnskabernes Selskabs Skrifter}
[Transactions of the Danish Society of Sciences], Fifth Volume, 1810.} whose title was:
\emph{Gives der noget Begreb eller nogen Ide om enslige} [i.e.\ single, individual]
\emph{Ting?} [Is There Any Concept or Any Idea of Single Things?]. In this investigation,
directed against Schelling, the author maintains that the individual beings are more than
fleeting images of the Absolute. No universal concept exhausts their nature. It is not the
species, but the individuals, that are the properly real, and there are scarcely fixed
boundaries between the species. It is the limitation of our understanding that makes it so
that we cannot penetrate to the distinctive, the single, and the individual. The universal
concepts are only aids to survey.

To this essay Ørsted refers, though he had already entered upon a similar line of thought
by his own way. It is on the ethical domain that he maintains the significance of the
individual. No universal moral commandments can exhaust what the circumstances demand in
the individual case, and in particular they cannot let the individual's distinctiveness
come into its full right. The spiritual takes on a peculiar shape in each individual human
being. He who would abolish his distinctiveness and seek to transform himself into an
abstract rational being would work at tearing up the root of life and making himself a
marrowless shadow-image. The contribution to the carrying-out of the world-plan that the
individual is to give, he can give only in the form peculiar to his spiritual nature.
Besides the tasks he has in common with others, he has his own peculiar task, his special
ethics. Not only has each individual his calling, but the common struggle against folly and
wickedness must everyone wage in his own way. If the meek man would take on the bold man's
role, and the latter the former's, they could both be sure of speaking and acting in vain.
Unfortunately we have here no means of construction or measurement. It is not even possible
to bring a definite, observationally given individuality under a concept: however long I
keep adding mark to mark, I shall not be able to exhaust the infinite quantity of nuances
and degrees. Individuality sets a limit to our striving to see through the connection
between unity and manifoldness.

No moral theory can set up concepts that are sufficient to solve all the tasks that life's
manifold entanglements bring with them. Thereby the significance of finding ideas and types
does not fall away; but the scientific work must here in the end pass over into an art.
Science can prepare this art, but not produce it. Two different drives work in us: to find
unity in life as in nature --- and to let life unfold itself freely and distinctively.
These two conflicting drives suit well a being whose vocation is a constantly progressing
development, not a fixed possession. ---

In Holberg we have already met a similar line of thought, and we shall meet it later in
other Danish philosophers. The essay on theory and practice in moral doctrine is not only
an interesting contribution to the theory of ethics, but also a characteristically Danish
contribution. As for Ørsted, he states 40 years later in his autobiography that he still
quite acknowledges the essay's content. ---

\begin{center}---\end{center}

Ørsted's literary activity was stopped for a long time in the year 1826. He had, in
\emph{Juridisk Tidsskrift} [Legal Journal], written a contribution to the well-known
ecclesiastical dispute between Grundtvig and H. N. Clausen, under the title: \emph{Behøver
den danske Kirkeforfatning en omgribende Forandring?} [Does the Danish Church Constitution
Need a Comprehensive Change?]. Ørsted answers No. He maintains, against Clausen, that it
was not necessary to alter the symbolic books --- just as little as it was justified (which
Grundtvig demanded) to exclude from an ecclesiastical or theological teaching office anyone
who did not take the Bible and the symbolic books literally. He here brought to bear a line
of thought similar to the one he had maintained in 1798 against the Bible and Reason
parties, and a line of thought analogous to the one that underlay the essay on the
government's right to abolish or alter private foundations. He thought that interpretation
must change with the changing times; literal assent was not necessary, provided only the
essential was acknowledged. As points where a literal conception was not necessary, Ørsted
names: the Fall, the Trinity, inspiration, and the eternal punishments of hell.

King Frederik VI found that a high legal official ought not to take part in the discussion
of a matter in whose treatment before the courts he might perhaps come to take part. Ørsted
had to promise to cease his literary activity.\footnote{On this matter see L.\ Koch:
\emph{Anders Sandøe Ørsted}, 1896, pp.~85--96.} ---

On Ørsted's later political activity I refer to L. Koch's fine biography, as well as to my
characterization of the Ørsted brothers (in \emph{Vort Folk} [Our People]).

In the [eighteen-]thirties he was ``the King's and the People's man,'' and as such
president in the estates assemblies. The turning point in his relation to his countrymen
came, as he himself has stated,\footnote{\emph{Af mit Livs og min Tids Historie}, IV,
p.~XLII.} through his appearance (in the Viborg Estates in 1842) in the matter of the
Danish language's right in the estates of Schleswig. He, who had so strongly impressed the
significance and right of experience and individuality, did not see the significance of the
awakening feeling of nationality and of the national differences. The oppositions --- the
religious, national, and social --- had sharpened in the course of the century, and it was
no wonder that even highly gifted men, who had grown up under the more harmonious conditions
at the century's beginning, could not take the right stand in the strife of these
oppositions. Ørsted's tragic fate as a historical personality has not, however, hindered
the recognition of his work as a researcher. In Bissen's fine statue it is the brooding
thinker who is represented.

\begin{center}---\end{center}

% ============================================================
%  9. F. G. HOWITZ   (Da.: "9. F. G. HOWITZ")
% ============================================================
\dfchapter{9. F. G. Howitz}
\label{ch:df-howitz}

\noindent With a single exception (Joh. Boye), those who philosophized here at home
since Holberg's and Sneedorff's time had been quite predominantly influenced by German
philosophical literature, and this was especially the case in the first decade of the
nineteenth century, when German Romanticism had the say. To this it also contributed
that in English and French philosophy there prevailed an ebb --- after Hume in England
and after Rousseau and Diderot in France. As a refreshing breath, then, comes the
appearance of the physician Franz Gothard Howitz in the year 1824. Howitz (born 1789,
died 1826) was professor of forensic medicine and wrote, on the occasion of a case
concerning an accused woman's imputability, an essay \emph{(Om Afsindighed og Tilregnelse}
[On Insanity and Imputability], a contribution to psychology and jurisprudence) that
awakened a dispute about ``the freedom of the will,'' in which A. S. Ørsted, Mynster,
Sibbern, J. L. Heiberg, and others took part. Howitz's opponents all stood on a Kantian
or German-speculative standpoint, and he found in the course of the dispute repeated
occasion to refer to English and French thought, as well as to Spinoza, whose standpoint
he most nearly makes his own. He calls attention to how much there still is to learn from
those older thinkers, and he thinks that even though one cannot deny German research its
significance, German literature at the beginning of the nineteenth century had taken a
direction toward the obscure and the fantastic that had harmed not only philosophy, but
also other sciences. Our national character is, moreover, different from the German; like
the Dutch, it may be said to form a kind of transition between the German and the
English. But above all it is a matter of letting experience come into its own against
speculation, and he points out that there are signs in Germany itself of a turn in that
direction, also in the philosophical domain.\footnote{\emph{Determinismen eller Hume imod
Kant. Et philosophisk Forsvar for Afhandlingen om Afsindighed og Tilregnelse}
[Determinism, or Hume against Kant: A Philosophical Defence of the Essay on Insanity and
Imputability]. 1824. Preface. --- It is the philosopher Beneke who is meant in the
last-cited words. Compare \emph{op.\ cit.}, p.~76, and \emph{Ultimatum} (1825), p.~50. ---
In \emph{Determinismen eller Hume mod Kant}, Howitz gives a translation of Hume's
excellent essay on liberty and necessity.}

When the dispute between Howitz and his opponents came to turn on ``freedom'' and
``necessity,'' it was nevertheless not this side of the matter that Howitz had laid the
chief weight on in his first essay. His main aim was to impress that there are
intermediate forms and transitions between insanity and ordinary human rationality, and
that there is not in nature the sharp difference the jurists were accustomed to assume.
His line of thought here is really a continuation of Treschow's and A. S. Ørsted's. Man
forms for himself, he says, sharply separated concepts, sets up definitions, and makes
divisions, but in the real world the boundaries of things melt together, and there is a
chain-connection brought about by continual transitions. The systematic naturalists often
find it difficult to fix the boundaries between different species or groups, as for
example between plants and animals; but the difficulties must become still greater where
we --- as where it is a matter of man's psychological nature --- do not have quite good,
tangible marks to hold to. There are many degrees between freedom and unfreedom, between
the human rational normal state and insanity, between imputability and non-imputability.
These oppositions hold only when one takes cases that are extremes in different
directions, but not for the intermediate cases. There are thus many degrees of insanity,
especially when one takes into account that insanity does not always consist in an actual
lack of reason. There can be morbid states where reason is in no way defective, but is
prevented from exercising its influence; --- fixed affects there are just as well as fixed
ideas. There can be suddenly bursting impulses that stand in the sharpest opposition to
the individual's whole character, and that it cannot check, although it may be clearly
conscious of its state. Of this kind is the \emph{mania sine delirio} described by Pinel
and the so-called \emph{folie raisonnante}, where the patient with great shrewdness knows
how to conceal his state and to answer all objections directed against his notions and
plans. And as there are many kinds of mental illness, there are also many transitions
between insanity and ordinary human reason. The lucid intervals bear witness to it;
likewise the lower degrees of all kinds of insanity, and the insanity that is limited to
quite particular domains, as well as the initial stage and the convalescence as regards
the individual mental illness. In ordinary people's half-sleep, in the morning reveries
continued with half-bound consciousness, which are so alluring especially for nervously
weak people, we also have transitional states; in this morning twilight, images of terror
from the dream can be transferred to persons actually present and lead to attacks on them.
Finally, so-called normal people's passions and follies often present features that make
it difficult to set a sharp boundary between the healthy and the sick.

And as insanity has its degrees, so too has freedom, when one takes this word in the
simple sense of the ability to do what one wills or wishes, independently of foreign
determination and foreign will. In this sense freedom has an undisputed reality --- but is
present in the individual cases in a highly different degree. It will depend on the degree
to which a person can raise himself above particular, momentary impulses. As everywhere in
nature force stands against force, so here too: a living being's self-activity is
different from the external, mechanical necessity; against animal desire stands rational
deliberation; against the egoistic drives stand the feelings that lead to work for the
good of the whole. The relation of strength between the various forces is highly different
in the individual cases and in the individual persons.

Howitz found that the current juristic conception, against which he came forward, suffered
from a false psychology, which in the last generation had been reinforced by Kant's
philosophy. It was especially the life of feeling that did not come into its own. For
Kant's strict morality of reason, all feeling, all pleasure and displeasure, all joy and
sorrow, all fellow-feeling and sense of honour, stood as ``sensuous'' or ``egoistic.''
Howitz impresses that feelings can be of highly different value, and that one cannot
without further ado put the joy of sympathy on a level with the pleasures of the palate.
Every feeling is connected with a need or a drive, and the individual's development depends
on the relation between the various drives. However exalted and noble an action may be, it
nevertheless becomes possible only through a feeling or drive conquering others. Spinoza is
right that a passion can be driven out only by another passion. The life of feeling has its
own laws, and reason --- the ability to judge of the relations of things --- is only
serving; it decides what means are necessary and what ways can be gone, but it does not
decide what becomes the final aim. Reason is the helmsman with his rudder and his
sea-chart, but feeling is the wind that fills the sail. Pleasure drives the work, interest
develops --- indeed, creates --- the genius, and the role that a person is to play in the
world depends more on his ruling character and temperament than on the logical perfection
of his faculty of thought.

Through the life of feeling and drive we come, according to Howitz, back to man's natural
ground. There is a connection between all expressions of psychic life, even though they
are of highly different value. This connection is indicated by the fact that the drive for
felicity stirs at all stages; the difference between the stages depends on the difference
between the objects by which the drive for felicity is satisfied. --- On this point Howitz
was, not without his own fault, misunderstood by his opponents, who, although they
distanced themselves from Kant's psychology, by drive for felicity thought merely of
egoism or the need for sensuous enjoyment. But his chief endeavour was to maintain the
lawfulness of psychic nature. And while the Romantics, such as Steffens, wished to read
nature from above downward, Howitz wished to go the opposite way,\footnote{Already Johannes
Boye in \emph{Statens Ven} and Treschow in his \emph{Historiens Filosofi} [Philosophy of
History] (1811) had taught a great process of development from the lower organic forms to
the highest. More scientifically, a couple of years before the appearance of Howitz's
writings, the political economist Chr.\ Olufsen, in an essay \emph{Om Menneskets Rolle i
den fysiske Verden} [On Man's Role in the Physical World] (1822), put forward the view that
``the life-force'' in nature works its way forward, from the strata and forms whose time is
past, to new forms better suited to the altered conditions. Although Olufsen does not seem
to assume a continuity between the various organic forms themselves, his presentation
nevertheless has interest for the history of the idea of development. --- Howitz can hardly
have known Olufsen's essay, but his conception could lead in the same direction.} to regard
man as a being of nature, sharing in the qualities that all organic beings possess, and to
ascend from the simplest, most elementary actions to the more composite and developed
expressions of will, constantly seeking the definite acting forces or motives.

It is in reality a sound and capable empirical psychology that Howitz brings to bear,
influenced as he is through the study of Spinoza and Hume. He was fully the equal of his
opponents, although some of these, especially Mynster, wished in a very little attractive
and still less well-founded way to make a name for themselves at his expense. --- With
great clarity he expresses himself, for example, on the relation between the spiritual and
the material, on the occasion of his having been reproached for inclining to materialism.
He holds to experience, which shows man equipped with both bodily and spiritual qualities.
Only by abstraction can one separate the spiritual from the material, or conversely. The
two kinds of qualities rise and fall with one another; but it is not correct to say that
there is here a causal relation present, since we can derive the spiritual from the
material as little as the material from the spiritual. We have here a being with two
different qualities; but because two qualities are different, they need not belong to two
different worlds. It depends perhaps on our faculty of conception that we must make such a
difference between the qualities. I find, says Howitz, no satisfactory philosophical view
of this matter except the identity-system, which declares body and spirit to be at bottom
the same substance, only that this presents itself in another way to the inner than to the
outer sense. --- Howitz refers not only to Spinoza, but also to Treschow, who had stated a
similar view, which we shall treat in the next section. Sibbern too embraced such a view.

\begin{center}---\end{center}

It was, as already remarked, not Howitz's intention with his appearance to teach
determinism, but to demonstrate that the question of imputability was independent of
speculative theories. This he wished to achieve by two ways. In the first place he assumed
that when there is present in experience a quantity of transitions between insanity and
spiritual health, and between weakness toward momentary impulses and receptivity for
reflection, one will not be able to point out any absolute boundary between freedom and
unfreedom. He wonders that Ørsted and Sibbern, those of his opponents for whom he has the
greatest respect, can acknowledge those many degrees and yet think themselves
indeterminists. --- In the second place he maintains that a jurist like Ørsted, who grounds
the justification of punishment on the necessity of deterrence, really neither needs nor
can use the assumption of a causeless volition. It depends on whether the individual's will
can be influenced by regard for the punishment; the justification and necessity of
punishment lie in the fact that a counterweight, a counter-motive, is to be brought about
against the motives that can lead to crime. Punishment is a means of prevention. It does
not then depend on whether the individual, at the moment of the action, could have resisted
the impulse; punishment aims at calling forth the motives that have shown themselves not to
be present. Only him who cannot be influenced by regard for punishment would it be
meaningless to punish.

When Howitz declares himself against indeterminism, it is not least because it makes
punishment and imputation impossible. Freedom --- in the sense in which Howitz assumes its
possibility --- is receptivity for motives and for reflection. He denies only freedom in
the sense of the ability, under the same inner and outer conditions, to make two different
decisions equally well (bilateral arbitrariness). But he calls attention to the fact that
those who defend indeterminism use the word freedom in three or four senses, which does not
benefit clarity. Particularly frequent is the confusion of causal freedom and strength of
will.\footnote{In my \emph{Etik} [Ethics] (chap.~5) I distinguish between six different
meanings of the word freedom.}

The word necessity too has different meanings. It does not mean the same as compulsion.
Determinism (or, as Howitz somewhat unhappily also calls it, fatalism) is not the same as
Turkish fatalism. Turkish fatalism lets, like indeterminism, that which happens be
independent of all preceding causes, while determinism maintains that we, through
understanding of the causes and application of this understanding, can prevent the evil and
further the good. Only thereby do education and formation of character become possible.
When knowledge of men is so difficult to attain, it is not because the causal law does not
hold in human life, but because here so many different causes work together.

\begin{center}---\end{center}

In his ethical conception Howitz essentially attaches himself to Hutcheson and Hume, in
that he sees in the sympathetic drives the foundation of the moral. These drives are formed
and clarified by the understanding of what the existence of social life demands. Therefore
a developed cognition is required in order that one may act morally. But with the
satisfaction both of the sympathetic drives and of the need for cognition there is
connected a feeling of pleasure, and in so far the drive for happiness may be said to be
the fundamental drive. If a mother did not feel greater satisfaction in --- despite her own
weariness and privation --- tending her child than in abandoning it, then mother-love would
be impossible.

As already remarked, in Howitz the drive for felicity and egoism do not coincide. The drive
for felicity can take on forms that are among themselves highly different --- indeed,
conflicting. Thus man's love of the good or of the generally useful, and his inclination to
strive out beyond his finite existence, are a peculiar kind of motive, different from the
sensuous impulses. ---

The views that Howitz brought to bear in the psychological and the ethical domain found no
assent in our philosophical literature. He stood as good as alone with them. Among his
opponents, Sibbern was the one who stood nearest to him, and Sibbern's later development let
this kinship come forth more strongly.

After the philosophical single combat [Holmgang], Howitz withdrew to his own domain, where
it was unfortunately not long granted him to continue his activity.

\begin{center}---\end{center}

% ============================================================
%  10. NIELS TRESCHOW   (Da.: "10. NIELS TRESCHOW")
% ============================================================
\dfchapter{10. Niels Treschow}
\label{ch:df-treschow}

\noindent Several times in the foregoing Treschow has been named; he was the one who
most thoroughly presented and assessed Kant's philosophy, and both A. S. Ørsted and
Howitz refer to his works. It is now the place to give a connected characterization of
him.

Treschow was born near Drammen in the year 1751 and became a student in Copenhagen when
not yet fifteen years old. His urge to think for himself was early awakened. ``I
followed,'' he says in his autobiography, ``Descartes' rule, without knowing this
philosopher, of beginning my scientific career with doubt about everything I had
learned. These doubts I hoped, however, at length to get resolved --- for I saw that
learned men no longer doubted.'' After very extensive studies, which in philosophy
comprised, besides the reigning Wolffian school, also the English philosophers of the
eighteenth century, he went again to Norway as a teacher at the school in Trondhjem.
Later he became headmaster in Helsingør and from there came again to Norway as
headmaster in Kristiania. Here he delivered, in 1796--1797, his lectures on Kant's
philosophy. It was probably chiefly these that brought it about that in 1803, when
Riisbrigh took his leave, he became, without application, professor of philosophy in
Copenhagen. The government would not have Steffens, that romantic visionary, as a
teacher for the young. The choice that was made was exceedingly fortunate. Treschow was
a lively and independent man, who exercised a great influence on the students\footnote{Carsten
Hauch has, in his \emph{Ungdomserindringer} [Reminiscences of Youth], given a
characterization of him as a teacher.} and was missed when in 1813 he went to
Kristiania's new university. He also gave lectures to larger circles with great
attendance. --- In Norway he later became a councillor of state (minister of religious
affairs) and lived his last years withdrawn and occupied with philosophical works until
his death in the year 1833.

\begin{center}---\end{center}

Treschow composed the first Danish empirical psychology: \emph{Om den menneskelige
Natur, især fra dens aandelige Side} [On Human Nature, Especially from Its Spiritual
Side] (1812). This book, which still has its interest, bears witness to a fine faculty
of observation and to a warm interest in the life of the soul and its development.
Although he is still in so far the child of the eighteenth century that he regards the
life of representation as the fundamental thing, so that all drives are to have their
origin from ``obscure representations,'' he nevertheless sets feeling and will as
independent sides of the life of the soul over against representation. The section on
the life of feeling is especially rich in good observations and remarks, and here again
especially the section on sympathy, admiration, enthusiasm, and similar feelings.

With regard to the relation between soul and body, Treschow lays Spinoza's conception at
the foundation. ``Man may indeed be regarded as a composite being, but not as composed
of soul and body, for both are only opposite sides of the same thing, in so far, namely,
as it can be an object for the inner and the outer sense.''\footnote{In his
``philosophical testament'' (\emph{Om Gud, Ide- og Sandseverdenen}, Kristiania 1831, II,
p.~51) Treschow says: ``The soul's thoughts and feelings would, if it were possible to
observe them by the outer senses, themselves appear to us as certain kinds of motion;
what they are beyond this, the inner sense teaches us, in the same way as by sight and
hearing we discover qualities of which by smell and taste we do not even get any
inkling, and conversely.''} The materialists assume only the existence of the bodily,
the spiritualists only that of the spiritual. Both are in the wrong toward experience.
Of neither matter nor soul have we immediate experience; it is in both cases by
inference that we reach, from the shifting phenomena, the thought of what lies at the
foundation, and it is in the end one and the same principle in which the ground must be
sought.

By this way of seeing, Treschow was able to give the scientific conception its full
right, without any romantic reinterpretation. Psychology could preserve its independence
as proceeding from a peculiar point of view of the same existence that natural science,
especially physiology, regarded from its point of view. It is on the very fundamental
principles of natural science that Treschow supports himself: cause and effect cannot be
different in essence; a motion can call forth only another motion; if there seems to
arise more in the effect than one can derive from the cause --- for example, sensation
besides motion --- then the explanation must lie in the fact that we have only
incompletely and one-sidedly conceived the cause.

It was this conception that Howitz attached himself to. It also underlies Sibbern's
psychology, as well as my own presentation of psychology. There is good Danish tradition
for it. Treschow seems on this point to be influenced by Hartley, besides by a direct
study of Spinoza; but his own thinking made the decision. ---

\begin{center}---\end{center}

Perhaps the most remarkable of Treschow's works is the already provisionally mentioned
essay (in the Transactions of the Society of Sciences, 1810) that discusses the question
whether concepts of individual objects can be formed. \emph{(Gives der noget Begreb
eller nogen Ide om enslige Ting?} [Is There Any Concept or Any Idea of Single Things?])
It is directed against speculative philosophy's inclination to move in universal ideas
and to regard the individual phenomena as imperfect expressions of them. Treschow
emphasizes that there is contained more in the individual objects than what can enter
into a universal concept. The special, concrete form or type that the individual
presents must be just as much a task to know as that which is common to many
individuals. In the real world there are only individuals, and our universal concepts
are really only aids to coming to know them. But since an individual, in its whole
particular distinctiveness, presents an inexhaustible content, it is a very difficult
task to find the idea that expresses what lies at the foundation of its shifting
states.\footnote{When Treschow here speaks of ``an idea that is not universal in the
respect that it comprises several things, but in so far as it expresses all possible
states of the same,'' he means the same as what I have in my \emph{Psykologi}
[Psychology] (V\,B, 9) called a typical individual-representation. --- Compare what H. C.
Ørsted called ``co-foldness'' [\emph{Samfoldighed}].} It is especially in the domain of
organic life and of history that this problem is posed. And there is a corresponding
practical problem: our value (``the whole excellence of our being'') depends on
individuality, and the aim of our striving must be to attain the highest degree of
individuality.

When individuals are the properly real, not mere means or forms for something universal,
it becomes a question what validity can be ascribed to the species-concepts. Should
nature really have formed fixed boundaries between the species? Are the species-concepts
more than aids we use to procure a survey for ourselves? Is it not the imperfection of
our understanding that makes such aids necessary, because we are not able to penetrate
into the whole distinctiveness of the individual? ---

Nearly contemporaneously with these thoughts appeared \emph{Lamarck}'s critique of the
species-concept (\emph{Philosophie zoologique}, 1809). The species-concept is for
Lamarck an artificial means that one must not reckon among ``the very laws and
operations of nature.'' Treschow did not know Lamarck, whereas he refers to older French
naturalists (Maillet, Buffon), whose ideas already went in the same direction. ---

When the species-concept is criticized in this way, it does not lie far to assume that
the differences between the various species have arisen through successive development.
And this thought also plays a great role in Treschow. It appears in his \emph{Elementer
til Historiens Filosofi} [Elements toward a Philosophy of History] (1811), which he first
presented in lectures in the winter of 1806--1807. On this point too Treschow entered
into opposition to Romanticism. When the latter spoke of development, it was only
symbolic: it is in the play of forms, in the fantasy of nature or of the spirit of
nature, that the transitions take place; the eternal is at once expressed in infinitely
many forms; but a historical transition, a development in time under definite external
conditions, is not assumed. Steffens, indeed, thus expressly denies a real development
in nature. Treschow, on the other hand, attempts in his way to give a natural history of
creation, in that he interprets the first book of Moses with the same boldness as
Holberg. Creation is a long process, advancing from stage to stage. An ``immediate
creation'' is an empty word. ``From materials that seem raw, because one becomes aware
of no form in them, there takes place an imperceptible transition to more definite
shapes\ldots\ Nature's course in those deviations or degenerations by which the previous
races and species are transformed into quite new ones is not only creeping, but so slow
withal that neither one's own experience nor history's all too little compass puts us in
a position to follow its almost imperceptible traces. After the lapse of millennia, the
least ambiguous remains of animals and men that can be regarded as ancestors of the
present, though fairly different, nevertheless give occasion for objections and doubts.''
--- Man has only gradually acquired the qualities that now characterize him. Man's
natural history is a part of the whole history of nature. He has for centuries led an
animal life, even of a highly imperfect kind; his development was due to an interplay
between inner and outer conditions. Before we could lift our heads to heaven and feel
ourselves akin to it, not only had our organism to be formed, but also the surface of
the Earth, the plant- and animal-world, had to undergo changes.

``Have we perhaps,'' asks Treschow, ``cause to be ashamed of such an origin? This
question can be answered quite briefly. We need not go so far back in time to find
stronger grounds for our humiliation. How many years is it since each of us stood far
below the animals, indeed even below the life of the meanest plant? Is it not more
important to think of what we now are and one day quite certainly shall become, than of
what we were --- as if the proud cedar would look down with scorn on the young shoot that
springs up at its roots, or we on the foetus whose heart has barely begun to beat in the
mother's womb? Everywhere the living has the same root: divine is our descent, eternal
and infinite our vocation.'' ---

Treschow does not have the idea of the struggle for existence, which, on the other hand,
as mentioned above, is found in Joh. Boye and comes forth some years later in Chr.
Olufsen, in an essay that in an interesting way stands between the romantic and the
Darwinian conception.

As is already seen from Treschow's cited statements, he lays weight on the analogy
between the individual and the development of the whole race. Both physically and
spiritually the individual goes through the development of the race. He stretches this
analogy very far and finds thus that the race, like the individual, has had its foetal
state, its childlike state of innocence, its boyhood, and so on. In the end it will also
get its old age. --- Such an analogy is, of course, arbitrary, and it fails at the
decisive points. Only figuratively can the race be conceived as a single individual,
already for the reason that in history there simultaneously take place different
processes that cannot be subordinated into a single progressing series. One must then in
each case --- as Sibbern later did --- lay weight on the fact that development takes
place sporadically from many different points at once. But thereby the problem itself is
pushed further out, since the great question becomes whether and how the many processes
can be united in one gathered stream. ---

\begin{center}---\end{center}

Behind the special theories there lies in Treschow a religious-speculative doctrine of
unity. The unity in everything is for him the last consequence of thought --- the point
we reach when we have step by step followed the series of grounds and consequences, but
where the proposition that there is a ground for everything can no longer be applied,
because it holds only for the manifold and the different.

Against Kant, Treschow had maintained that we can connect in our thought only what is
already connected in nature. It is the contradictions and difficulties that appear in
experience that drive us forward to the idea of a highest unity. It is the necessity of
connecting things that demonstrates the validity of this idea. From the world of sense
we thus come to the world of ideas, and from this to God. But it is not three different
worlds that we here have before us. ``The forces of nature are not different from the
Godhead's own, except in so far as this, regarded in itself, is infinite, but reveals
itself in those in a finite way, and in its infinity only as a ceaselessly increasing
magnitude. In this sense one can say that matter has formed itself, has gradually
advanced from the less perfect to the most perfect human form. One can likewise say with
truth that God, after having brought forth matter itself, or by its qualities made his
own being as intuitable as it could become, has from it again brought forth all things.
These turns of phrase, which most people hold to be poles apart, nevertheless really
mean the same thing and only make known two human modes of representation.'' --- Like H.
C. Ørsted, Treschow was convinced that what philosophy calls the highest unity is the
same as what religion calls God. An immediate revelation he did not assume. Not even in
his work of old age \emph{Om Gud, Ide- og Sanseverdenen} [On God, the World of Ideas,
and the World of Sense] (1813), where he develops his views in a somewhat dogmatic way
and takes more account than before of theological representations, does he abandon his
purely philosophical standpoint, which had, after all, such depth and breadth that it
could let scientific hypotheses and psychological observations unfold themselves freely,
without needing to fear coming into conflict with them. ---

After Treschow's departure from Copenhagen, it was the aesthetic and theological
questions that for a long time occupied interest. But Treschow's successor at the
university, of whom we shall now speak, worked essentially in the same spirit as he,
although he too was strongly influenced by the spirit of Romanticism and only by the
constant striving for health and truth that marked his long life worked his way out of
its one-sidednesses.

\begin{center}---\end{center}

% ============================================================
%  11. FREDERIK CHRISTIAN SIBBERN   (Da.: "11. FREDERIK CHRISTIAN SIBBERN")
% ============================================================
\dfchapter{11. Frederik Christian Sibbern}
\label{ch:df-sibbern}

\noindent When Treschow went to Norway in 1813, Sibbern (born 1785, died 1872) became
his successor. Baggesen had thought of becoming professor of philosophy on this occasion
and expresses himself, some years later, very resentfully at having been passed over for
a ``witless boy.''\footnote{See Kr.\ Arentzen: \emph{Baggesen og Oehlenschläger}, VI,
p.~35. --- That Sibbern came forward on Oehlenschläger's side in the great dispute did
not, of course, make Baggesen milder.} This boy was, however, to grow to manhood and
become one of the most distinctive figures in the history of Danish thought.

Sibbern's position in our intellectual life is characterized by the fact that he, who in
his youth was strongly gripped by the spirit of Romanticism, both in a philosophical,
religious, and poetic respect, through tireless work with thought and with his whole
personality opened himself to new ways of seeing and continued his development right up
into old age, where for most it is once and for all concluded. It was especially his
psychological sense and his warm sympathy that opened his eyes to the realities of life
in their manifoldness and distinctiveness, each for itself. In place of Romanticism's
speculation entered a view of the world grounded on empirical psychology and natural
science; in place of the orthodoxy to which, from disgust with rationalism, he had
attached himself, there developed a free religious view of life; in place of
Romanticism's conservatism entered a constantly growing sympathy with the striving of
the people, especially of the ``lowly folk,'' and he was one of the first who here at
home had a sense for the social question and spoke out for a radical solution of it.
Common to all stages in his development was the lively sense for all that is human, the
openness, and the precious naïveté. In his outward bearing there was a certain
awkwardness, and there came besides, early on, a pedantic and stiff form into his speech
and writing, a shell that for many concealed the kernel.\footnote{``For me,'' he says in
a letter, ``there has always been a great step from what stood before my inner sight to
getting under way with bringing it forth before me in a satisfactory manner.'' \emph{Breve
til og fra Sibbern} [Letters to and from Sibbern], II, p.~230.} The liveliness in his
writings was hindered --- especially in the later editions --- by a strong inclination to
divisions and classifications.

His own experiences and his own temperament became in several respects decisive for the
psychological conception that played so great a role in his philosophy. In his youth he
had a period of ferment, in which a hopeless love was the most strongly acting element.
He learned here from his own experience to know the conflict between rising desire and
admiring joy. In the first part of \emph{Gabrielis' Breve} [Gabrielis's Letters] he has
depicted this crisis, often with the same words that are found in his own letters from
the same time. In his mind bitterness and dejection constantly strove with cheerful
optimism, and he thereby felt the opposition between the times when spiritlessness and
barrenness prevail and the times when ``the Muse comes to visit.''

His greatest influence on the students he probably exercised in the [eighteen-]twenties
and thirties. For those who came into contact with him at the close of his 55-year
university activity, it was often difficult to penetrate to the kernel behind the
somewhat stiffened bearing. And yet the kernel was constantly alive. Even at 86 years old
he stated that he felt summeriness in himself, and that one could not, indeed, add a
cubit to one's stature, but perhaps an inch, when one holds oneself erect!\footnote{In my
\emph{Mindre Arbejder} [Minor Works] (First Series) I have given a fuller characterization
of Sibbern's personality. As regards his doctrine too, I can refer to that presentation;
the one given in what follows is on some points fuller, on others more concise. In that
characterization the influence of Steffens is especially emphasized. To this must be added
Treschow's influence, which has not been slight, and which Sibbern himself repeatedly
emphasizes. And both influences went in the direction of the significance of the individual
as opposed to the abstract universal, the significance of the real as opposed to the
formal.}

Sibbern was first and foremost a psychologist. He had a great faculty for observation and
great joy in making observations in the psychic domain. He compares his psychologizing
with a botanizing. And as for the botanist there is no weed, so the psychologist is
interested in all that is human, whether it be great or small, good or evil. He was
convinced of a close connection between the psychic and the bodily. It is one and the same
fundamental activity that appears as psychic and as bodily, one and the same life that
expresses itself under two forms. If one would assume that from the first there took place
only a bodily process of motion, one would not be able to understand how it in its
continuation could produce thoughts. The bodily can produce only a bodily thing.
Conversely, one cannot understand how thoughts can produce bodily motions, and the psychic
can therefore likewise not be the original, but must have its ground in something deeper,
which is at the same time able to produce motion. Out of the same fundamental source or
fundamental activity by which the material formations take place, the spiritual activity
too takes place.\footnote{\emph{Om Forholdet mellem Sjæl og Legeme} [On the Relation between
Soul and Body]. 1849, pp.~84--86. --- \emph{Spekulativ Kosmologi} [Speculative Cosmology].
1846, p.~41. --- In Treschow and Howitz we have already met a similar conception. Sibbern on
the whole refers repeatedly to Treschow as his predecessor. See the preface to
\emph{Psykologisk Patologi} (1828) and several statements in his work of old age
\emph{Meddelelser af et Skrift fra Aar 2135}. --- What places Sibbern in a certain
opposition to the science of our time is his vitalism, that is, his doctrine of a peculiar
life-force, of an activity coming from an inner source. Yet he again sets this life-force
in inner connection with the force acting in all inorganic nature, thus already in the
primeval nebula.}

The life of the soul, like real life, is placed in constant change and becoming. Life
constantly works its way forward to new forms; not only from the first, but throughout the
whole course of life, work is done in the deep ground of the soul; new dispositions and
impulses arise. Of this bear witness not only what is of genius, but also the sudden
awakenings that can arise from unconscious origins. It is life that will come forth in us.
There is much that belongs to the life of the soul and that works before consciousness
awakens. The life of the soul is, just as little as bodily life, formed once and for all;
the process of soul-formation constantly goes its way, whether we immediately notice it or
not.

The soul does not exist as a being for itself, apart from the life of the soul, the psychic
functions, any more than thought exists apart from the process of thinking. The words ``in
itself'' can only mean: ``in its fundamental relations.'' Fichte rightly rejected the
notion of a soul, when one thereby understands a thing, a something that \emph{has} life or
\emph{has} consciousness. It is at bottom a kind of materialism that expresses itself in
such a notion; one makes ``the soul'' into a being that could be met with somewhere in
space, for example somewhere or other in the brain.\footnote{Compare the essay \emph{Om den
menneskelige Sjæls Subsistens i Liv og Død} [On the Subsistence of the Human Soul in Life
and Death] (\emph{Filosofisk Arkiv og Repertorium}, 1829).} The living or acting in which
the soul consists can constantly undergo changes, and the causes of these changes need not,
as said, be found in the conscious life of the soul. Throughout life the soul can grow by
influences from the deep natural ground that bears it. --- Sibbern's conception here recalls
what especially William James in the most recent time has brought to bear: that new
increments can enter the life of consciousness from something that lies beneath the
threshold of consciousness (the subliminal).

But the development that thus takes place does not take its start from a single point in the
soul, so as thence to spread to the rest of the life of the soul. There stir at once
different, often mutually conflicting dispositions and impulses, ideas and drives.
Development takes place \emph{sporadically}, and it then depends on whether a whole can come
out of it. Both the ``lower,'' that which ought to be driven out, and the ``higher,'' that
which ought to be furthered, arise first as a flash, a spur, a captivation. And from this
there follows then stir in the mind, ferment, and often a grievous inner struggle or debate.
It is thereby not purely individual forces that are in motion; for the living and acting in
which the soul's essence consists is a part of the universal activity that stirs in all
life, in the whole of nature. And both of the purely individual and of the universal it
holds that it appears sporadically and often conflictingly, before unity and harmony can be
reached.

``The more deeply life opens itself in a temperament, the sooner it can come to pass through
a chaotic mixture and stirring of the elements in its interior, and the world that is to
build itself up in it, it must now itself take part in creating, as if it were to be
re-created from the ground up. And since herein not merely the content of life in general
stirs according to its essence --- yet in its various fundamental activities as from
different sides and as in conflicting elements, which first seek the harmony --- but now
also the I-ness stirs according to its I-ness, so it is quite natural that in the inner
struggle a certain morbid entanglement in oneself, with alternating self-confidence and
self-doubt, can easily gain the upper hand in certain periods of life as a natural critical
transitional state. This must then be passed through with confidence and self-encouragement,
but also with resigned biding and devotion to the all-embracing all-activity, which just
here shows its power in the individual, and whose vessel one just here feels oneself to
be.''\footnote{\emph{Psykologisk Pathologie} [Psychological Pathology]. 1828, p.~223. ---
This work forms the second part of the first edition of Sibbern's psychology. It was
reprinted in 1885 (the centenary of his birth) under the title: \emph{Læren om de
menneskelige Følelser og Lidenskaber} [The Doctrine of the Human Feelings and Passions]. The
word ``Pathology'' here means, namely, not the sick life of the soul, but on the whole the
life of the soul as it appears in feelings, drives, and emotions. --- It is the fullest and
liveliest psychological presentation that Sibbern has given.}

This statement shows Sibbern's distinctiveness: personal experience has opened his gaze to
an essential side of the conditions of psychic life, and it is then laid at the foundation
of a general theory; at the same time it is not merely the psychologist, but the practical
philosopher of life, who has the word.

While in the cited statement he distinguishes between the universal activity and the
I-ness, he dwells in other places precisely on such states where both go into one, as in
instinctive and inspired activity, or in full devotion to and absorption in a cause. Here
there can be no talk of motives, while one may well speak of movements of feeling. Only
where one does not go up into immediate activity --- when, that is, the action does not
``understand itself of itself'' --- only there does there arise consciousness of motives as
something different from the will. Full immersion is then not present.

The life of feeling Sibbern depicts best. He has a clear eye for the significance of
contrasts, for the opposition between desire and admiration, for the conflict of different
feelings and the close connection they could enter precisely by means of the conflict. In
the last-named respect he distinguishes between feelings' ``mixture'' and their
``mixedness.'' When two objects (or an object's two different sides) at once set the soul
in two opposite movements of feeling, there occurs either an external mixture of the two
feelings, in that they as it were lie side by side, or the mind swings between them; in both
cases the conflicting elements are still unreconciled. It is otherwise when the opposition is
dissolved into unity. Then there arises a mixed feeling, in that the two elements work so
into one that together they produce a single total-feeling. Thus when cheerfulness is won by
the overcoming of fear, or courage and desire are heightened by the very resistance that is
met. It is not the mere contrast that acts; it is a new feeling, to which both the
contrasting feelings have given their contribution.

\begin{center}---\end{center}

What Sibbern had emphasized as distinctive of psychic development, he found again in the
whole of existence. The idea of sporadic development is a chief thought in him and of
decisive significance for his whole view of life, or for his whole mood toward life.

Sibbern expresses himself against every attempt at an a priori construction of nature, among
others against Kant's attempt in this direction.\footnote{In particular Sibbern objects to
Kant that he conceives gravity (attraction) as a fundamental force, although the possibility
is not excluded that it may be the result of a process, may spring from something far more
complicated than is seen in the simple phenomenon. \emph{Spekulativ Kosmologi}. 1846, p.~36.
--- This work contains in broad outline Sibbern's general view of the world. His conception
of philosophy he has presented in his work \emph{Om Filosofiens Begreb, Natur og Væsen} [On
the Concept, Nature, and Essence of Philosophy]. 1843.} We are still, he says, far too far
from having penetrated so deeply into the study of nature that an endeavour to construct for
oneself a system of nature a priori, merely according to the fundamental ideas to which we
have now come, could be anything but highly precarious. If one is to reach a philosophy,
this must develop under constant comparison (conference) with the given. It is the goal, not
the beginning. Even though that which acts in the whole of existence also acts in our
thought, this our thought can nevertheless develop only through interaction and conflict
with the other elements in existence. Our thought is only one of the many elements of
existence. Here there shows itself precisely \emph{the sporadic}.

Experience bears witness in all domains that everything in nature from the first arises
sporadically. In the primeval nebula, from which our planetary system has developed, and
which perhaps is itself the last, purely dissolved residue of an earlier system of worlds,
effects went out from different parts, and through this interaction development advanced. In
crystallization --- for example in the formation of ice --- needles shoot forth from
different points so as at last to form a firm connection. During foetal development the
process of organization, especially the formation of bone, begins from different places and
then meets so as to join together into a whole. In the life of the soul there move, as we
have already seen, many different drives, impulses, and thoughts, if not chaotically, then
at least by different ways, until the right harmonious absorption into totality becomes
possible. In a similar way it goes with men in society, and with the various societies
within humanity.

Through the debate or struggle of all against all occasioned by the sporadic, development
now advances. The condition for a thoroughly organized and rich living is that it has
unfolded itself successively, and that it has arisen through the harmonization of
sporadically arisen elements. Time and separation are thus essential presuppositions for
development. --- Sibbern here enters into sharp opposition to Romanticism's philosophy, which
assumed only a development in a purely logical sense, and for which development was an
unfolding of an absolute principle without the presupposition of real differences. With the
emphasis on the sporadic is connected Sibbern's great sense for the individual: each single
individual is the starting point for new, distinctive processes, is one of the elements of
existence, a contribution to the great debate of life. By the emphasis on time's decisive
significance for all that exists, Sibbern shows himself as an eager philosopher of
development. He applies his thought of development also to man. Humanity has developed out of
animality. A certain animal formation has, from a lowly beginning, through a series of
generations, within which the process of formation was constantly carried further, reached
the point where now its last offspring became a being that, so to speak, opened the eye of
spirituality. The process that has thus taken place had again its origin from the first stir
in the primeval nebula, from which advance was made to a point where the development of life
could begin.

A first beginning we do not trace, and of a final end we can speak only for the part of
existence in which we most nearly find ourselves, and which we know. And this very part is
probably still in its first stage, on the way to its proper beginning! --- An absolute final
end would presuppose a time whose infinity corresponds to the inexhaustibility of the
content of existence.

Sibbern's doctrine of development points back to Lamarck; but in the significance he ascribes
to the great debate or struggle of all against all, he points forward toward Darwinism and
the theory of mutation. In Danish literature he has Boye, Treschow, and Olufsen as his
predecessors. Indeed, already in Holberg there showed themselves hints of a great process of
development in existence.

\begin{center}---\end{center}

Characteristic of Sibbern is both his conviction that new experiences can constantly be made,
and his no less firm conviction that, despite the scattered and chaotic in the emerging
phenomena and processes, harmony will develop between them. The long, restless transitional
periods, which often present the image of something distorted and abnormal, can awaken doubt
and put patience to the test. But the great debate of life is, after all, the condition for a
powerful formation of unity. ``The richer the finitude, the richer the infinity; the greater
the manifoldness, the fuller the unity! Everything separates itself so as in a so much more
intensive way to go together again into unity!''

Here Sibbern's philosophical conception goes into one with the religious. In the
presupposition of an inner connection in all things, on which all science and all scientific
striving rests, he finds testimony of an ideal principle, a fundamental being, a \emph{natura
naturans}, as the innermost present and underlying thing in all things. This all-active being
bears in itself a creative deep, out of which ever new things well forth, which is especially
seen when the result of a composite process contains more than follows from the co-operation
of the elements. The word God designates this omnipresent and all-active fundamental being,
in so far as it is thought as centralized in itself and becomes an object of reverence and
confidence. There is on this point a limit to Sibbern's philosophy of development, for he
cannot conceive that God --- who is, after all, according to his conception, a personal being
--- himself develops. The all-active is eternal and unchangeable; it is only existence's
world-side, not its Godhead-side, that develops. It is, from Sibbern's own standpoint, an
inconsistency; the analogy with his conception of the concept of the soul ought to have led
him to see the untenability of this distinction between the unchangeable and the changeable.

As regards the relation between philosophy and positive religion, an interesting development
took place in Sibbern. In his youth he was strongly taken by positive Christianity. He found
rest and refreshment in it for his soul, and he found in it a great fullness of spirit that
rationalism had not let come into its right. But he maintained already at this time the
independence of philosophy. In the first place he called attention to the fact that
dogmatics has always built on philosophy: if one wished to draw out of dogmatics, as it has
formed itself right from the earliest times, everything that has sprung from philosophizing,
it would collapse. In the second place he maintained that philosophy cannot stand otherwise
toward Christianity than it stands toward all psychological, biological, and physical
phenomena. One must indeed distinguish between the relation of theological assumptions to the
ecclesiastical tradition and the real possibility, as well as the actual character and
significance, of the facts that those assumptions concern. Philosophy stands toward faith and
Christianity as it stands toward any other subject whatever.\footnote{See the essays
\emph{Bidrag til Bestemmelse af Filosofiens Forhold til Teologien} [Contribution to the
Determination of Philosophy's Relation to Theology] and \emph{Hvad er Dogmatik?} [What Is
Dogmatics?] in \emph{Filosofisk Arkiv og Repertorium} (1829--1830).}

Already now he made his reservations toward the ecclesiastical orthodoxy that, in opposition
to rationalism, regarded belief in the Devil and the assumption of an eternal damnation as
necessary constituents of Christianity. His later philosophical research and personal
life-experience led him step by step further. --- In his ``Cosmology'' he conceived God's
becoming man not as an event first taking place at a single definite point of time, but as a
constantly ongoing process, which at a single point reached its highest, its culmination; ---
as humanity developed out of or in animality, so again the Godhead developed out of or in
humanity. The Incarnation takes place ``from eternity'' --- that is to say: now and always!
In this speculative reinterpretation of the ecclesiastical dogma there expresses itself
Sibbern's need to get the closest possible connection between his religious representations
and his philosophical ideas. --- His personal life-experience also led him away from
orthodoxy. It had always, indeed, been the spiritual nourishment he found in Christianity
that had drawn him to it. And now he found too much that was external in the current,
handed-down conception of Christianity. When personal experience and assent are the decisive
thing, then one must let everyone go the ways that suit him. One must let everyone keep that
by which he feels himself furthered in respect of the one thing that is needful, and let him
work freely in his direction, if only he will grant us freedom for the same. Let everyone
keep his fetish! Unity in the religious domain has value only when it comes as a natural
effect of the fact that the knowledge of truth penetrates in all. The ways of presumption,
the external arrangements, the letter-worship, lead only to crushing the independent life. By
many ways spiritual refreshment streams in --- from humanism and naturalism in their way,
just as from orthodoxy, and from science, art, and free love of man, as well as from
religion. When one dogmatizes, one tears the truth loose from its living root and makes a
mummy of it, in which shape it is then handed down from generation to generation. --- Let one
not call this subjectivism, for everything that has value and is later formed as objective
doctrine has had its first origin in the innermost of the soul. Precisely in the most inward
movements in the mind, great forces of existence can stir and strive to work their way
forward. The doctrine of sporadic development finds application here too. It is the way of the
spirit. And besides: what has found recognition in a great circle of men, for example an
existing Church, does not thereby become more objective than what in the individual makes
itself subjectively felt. The objective means the valid, and this does not depend on external
formulations and organizations.

It was in two university programmes from 1846 and 1847\footnote{\emph{Om den christelige
Ytringsfrihed i kirkelig Henseende} [On Christian Freedom of Expression in Ecclesiastical
Respect]. 1846. --- \emph{Bidrag til at opklare den christelige og kirkelige Frihed}
[Contribution to the Elucidation of Christian and Ecclesiastical Freedom]. 1847.} that Sibbern
put forward these thoughts, which on an essential point agree with the proposition that Søren
Kierkegaard at the same time proclaimed: Subjectivity is the truth! --- Sibbern has undoubtedly
come to his standpoint by his own way --- and besides, he draws far more radical consequences
from the proposition than Kierkegaard did.

But it did not become Sibbern's last word concerning Christianity. In close connection with the
views on ethical and social questions that developed in him during his constantly continued
life of thought and by means of his lively sense for the human, he made in his later years a
far sharper distinction between the perishable and the abiding in Christianity. In the
fantasy-image of a future society that he has given in his remarkable work \emph{Meddelelser af
Indholdet af et Skrift fra Aaret 2135} [Communications of the Contents of a Writing from the
Year 2135] (1858--1872), he directs stern reproaches against Christianity. It has in all its
time far more fostered and nourished presumption, lust for domination, and inflamed hatred than
the spirit of love of man. The men of the future will therefore not call themselves Christians;
they will not bear a name together with those who for centuries had committed presumptions and
atrocities in the name of Christianity, and that not only in Catholicism (the Church of sole
rule) but also in Protestantism (the Churches of protestation). They distinguish between
Christian-mindedness [\emph{Kristelighed}] and Christendom-as-establishment [\emph{Kristenskab}].
They read in the great book of nature and life, not only in the Bible, though from this too they
take what can refresh spirit and heart. They live out of a force that stirs in their interior;
their own striving toward the future and the divine activity go into one. Jesus is for them the
man in whom the fullness of light and life first stepped forth fully and wholly in humanity. ---
Through Christianity there has, to be sure, despite Christendom-as-establishment, gone a stream
of love of man. But how has one not often proceeded toward the heathen? The many little ones
among the heathen who had a God-fearing mind, love of man ought just as well to have avoided
offending, as the little ones who believed in Christ. Polytheism had its good ground. Let
everyone remain by his faith, when it is to him a confidence and a furtherance of that which
duty and human care demand! --- If it is asked whether anything can be known about another life,
the men of the future (the New Europeans) answer that whether there is an immortality or not for
individuals, they will live life as they live it; yet it seems that the all-active stirring in
all leads to belief in immortality. ---

There is a deep connection between Sibbern's religious ideas, as they gradually develop, and his
philosophical views. His belief that all development takes place sporadically, from many
different points, and yet leads to a harmonious result, must naturally lead him to lay the
weight on the individual personalities with their need and their experience, and on that which
could work its way forward in their interior. Therefore he speaks so strongly against all
presumption --- and he will not himself be presumptuous either (which his books also --- by
their awkward form --- were far from being able to be). He therefore says at the beginning of
his depiction of the future that no one ought to read it if he notices that ``it will confuse
for him that in which he has his good and secure hold.'' --- Scarcely in any Danish author has
psychological sense and feeling for humanity been connected with so much freedom of spirit as in
old Sibbern. In his strange books there is, in a more positive way than anywhere else, a
reference to a human solution of the religious problem.

\begin{center}---\end{center}

The conception that Sibbern had of philosophy's task, method, and significance is closely
connected with the very character of the views he held on psychological, cosmological, and
religious matters. He has expressed himself on his conception of philosophy in general partly in
his work on \emph{Filosofiens Begreb, Natur og Væsen} [The Concept, Nature, and Essence of
Philosophy] (1843), partly in his work on \emph{Hegels Filosofi betragtet i Forhold til vor Tid}
[Hegel's Philosophy Regarded in Relation to Our Time] (1838).

For Sibbern it is philosophy's highest task to lead to an all-embracing view of the world. But
this is to be built on what is given in experience. And the first task then becomes to undertake
a going-through or an analysis (an ``explication'') of the given. Here philosophy presupposes not
only the other sciences, but itself undertakes such a going-through in the psychological,
logical, aesthetic, ethical, and religious domain. Only when this analysis has been undertaken
can the speculative construction, which will derive everything from a fundamental idea, begin its
work. He blamed Hegel because he wished to go straight to construction and develop all the
content of existence in a single continuous series of thought. And yet Hegel himself, by writing
his \emph{Phänomenologie des Geistes}, had admitted that a foundation and an introduction were
needed before speculation could gather speed. It is the given content of existence that
philosophy will see in the light of the idea --- and it must then first convince itself that this
content is really given. Next it depends on whether the content lets itself be led back to one
fundamental thought. Here, on account of the manifoldness and difference of the content, a long
debate must be carried on before a system can form itself. One thing is the fundamental idea ---
another is how we come to it! From a single abstract thought nothing can be derived, and Hegel's
apparently purely logically advancing line of thought has then also at every point, by
surreptition (subreption), taken up a content given in experience. But experience does not let
itself be ordered in a single series; it presents not only elements that can be placed in
relations of subordination and superordination to one another, but also coordinate elements.

Sibbern's critique of Hegel's philosophy contains many apt remarks and hits several chief points
in this speculative theory, which in the [eighteen-]thirties and forties made itself not a little
felt and found adherents in men such as Heiberg and Martensen. ---

That a speculative view of the world --- on the basis of experience, analysis, and debate --- is
possible, gets according to Sibbern its explanation from the fact that what constitutes existence
in its innermost also constitutes our spiritual existence in its innermost. He thus gives his
theory of knowledge (as regards the highest kind of knowledge) a religious-metaphysical
foundation. In his ``speculative cosmology'' he has himself made an attempt, whose fundamental
features we know from the foregoing. But the speculative theology that he still announced was not
completed, partly no doubt on account of the decisive change in his religious views that occurred
at the close of the forties, partly on account of the fact that the social question from that
time on, as we shall see in what follows, so greatly took his interest captive.

A complete systematization of philosophy was now, according to Sibbern, likewise not possible.
And the reasons he adduces for this are in a high degree characteristic of his whole standpoint
and his view of life.

Experience is scattered and manifold. The world does not open itself to us from the first from
its innermost alone, but from many points at once, and besides it constantly advances further.
This sporadic thing will never be able to be quite gathered into one system.

And the researcher is besides himself one of these points. Philosophy is one of the many
sporadically emerging processes in the world. How, then, should it be able to survey the whole?
--- The disputes between the philosophers are signs that each has taken his starting point from a
different moment in existence, which they emphasize with one-sided force and in a one-sided
direction. Nowhere is it seen more clearly than in the sphere of philosophical debates and
religious movements that our spiritual life still finds itself in the first process of formation.

Even if now, at a given time, a system has gathered the chief moments of the content of existence
and shown their origin from a fundamental idea, yet --- on account of the constant advance of
experience and life --- the work must in every new age be taken up anew. For while philosophy is
systematized, life will have opened more of its fullness, and new knowledge and new points of
view will then become necessary. The debate must begin again. --- Therefore Sibbern was convinced
that things could not rest with Hegel's philosophy, and he defended, against Heiberg, Poul Møller,
who had for a time been a Hegelian, but later felt the need for another view, on which account
Heiberg called him a deserter. Sibbern maintained that what Poul Møller had striven after was a
richer content of experience than the one Hegel had systematized.

Finally, philosophy has for Sibbern only ``intermediary'' significance --- significance as an
intermediate link. The reckoning that philosophical reflection undertakes is to lead to a so much
the deeper appropriation of life's real powers, art and poetry, love and religion, as well as
real nature, the whole realm of experience. Philosophy can be thought's, but not life's, highest
form, already for the reason that thought is constantly only one of life's elements. At the close
of his philosophy the thinker constantly stands over against the other elements of life.

\begin{center}---\end{center}

On Sibbern's special treatment of logic, aesthetics, and ethics I do not here enter, much of
interest and ability though it presents. On the other hand I will, to conclude his
characterization, briefly treat his political and social views.

From the first days of the estates period right up into his advanced old age, Sibbern took part
in the political debate. It was, after all, one of his philosophy's chief propositions that all
development advances through a debate of all against all. His own contributions to the political
strife did not come to play any great role; but they were in a high degree characteristic of
himself by their spirit and tone, and they bear witness to how clear an eye and how warm a heart
he had.

There are in a political respect certain chief thoughts that he constantly held fast. He laid
great weight on a thorough and many-sided consideration of matters, as well as on the possibility
of an energetic government. The ideal was an advisory assembly and a strong government animated by
warm care for the common good, especially for the good of the ``lowly folk.'' Again and again he
looked back to the reforms of the enlightened absolutism, especially to ``the men who with
enlightened and purified will stood for the peasant cause.'' What the Liberals called freedom was
in reality a share in power. When one demanded guarantees, one overlooked that all constitutional
guarantees are in the end moral. Besides --- everything is to be controlled, including the
controlling representation! --- In the [eighteen-]thirties and forties Sibbern was therefore
against a constitution, a ``paper pope.'' (\emph{Patriotiske Intelligensblade}, 1835. ---
\emph{Dikaiosyne}, 1843.)

But if one was to have a constitution, then it should be as free as possible, stand open to as
many as possible. He did not find that in Denmark there were conditions for an opposition between
an aristocratic and a democratic chamber. Therefore he held to the June Constitution and spoke out
both against the Ørsted ministry and against the constitutional change of 1866. (\emph{Om og i
Anledning af Kampen mellem de tvende højeste Statsmyndigheder i Danmark} [On, and on the Occasion
of, the Struggle between the Two Highest State Authorities in Denmark], 1854. ---
\emph{Samfundsbetragtninger} [Reflections on Society], 1865--1870. And a series of pamphlets from
1865--66.)

But already very early the social question was for him more important than the
political-constitutional one. In this respect his university programme from 1849 (\emph{Nogle
Betragtninger over Stat og Kirke} [Some Reflections on State and Church]) is of interest. He here
distinguishes sharply between State and People. The State's power and organization were to abolish
the war of all against all that is supposed to prevail in the state of nature. But has it not
rather made the war stronger, sharper, more all-embracing and incisive, partly between the states
or peoples among themselves, partly within the individual people? It is within the people that the
struggle for gain, even for simple subsistence, gathers speed. The gospel of competition leads to
the most nervous struggle of all against all. Competition has in a high degree promoted
industriousness; but which are the gods for whom altars have thereby been raised, and whose
worship and reign has thereby been furthered? It goes not only hard on those who struggle for the
necessary, but works also nerve-weakening and consuming on the wealthy. All ought therefore to be
agreed on setting another organizing power in place of the blind and raw competition. The highest
justice is distributive justice, which aims at letting all get access to the common possession of
the whole of humanity. Birth ought here not in advance to exclude or favour. What the individual
himself acquires is due in large part to the social conditions under which he works, so that
self-acquisition cannot ground an absolute right of property. It must be the State's task to
organize labour in such a way that all can get a share in the yield and that accidental
head-starts are excluded. In the absolutism's work for the peasantry's subsistence and
development we have an example of such an organizing activity. Only a pity that it did not come to
the benefit of the cottar class! --- Statesmen must in time familiarize themselves with the idea
of communism, which has reality as its foundation and reality as its aim. Even if it has darkness
behind it and darkness before it, it has nevertheless not, like the constitutional system,
illusion or half-heartedness before and behind.

Here too Sibbern applies his doctrine of the sporadic character of development. There are many
scattered efforts and endeavours in the direction of the one right goal; but since connection
between them is lacking, it can look as if the whole were absurd. As men have shown themselves
disposed, indeed bent, on believing the most unreasonable things, so they have also an inclination
to go the most mistaken ways; absurdity seems to have been humanity's cradle and nurse. But if one
takes the foetus at an early stage, it too seems distorted and unnatural, and yet a harmonious
human figure comes out of it. And if one takes a boy in the awkward age, one often sees only
unruliness and refractoriness, restlessness and change, and yet a good and satisfying existence can
come out, when the metamorphosis has been passed through. With these analogies Sibbern consoles
himself: humanity is perhaps still in the awkward age --- if it does not even have its birth first
still to come!

When Sibbern published this programme, he had already begun the working-out of his remarkable
utopia \emph{Meddelelser af et Skrift fra Aar 2135} [Communications of a Writing from the Year
2135]. The first preface is dated 1846.

In this work Sibbern casts a critical retrospect on the nineteenth century in Europe, especially in
Denmark, particularly as regards the political, economic, and religious conditions. He in part
repeats his critique of constitutionalism and parliamentarism; it is besides here that his
strongest statements about ``Christendom'' [\emph{Kristenskab}] and the Church-system appear. We
shall now especially dwell on the social-economic future ideal that he lets his pity and
indignation, awakened by the meaninglessnesses, distortions, and sufferings of human life, dictate
to his fantasy.

At the close of the nineteenth century souls languished. All joy and freshness of life was gone.
One groaned under the government- and parliament-nuisance, under the property-bungling, under the
pressure of laws and taxes, as well as under the presumption of the Church-doctrines. Everywhere
external forms and stiff dogmas had forced themselves in between man and the inner realm of the
spirit. The State was to abolish war, but had only set inner war in place of outer. There were many
people who were not taken up into the common life and did not enjoy its advantages; they had become
outcasts --- but one had oneself cast them out. It was the lack of true human care that had made
the penal and poor systems, the European states' \emph{partie honteuse} [shameful blot], necessary
and possible.

There came now a state of dullness and sleep over men. Everything stood still or was in dissolution.
And at the same time natural causes too worked toward the great rest, in that subterranean vapours
rose up and made all activity impossible. When one then awoke, one was clear that one would not
introduce the old State- and Church-system again, as little as the legal, commercial, and property
system. Labour was organized in such a way that all took part both in bodily labour and in
spiritual striving and living, and that souls were not crushed, but developed and enriched through
the labour. Soul-welfare became the first and most important thing. There were public storehouses,
in which all produced products were kept, and from which everyone could get what he needed. No one
took more than he needed --- for what should he do with it? --- With a multitude of naïve details
the future-image is furnished. Everywhere there expresses itself a remarkable freedom of spirit and
a great sense for the inner life of the soul and its need.

This social fantasy of the future stands in an interesting way in opposition to Romanticism's
absorption in the past and its idealizing of the vanished. Hope had stepped into memory's place,
and that in a man who had himself with warmth lived through the time of Romanticism and been
animated by its ideals. It was, besides the constant inner striving, the fresh sense for the human
that led the old thinker to hold fast to the ideals of his youth in so free and living a way that
they could clothe themselves in quite other forms. He had begun his youth with looking back and
with inwardness immersing himself in what past and present had brought of forms of thought and
life; and precisely because he so faithfully worked in these, it became possible for his old eyes
to look forward with so great confidence.

It is, of course, only a reference to a problem, no solution, that Sibbern gives. His socialism is
utopian, not only because it depicts conditions that quite differ from those we know, but also
because he does not at all hint at any natural and historical development from the present's to the
ideal future's conditions. It is, after all, easy enough to let all sleep themselves away from all
inconveniences and sleep themselves to the new perfections. By his tale of the great slumber he has
indeed frankly admitted that he can point out no transition.

Utopian socialism stands in a peculiar opposition to so-called scientific socialism, which
precisely thinks it can point out the necessary historical development that is to lead from the
present ``capitalist'' period to the future's socialism, but which in return prefers to refrain
from entering more closely upon depicting the future ordering in detail.\footnote{Compare, on the
various principal forms of socialism, my \emph{Etik}\textsuperscript{3} [Ethics, 3rd ed.], XXVI, 7.}
--- Later on I shall have occasion to treat a young Danish author who, nearly at the same time as
Sibbern, set up a similar social future-ideal, with far less sense for personal inwardness and for
the need of the life of feeling, but also with a clear understanding that the way toward such an
ideal did not go around, but precisely through, competition and parliamentarism.

\begin{center}---\end{center}

It is a comprehensive work of thought that a retrospect on Sibbern's long life shows us.
Unfortunately he did not succeed in any single work in giving his ideas so clear and rounded a form
that they could maintain for themselves a constant and easily accessible place in our literature.
Now only he who occupies himself with the history of Danish intellectual life will stop at his
books. But it is then a peculiar freshness, heartiness, and depth that meets the reader, and that
--- despite the strange form --- secures the old thinker an unfading memory.

\begin{center}---\end{center}

% ============================================================
%  12. POUL MØLLER   (Da.: "12. POUL MØLLER")
% ============================================================
\dfchapter{12. Poul Møller}
\label{ch:df-moller}

\noindent It is especially Treschow and Sibbern who have given Danish philosophy in the
nineteenth century its distinctive stamp. The predominant thing in them was psychological
sense and interest, weight on empirical knowledge, emphasis on the significance of
individuality and of the differences of individuality, critical circumspection toward the
speculations. The greatest figure in the history of Danish thought, Søren Kierkegaard, will
show us the same stamp in a new and distinctive shape, and moreover as it is formed by the
sharpest contrast to the most speculative of all speculative systems, the Hegelian, which
had gained great influence here in the country, especially among aestheticians and
theologians.

But before we speak more closely of Hegelianism here in the country and of its passionate
opponent, we must dwell on an intermediate link between Sibbern and Kierkegaard --- an
intermediate link through which Kierkegaard is closely connected with earlier Danish
thought, an intermediate link that has its own independent interest.

Every Danish reader grows warm about the heart at hearing Poul Møller named, the most
Danish of all Danish authors. Here we shall occupy ourselves only with a single side of
him; but we shall there find him again whole and clear, far from --- as one has sometimes
thought --- merely in an accidentally called-forth position. Neither as poet, philologist,
nor philosopher has Poul Møller left behind anything but fragments, but in these he has laid
his whole personality open to the day. The half can indeed sometimes be more than the whole.

Thirty-two years old, Poul Møller went as professor of philosophy to Kristiania (1826), and
after having worked here for five years, he became Sibbern's colleague until his death in
1838. --- On his deathbed he sent, through Sibbern, a greeting to ``young Kierkegaard.'' The
connection between the three men gets its expression herein.

In his philosophical conception Poul Møller was for a time strongly influenced by Hegel's
speculative system. Sibbern maintained later, against Heiberg, that it followed from Poul
Møller's nature that he had to take an interest in a system like Hegel's, but that it no
less followed from his nature that he could not hold fast to such a system. In Poul
Møller's ``Strøtanker'' [Stray Thoughts], the most important contribution to knowledge of
his philosophizing, we find confirmation of this statement.

He saw the significance of speculation in this: that one point of view is consistently held
fast and carried through. In nature (experience) manifoldness and many-sidedness reign. But
a one-sided system, even were it a pure fiction, yields more real formation, when one
studies it, than occupation with a heaping-together of reflections. One learns to think with
organic consistency. He who feels no need to order his system of thought must either have
much genius or little thoughtfulness. It is of great significance to give oneself an account
of one's view of the world in its whole connection. Consistency excludes confusion, which
arises through reminiscences from a foreign view of the world. --- Aptly is here emphasized
the connection of consistent thinking with the unity of personality. It is characteristic
that it was especially this side that interested Poul Møller in his speculative period.

That he did not remain at the speculative and systematic standpoint had first and foremost
its ground in the recognition that the problems did not find their solution by the way of
speculation, but really only showed themselves to lie at other points than one had assumed:
``By speculative philosophy one obtains not so much an answer to the questions that ordinary
human understanding presents, as one comes to recognize that the questions themselves, as
grounded on concepts that have no truth, must fall away.'' --- But next he missed in
speculative philosophy a sense for the significance of individuality. The various
individualities are discovered only by the way of experience, and why should that be an
imperfect cognition? I do in fact thereby widen my knowledge, though it does not happen by
an a priori way. And if only a priori cognition is the true one, individuality cannot be
cognized at all. --- Poul Møller here stands by thoughts that especially Treschow had
impressed, and that had played an essential role for A. S. Ørsted and Sibbern.

In his essay \emph{Om Muligheden af Beviser for Menneskets Udødelighed} [On the Possibility
of Proofs for Man's Immortality] (1837) Poul Møller develops these thoughts somewhat more
fully.

A completed knowledge is unattainable. The empirical sciences never become finished with
comprehending the existing individuals and their modes of existence in a complete system of
concepts. And even if such a conclusion were possible for a certain limited domain of
experience (the terrestrial globe, the human race), that comprehending can nevertheless take
place only by means of universal concepts, and the full intuition is lost. One must hold to
species and genera --- but what an infinite content does not then vanish for us with the
individuals!

If one leaves the way of experience and holds to the most abstract points of view, one can
thereby perhaps bring the general fundamental relations of existence to clear consciousness,
but the whole infinity of determinations (qualities) cannot thereby be exhausted.

A third, higher mode of cognition, which at once lets the individual and the universal come
into their right, does not stand at our disposal.

Nevertheless Poul Møller thinks that a view of the world can be developed in which not only
the reality and significance of individuality is maintained, but also personal immortality
can be demonstrated. He admits that such a view of the world is still only in its becoming;
but he does not see that he himself has declared it impossible, since that ``third, higher''
cognition is not attainable. He appeals to a science whose impossibility he himself has
proved, and which he moreover treats in very vague terms. --- The matter was that there
stirred in him --- especially under the impression of a great sorrow --- a personal need for
faith, which made him now place himself in a more sympathetic relation to positive religion
than before. ``The Christian tradition'' is often invoked in the said essay, and it is
demanded of the philosophical view of the world that it shall let this tradition come into
its right. How far this is possible, he did not come to investigate thoroughly. ---

The break with speculative philosophy has scarcely been without pain and vexation, because
such great expectations had attached themselves to it. But his sound humour again gained the
upper hand. ``For a time,'' says Kierkegaard, ``he spoke almost with indignation of Hegel,
until the sound humorous sense that was in him taught him to smile, especially at
Hegelianism, or, to recall Poul Møller still more clearly, to laugh right heartily at it. For
who has been in love with Poul Møller and forgotten his humour; who has admired him and
forgotten his health; who has known him and forgotten his laughter, which did one good, even
when it did not become quite clear to one what it was he laughed at; for his distraction
sometimes brought one into perplexity.'' --- Yet he held fast to the admiration for Hegel's
great work of thought. But it was now so, he thought, that every giant in philosophy was to
fall by dwarfs. And he had no great respect for those who sought ``to go beyond Hegel.'' One
of them he compares to a rat that digs its way out through the wall of a cathedral. More than
a tent for use on a journey could for the present not be had in place of the proud building.
Such a tent he himself was working on just before his death (a doctrine of the categories or
fundamental concepts); as far as one can see, it would scarcely have been suited for a long
journey.

\begin{center}---\end{center}

Although Poul Møller referred to tradition as that which should also come into its right for
thought, it was nevertheless his definite conviction that only that spiritual content has
value which is not merely appropriated, but produced anew in each individual. In his
aesthetic conception, as we come to know it from his literary reviews, he distinguishes
between the original art that springs from life and nature, and the derived art that is
predominantly, if not exclusively, inspired by the preceding art, not by one's own immediate
relation to reality. Poetry is a fruit of the people's life. ``What a people and its
individuals have experienced yields, in a certain sense, the material for its higher life in
the art of poetry. But when a rich poetic literature has developed, art gradually acquires
more and more independence, in that the chief source of new poetic works becomes ever more
the existing poetic literature, and the original poetry, which with freshness and originality
proceeded from real life, becomes an ever rarer phenomenon\ldots\ When in a disproportionate
number of a people's individuals there arises an inclination to lead a poetic spectator-life,
then a poetic period must soon be over, and the ore is then surely no longer fetched from an
abyss that is accessible only to genius.'' (\emph{Recension af Extremerne} [Review of
\emph{Extremerne}].)

Like Romanticism's philosophy, its poetry too had thus, for Poul Møller, come to a critical
point. A definite programme he could not give for the future --- but as regards both domains
he was sure where the ore was to be fetched. In a living personal experience, in close
connection with life, thought and poetry were to fetch the new force. And in some of his
poems, as well as in manifold of his thoughts, he has himself given models --- unfortunately
only little carried out. As regards his philosophizing, there is especially one circle of
thoughts that occupied him in his later years, and that, fragmentary as it is, nevertheless
marks his most important contribution to Danish thought. On it we must dwell somewhat more
fully.

\begin{center}---\end{center}

When Poul Møller demanded a closer connection between thought and experience than Romanticism
had testified to, it sprang for a large part from his lively sense for the original, genuine,
and fresh. He demanded above all that thought should be connected with man's own innermost
being, have its root in this and express this. It is a great loss for Danish literature that
the essay \emph{Om Affektation} [On Affectation] that he had contemplated was not written.
Only some stray thoughts and a sketch on a few sheets are extant. Poul Møller had caught
sight of what Holberg already so strongly emphasized, the frequent confusion of the shadow
with the body.

It was, however, not only from an epistemological and psychological interest that he occupied
himself with this subject. Nor was it only from an ethical interest in honesty and
uprightness in the life of thought, in personal truth.\footnote{The expression ``personal
truth'' occurs, as far as I know, for the first time in Poul Møller. Synonymously with it he
uses the expression ``the truth of the personal life.''} It was at the same time, as Vilhelm
Andersen has shown, personal disappointments and bitter experiences that caused weariness of
life and called forth a skeptical and defiant mood, which he seems to have thought of giving
expression in a poetic presentation, for whose chief figure he had chosen Ahasuerus. There
seems to be an inner connection between the preliminary works for ``Affectation'' and for
``Ahasuerus.''\footnote{Vilhelm Andersen: \emph{Poul Møller}. 1894, p.~244. --- I cannot
follow Vilh.\ Andersen in the influence he thinks the Ahasuerus reflections had on
Kierkegaard. On the other hand it is to me (as I have already stated in my book \emph{Søren
Kierkegaard som Filosof} [Søren Kierkegaard as Philosopher]. 1892, pp.~24--27) beyond doubt
that Poul Møller's discussion of affectation had a great significance for Kierkegaard,
although he gives Lessing all the credit. And behind Poul Møller stands again Sibbern. ---
What Poul Møller --- see below --- calls ``alternating affectation'' leads the thought to
Kierkegaard's ``rotation-drive'' [\emph{Vexeldrift}] (in the first part of \emph{Enten--Eller}
[Either--Or]).} Poul Møller belonged to those who preferred to think with the whole person.
It is then certain that the line of thought in him that is now to be brought forward was of
decisive significance for him.

Affectation is at once falseness and self-deception. On the other hand, falseness for itself
and self-deception for itself are not affectation. The self-deception can be quite unmerited
and is then not affectation. And he who dissembles with full consciousness is likewise not
affected. For affectation there is required a striving to be what one cannot be; but such a
striving is impossible when one does not for a while imagine that one is what one is not.

``The person who has affectation in his life determines himself in so far not with complete
moral freedom; his actions have not their origin from the true self, which is his free,
ethical will. His will is determined by some merely natural aim, whereby he is brought to
represent a foreign person or to take on a false role that has not been assigned to him in
life.''

I understand this thus: There is, besides the shifting representations and interests in every
person (and in each single period of life), a central content, certain deep-lying thoughts,
aims, and interests --- a will that need not express itself immediately in every moment, but
constantly lies at the foundation. It is what I in my \emph{Psykologi} (V\,B, 5) call the
real self. Only the spiritual expressions of life that accord with it, and especially those
that spring from it, are in truth free and genuine; only in them is there personal truth.
But often the central thing conceals itself; it is overwhelmed by external influences, or one
shuns the exertion that can be connected with bringing it to bear. --- It is a less felicitous
expression that Poul Møller uses when he speaks of ethical will and moral freedom. For it is
evidently a matter of a relation between the inner and the outer, which is the same whether
the inner (the real self) is ethical or not. What matters is that man's inner really lays
itself open to the day in thought, word, and deed. In that consists personal truth. It
becomes a question for itself whether the inner that reveals itself has moral value or not. ---

Poul Møller sharpens his conception in a proposition that he knows will be regarded as
fantastic: ``No expression of life has truth unless there lies in it creative self-activity.
The utterance, the presentation of a thought really has truth only when the thought is in the
same moment produced. Afterward there is already mimicry in it. --- But nonetheless it is
human to reproduce thoughts; for only a Godhead is always creating. Man must in the realm of
thoughts be reckoned among the ruminant animals.''

Although he makes a reservation, he nevertheless holds fast, as the ideal, to the highest
self-activity in thought and word. And he expects much from an approximation to this ideal:
``If every single man, without fearing censure for simplicity, judged of things as they
presented themselves to him, then glorious phenomena would come forth.'' But out of
comfort-seeking or cowardice one does not let one's personality come forth; one does not
believe in one's person's infinite depth. One is ashamed to stand naked and therefore dares
not leave the traditional representations. One does not see that a thought that is uttered
with heartfelt conviction cannot be quite incorrect --- that it contains a point of view, even
if this point of view is wrongly regarded as the only possible one. ---

One sees that the Ahasuerus-mood in Poul Møller concerns only the frequent disproportion
between the inner and the outer. He is convinced that there are hidden treasures in men's
interior; it is only his sorrow that they do not come forth.

In the light of the demand for creative self-activity, the reference to the Christian
tradition in the essay on immortality appears as a contradiction, or at least as a problem:
how can a tradition contain the truth, when the truth essentially depends on man's own
self-activity? This problem Poul Møller did not take up. For Kierkegaard it became precisely
the chief problem. He wished to show that precisely the most inward self-immersion, the
strictest carrying-through of the proposition that subjectivity is the truth, would lead over
to the ideal that is contained in ``that handed down from the fathers.'' ---

There can, according to Poul Møller, be different degrees of affectation. --- In the
\emph{momentary} affectation it is inadvertently that one oversteps the boundaries of one's
personality. One has for a while intoxicated oneself in something one has not yet quite
appropriated. Especially in youth one easily makes oneself guilty of small, half-voluntary
dishonesties of this kind. It can even be a good sign, since it bears witness that one does
not egoistically shut oneself up in oneself. And a deeper self-understanding can be the
consequence, when the personal untruth is discovered. --- The \emph{fixed} affectation sets
in when the assumed role is held fast, becomes habit, and is pursued with consistency. An
involuntary resistance is made against being torn out of the self-deception; the personality
is falsified; one makes oneself into a changeling. An honest lie is not so bad as this
self-falsification; for the fully conscious lie one can more readily lay aside than the false
role one has accustomed oneself to play. --- The \emph{alternating} affectation is present
where every inner kernel is lacking. One forms for oneself in every moment a new personality
so as to lay it aside in the next moment, either out of whim or out of regard for one's
surroundings. In its completion this kind of affectation is one with a moral suicide. ---

As already hinted, it is not seen that Poul Møller has entered upon the problem of how one ---
to use his own expression --- can pass from being a ruminant to being self-creating. Man must
take in spiritual nourishment as well as bodily, and in both respects the nourishment must
correspond to the organism, be able to be made into flesh and blood in it. There becomes then
a question how far all nourishment that is offered can do service in this respect. A little
hint that shows in what direction Poul Møller might perhaps have applied his idea of personal
truth in the religious domain is given by the following stray thought: ``The modern idols,
which are made of words, stand in this respect behind the antique ones of wood and stone,
that they leave thought less freedom than these. Every antique artist nevertheless gave a
Zeus after his own ideal. Now every theologian is to recite certain formulas about the
highest being.'' --- A philosophical application is contained in the short stray thought:
``Affectation is symmetry in philosophical systems.'' The meaning is evidently that one
sacrifices the original and immediate nature of observations and thoughts in order to get them
to fit into the systematic form; thus, according to Hegel, thought was constantly to move
forward in three-membered relations, and the content was adapted to a form that did not form
itself involuntarily. --- In the aesthetic domain, the derived poetry was a sign of
affectation: ``Modern poetry seeks life in poetry at single places and picks them out. Hence
its affectation.'' ---

It is, despite the fragmentary form in which Poul Møller's thoughts are accessible to us,
nevertheless a characteristic type of Danish philosophizing that appears in them. He points
back to Holberg and forward to Kierkegaard.

\begin{center}---\end{center}

% ============================================================
%  13. HEIBERG AND MARTENSEN   (Da.: "13. HEIBERG OG MARTENSEN")
% ============================================================
\dfchapter{13. Heiberg and Martensen}
\label{ch:df-heiberg-martensen}

\noindent In the foregoing Hegel's philosophy has already been treated several times, and
it is now the place to dwell on the influence it won in Denmark.

Hegel sought to connect Kant's sober striving for foundation with Schelling's and Steffens'
romantic immersion in the interior of nature and spirit. He wished to apply a definite
method and put an end to the Romantics' mystical wallowing and fantasizing. It became his
fundamental thought that when the human spirit has wandered through the long way of
experience and history --- a wandering that he himself, in a very romantic way, described in
\emph{Phänomenologie des Geistes} --- there opens for it a circle of thoughts in which
everything from the highest to the meanest can and must find its place. And within this
circle of thoughts there is so intimate a connection that by beginning with a single
definite thought --- were it even with the most abstract of all, the concept of being in
general --- one comes with logical necessity over to all the other thoughts. Every single
thought releases its opposite, and then it becomes the task (a task that, according to
Hegel, can always be solved) to find a thought that, as a higher unity, comprises both
members of the opposition in itself. Every thought gets through this dialectic its
distinctive place in the whole world of thought. It must give up its absolute independence
and its external opposition to other thoughts, and in so far it perishes; but its
distinctiveness is secured to it as a moment within the higher unity.

But this whole movement and struggle of thought is now, according to Hegel, not merely
something that takes place in the human spirit. It is the innermost essence of existence
that thereby reveals itself. By an analogical inference of which he is not conscious, Hegel
conceives everything that happens in the world of nature and history as an expression of a
cosmic dialectic, through which the individual beings and events show themselves as members
in a rising series of higher unities. What in experience may look like destruction is
transition to being a member in a higher connection. And through this inner conflict, which
only for sensuous conception appears as a process in space and time, there expresses itself
the life of the world-spirit, of the Godhead. The Absolute lays itself open to the day in
the oppositions of the relative forms as well as in the connection that the oppositions and
their constant abolition demonstrate.

Hegel shares Romanticism's conviction of the inner unity of spiritual life. Religion, art,
and philosophy are for him expressions of one and the same thing, namely of that inner unity
which makes itself known in the connection and lies at the foundation in the conflict
itself. But of these forms of spirit, in which the Absolute becomes manifest, philosophy is
the highest, because it thinks and comprehends where religion and art only represent, feel,
divine, or intuit.

\begin{center}---\end{center}

It was Johan Ludvig Heiberg (1791--1860) who first introduced Hegel's philosophy into
Denmark. He had studied it with Hegel himself in Berlin and had entered into a personal
relation with the master. On the return journey from Berlin there came a moment when the
fundamental thoughts of the system stood before him all at once with full clarity and in
their whole connection. And Heiberg held fast his whole life to these fundamental thoughts
and to the conviction that only by them could there come unity into spiritual life, and that
especially art and religion could come to stand as fully justified forms of spirit only by
being seen in the light of them.

It was in the dispute about the freedom of the will that Howitz had called forth that
Heiberg first appeared as a philosophical author (\emph{Om den menneskelige Frihed} [On
Human Freedom]. Kiel 1824). Later it was especially aesthetic questions that gave him
occasion to apply the new philosophy, in particular to assign the various kinds of art their
mutual place (\emph{Om Vaudevillen} [On the Vaudeville]. 1826. --- \emph{Om Malerkunsten} [On
the Art of Painting]. 1838). Yet he also treated speculative theology (\emph{Treenighedslæren}
[The Doctrine of the Trinity]. 1837). And he has besides, for use in his lectures on logic at
the military high school, had printed a presentation of logic after Hegel's method. --- Some
of the said essays appeared in a periodical he began (\emph{Perseus. Journal for den
spekulative Ide} [Perseus: Journal for the Speculative Idea]), but of which only two numbers
appeared.

The place that Heiberg, following Hegel, assigned to philosophy within spiritual life is seen
especially from his work \emph{Om Filosofiens Betydning for den nærværende Tid} [On the
Significance of Philosophy for the Present Time] (1833).

What the age suffers under is, in the first place --- Heiberg thinks --- an overpopulation of
ideas. Alongside art, poetry, and religion, scientific and political interests and points of
view make themselves more and more felt. But in what order, in what mutual relation, are all
these ideas and interests to be placed toward one another? Here only a strict comparison of
the various fundamental thoughts among themselves can bring the decision. Already in a letter
to H. C. Ørsted from the year 1825 Heiberg emphasized it as Hegel's great advantage over
other philosophers that everything is found in his system and finds its definite place and
the validity that belongs to it. --- But the age does not always present a disordered
manifoldness of ideas. There often appears a one-sidedness, in that a single idea or interest
--- for example the political or the religious --- sets itself on the high seat as the one
decisive and determining one. Here too the remedy is only to be sought in an investigation of
the mutual relation of the individual ideas or interests. No idea or interest that has real
significance may be pushed aside for others, and it is in the end the truth, life as a whole,
that is to assign each its place. --- By both ways we are thus led to philosophy as the
carried-through doctrine of the fundamental thoughts (the categories) in their mutual
relation. Philosophy presents in connection the idea of truth, or the world of ideas, to
which all the individual interests and endeavours must in the end appeal their cause.

But philosophy nevertheless brings nothing absolutely new. It is constantly only a reflection
on, or a cognition of, what is already given. It is in so far contained in sound human
understanding. It is the cognition of the truth that is contained in the various human
experiences and endeavours --- in nature and the state as in art and religion. Philosophy
unfolds itself only when it is a matter of reflecting on a period of development that is
concluded; by transposing the developed content into the form of thought, philosophy makes it
eternally present.

And by opening our eyes to what we already possess, philosophy also brings us help for a
sickness in spiritual life that is no less serious than the chaotic or one-sided character of
the thoughts. The age suffers from an opposition between ideal and reality, between the
Infinite and the Finite. The finite and conflicting endeavours have pushed the Infinite and
Eternal out into a distant background. What we are to learn, and what philosophy can teach us,
is to see how the Eternal and Infinite lies precisely at the foundation of the individual and
finite endeavours, so that what we seek is given us already in our very seeking, not merely as
a distant and beyond goal. Hegel's philosophy reconciles, like Goethe's poetry, the ideal with
reality, our demands with what we possess, our wishes with what has been attained. ---

That was Heiberg's programme. We now understand the interest he had in the abstract conceptual
developments in the Hegelian logic:\footnote{Heiberg has made a change in the logic that is
undoubtedly a correction of Hegel's presentation. He calls attention to the fact that becoming
cannot, as in Hegel, occupy the place of a ``higher unity,'' since becoming means unrest, an
external mixture of being and not-being. In place of Hegel's first triad ``Being--Nothing--Becoming,''
Heiberg then sets the triad ``Being (which, as purely abstract and contentless, is
Non-Being)--Becoming--Existence [\emph{Tilværen}],'' where Existence is the result of the
becoming into which Being passes over when it is seen that it contains the contradiction of
being ``Non-Being.'' \emph{Perseus} II, p.~44. --- Of course this correction too rests on
Hegel's own presuppositions.} it was in this abstract world of thought, in which each single
domain of life and experience was represented by the fundamental concepts that expressed its
essence --- it was there that the struggle between the various endeavours was in the end fought
out. In Hegel's system logic was now only the one third; the second part was the philosophy of
nature, in which the abstract ideas appear in the forms of space and time; the third part was
the philosophy of spirit, in which that which is in time and space externally separated is
gathered in the unity-forms of spirit, in the end in thought itself, whereby the system in its
conclusion returns to its beginning and thus forms a closed chain.

Heiberg, who everywhere seeks to find triads and, where he finds triads, seeks to give them
purely abstract expression, compares the system's three parts with Hell, Purgatory, and
Paradise in Dante's poem. He himself has then most liked to dwell in Hell (logic) and in
Paradise (the philosophy of spirit); with Purgatory (the philosophy of nature) he has not
occupied himself. It is perhaps a sign of his criticism; for the philosophy of nature was, as
Sibbern rightly remarked in his in many respects interesting work on Hegel's philosophy, the
system's \emph{partie honteuse}. On the other hand he has, as already remarked, worked in the
domain of aesthetics and philosophy of religion.

When Sibbern in the said work called Heiberg a philosophical dilettante,\footnote{In return,
Sibbern's books were given, in \emph{En Sjæl efter Døden}, a place of honour in the library in
Hell (which here is the home of eternal boredom).} it was in so far justified as Heiberg is to
such a degree overwhelmed by the Hegelian dialectic that he often applies it in a very external
way, without regard to the distinctiveness of the subjects. In Hegel himself, however, one
finds a deeper sense for the realities that he thought he had given thought-expression to in
his concepts. Between the master and the disciple there shows itself here a peculiar difference
both in the domain of aesthetics and of philosophy of religion.

Heiberg attempted a systematization of aesthetics, even before Hegel's aesthetics had
appeared. His interest here aims at a division of the kinds of art --- in analogy with the
abstract logical categories, and in such a way that the triad --- an element, its opposite, and
the higher unity of both --- constantly comes forth anew. Beauty appears first as beauty of
form and beauty of nature; then as beauty of art; and finally as poetry, which forms a
transition to philosophy. Within poetry there again (in \emph{the essay On the Vaudeville})
enter lyric and epic in opposition to each other; the higher unity of both is the drama, within
which there is also a mutual opposition between lyric and epic drama (the former depicting
character, the latter situation) --- without the higher unity here being indicated. The whole
development serves, namely, to give the vaudeville (as a kind of situation-drama) its definite
place, and the investigation therefore ends here. Later (in \emph{the introductory lecture to
the logical course}) one ascends from tragedy through comedy up to the lyric drama represented
by Calderón as poetry's highest form.

Against Oehlenschläger's objection that such divisions into fixed kinds are to the detriment of
poetry's free development, Heiberg answers that one honours the great poets in the best way by
regarding the various forms they have used as something fixed and definite, not as something
arbitrary. One demonstrates the deep reason that reveals itself in the poetic genius, and
thereby secures poetry its high rank in life. But it was nevertheless a strange formal
classification that Heiberg here fell into. Hegel went another way, in that he conceived art in
its connection with historical development and thereby got the great triad of oriental,
classical, and romantic art. He moved in his dialectic predominantly forward through historical
characterization, and the abstract categories step into the background. His aesthetics Hegel had
fetched from the same source as his system on the whole; but Heiberg had first familiarized
himself with the system and then sought afterward, in a schematic way, to apply it to the
special domains.

Something similar holds of his development in the philosophy of religion. As for Hegel, so for
Heiberg too, religion and philosophy have the same content, only the form is different: the
Absolute (God) appears in the one in the form of representation (imagination), in the other in
the form of thought. According to \emph{En Sjæl efter Døden} [A Soul after Death] one can come
into the kingdom of heaven in two ways: either by immersing oneself in Christ's historical
figure, or by having speculatively ``grasped God.'' Only this latter form is the perfect form,
and to it one must therefore again and again go back, if religion is to have validity. For only
that is valid which can take on the form of thought. Philosophy does not prove God's existence,
which cannot be proved at all when it is not presupposed from the first. But philosophy makes
the representation of God, which is presupposed as given, into a concept of God, thinks God as
the Absolute, which has no opposition over against it, because it is itself the unity that makes
all oppositions and all higher unity of oppositions possible. In the Christian doctrine of the
Trinity Heiberg finds an agreement with Hegel's philosophy: ``The constant trilogy running
through the Hegelian system is the very reflection of the Trinity in the realms of thought,
nature, and spirit, and the absolute condition for all philosophy.''

If God is to be thought as infinite, one must by infinity not think of an unlimited extension in
space and time, which is only a constant transition from one finite thing to another, a sum of
finitudes; that is what Hegel calls the bad infinity.\footnote{Heiberg thinks that the Copernican
system has worked to entrench the notion of the bad infinity. ``One has become unfair toward our
globe, as if one came nearer to God by living in a larger and more magnificent house, as if our
knowledge of God were not a knowledge of the Highest and thus could not be surpassed by any other
knowledge, at least in quality.'' (\emph{Perseus} I, p.~27.) --- Compare Hegel's
\emph{Encyclopädie}, §~280, where the planet is the higher unity of the sun and the moon (or the
comet). --- What takes place on it must accordingly be the most important thing! But there are so
many planets!} The true Infinite is that which in itself has a fullness whose elements stand in an
inner, necessary relation to one another. Thus there is in self-consciousness an inner relation
between subject and object; they are at no point outside of or indifferent to each other. There is
here a circle running back into itself, in which nothing is fixed or unchangeable, first or last,
cause or effect, means or end. But the true Infinite cannot stand in an external, and thus finite,
relation to the finitude that spreads out in time and space. If the Infinite is not to have its
barrier in the Finite, and thus be limited by it --- whereby it would precisely cease to be
infinite --- then it must itself reveal itself in finitude, so that this is a necessary form
through which it reaches full definiteness. And that is precisely what Christianity teaches:
through the becoming-man, God himself steps into finitude; the second person in the Trinity
precisely designates the difference and the finitude, in so far as this has a place in God; and in
the third person, the Spirit, the perfect harmony of the Infinite and the Finite is expressed. ---
By this conception Heiberg thinks he is beyond the dispute between the Rationalists and the
Orthodox. While the former set the Infinite outside the Finite, the latter set the Finite outside
the Infinite. Both are self-contradictory.

In Hegel too there occur logical abstractions as interpretations of religious assumptions. But one
nevertheless often notices in Hegel a personal life-experience or a historical immersion in the
struggle and pain that ``finitude'' causes, before it can be ``abolished.'' Heiberg passes too
quickly from experience to pure logic. Not for nothing did he set Calderón far above Shakespeare
as poetic representative of humanity, because the former dissolves everything into the pure form
of infinity, while the latter gets stuck in characterizations and events that awaken wonder and
terror, but do not point out beyond the finite! --- And for Heiberg himself the purely logical
interpretation was decisive for the possibility of maintaining a religious standpoint. He writes
to H. C. Ørsted (1825) that he would be an atheist if his philosophy did not teach him the spirit
revealed in nature and history. It has surely been his earnest with this confession; he has been
convinced of the reconciliation of faith and knowledge that his philosophy taught. A faith that
could not thus be transposed into conceptual knowledge would for him be impossible. He belonged
most nearly to what in the [eighteen-]forties was called the Hegelian Right, not to the Left,
which with Strauss and Feuerbach drew radical consequences from Hegel's philosophy of religion.

The especially abstract character that Hegel's philosophy got in Heiberg's hands finds its
explanation not only in the fact that he had not, like the master, laboriously fetched the ideas
up from reality's pits, but also in the fact that in his poetry he found the best expression for
the more personal and living side of his spirit. In the vaudevilles it was finitude that ---
transfigured in jest and lyrical mood --- maintains its place as a moment in the Infinite; but in
his thought-poems it is an intimate union of ideas and feelings that get an expression unique in
our literature. Many abstract volatilizations can be forgiven him whose view of life has led him
to write the beautiful lines:

\begin{verse}
The God you seek is your own God;\\
what would you more?\\
He steps not forth as an object\\
among others, several.\\
You feel him, every time your heart beats\\
with holy desire;\\
you hear him in your best thoughts\\
as an eternal voice.\\
In seeking, you have already found him;\\
in asking, the answer is already won.
\end{verse}

When Heiberg in the year 1841 published his ``Prosaiske Skrifter'' [Prose Writings], he recalled
in the preface what resistance the Hegelian philosophy had met, in that the appointed teachers
had first neglected it and then sought to keep it under strict quarantine [a jab at Sibbern and
Poul Møller]; but he rejoiced that ``what had been decried as a matter of fashion and an empty
formula-system has now, through a younger worker's powerful, brilliant activity, become \emph{causa
victrix}.'' --- This younger worker was Hans Lassen Martensen (1808--1884), who just then, with
great attendance, delivered his lectures on speculative philosophy and theology at our university.

In his student years Martensen had been strongly taken by Hegel. He was especially enthusiastic
over the prospect that Hegel's philosophy gave of a speculative theology that could stand as a
higher unity over against rationalism and orthodoxy. To the Rationalists he reckoned Schleiermacher
too, partly because the latter sought a foundation for his theology lying in human nature (the
feeling of dependence), partly because he attached himself to Kant's impressing of the limit of
scientific cognition. When Schleiermacher visited Copenhagen in the year 1833, the young
theological candidate sought him out and discussed the relation between dogmatics and philosophy
with him. ``Once,'' Martensen tells in his \emph{Life}, ``I asked quite naïvely whether he assumed
that a philosophical cognition was possible of God's essence in itself, of the inner, eternal
life-process in God\ldots\ He answered quite calmly: Ich halte es für eine Täuschung [I hold it to
be an illusion].'' This did not, however, warn the young, enthusiastic, speculative theologian.
Even if one granted Schleiermacher, he thought, that we can only think in oppositions, this does
not exclude a scientific concept of God, since God would not be the living God if there were not
\emph{inner} oppositions in his essence --- which is precisely what the dogma of the Trinity
maintains. Here he attached himself to Hegel. But that he broke with Hegel came from the fact that
the latter, like Kant and Schleiermacher, wished to go the way of purely human cognition. In his
dissertation ``On the Autonomy of Human Self-Consciousness in the Dogmatic Theology of Our
Time,''\footnote{\emph{De autonomia conscientiæ sui humanæ in Theologiam nostri temporis
introducta}. Hauniæ 1837. --- A Danish translation appeared in 1841.} from which one has dated
``a new era'' in theology, he maintains that there is no purely independent theory of knowledge,
but that the foundation of all cognition is the consciousness of man as a being created by God.
This Kant and his successors (Schleiermacher and Hegel) had overlooked. Only when consciousness,
by virtue of an original co-knowledge (or conscience), proceeds from the first from the source of
all reality, can its knowledge become more than purely formal, and then a cognition of the
Infinite and Absolute becomes possible, although man himself is a finite being.

The independent (autonomous) human knowledge never gets beyond itself, never grasps the actual
reality, but halts at an insuperable opposition between knowledge and power. Only through the
relation to God, in whom knowledge and power are eternally connected, do we get beyond this
opposition. God is autonomous, man cannot be. To be sure, it depends on man's will whether he will
acknowledge that original knowledge of God dwelling in the human spirit; but just as sure is it
that if one does not begin with God, one begins by denying him. Faith is the presupposition of
cognition. \emph{Credo ut intelligam} (I believe in order to be able to understand), as Anselm
said.

Martensen admits that theology must presuppose a theory of knowledge, and calls attention to the
fact that every age's dogmatics is dependent on that age's conception of human cognition. But he
gets easily over this admission again by the simple trick of making all theory of knowledge into
theology. Philosophical theory of knowledge, on the other hand, builds on an analysis of the
actually given science and is an attempt to present its presuppositions in their connection with
the nature of the human spirit. Martensen has neither in his initial work (in several respects
the most thought-through he has produced) nor later entered upon showing how one can, from his
dogmatic presuppositions, come to the actually given science's methods and results, not even how
his presuppositions can subsist together with these. And yet he himself feels a sting here, in
that he declares that the difficulty is not to understand that God is, but to understand that
there can be anything \emph{besides} God, since it seems to follow from God's absolute nature that
there can be nothing besides him. This difficulty is, on the other hand, already fallen away by
the creation-dogma on which the theological theory of knowledge builds. --- Of course all riddles
fall away when one from the first proceeds from the assumption that they are solved.

Nevertheless Martensen lays great weight on the dogma's being subjected to a closer analysis. It
is precisely theology's task to give a speculative unfolding of what \emph{lies} in the
dogma.\footnote{\emph{Munus theologiæ dogmaticæ in iis quæ dogmate implicite continentur
speculative explicandis versatur} [The task of dogmatic theology consists in speculatively
explicating what is implicitly contained in the dogma]. \emph{De Autonomia}, p.~29 f.} If one does
not develop the speculative ideas contained in religion, then one gets on the one hand a science
that stands outside religion, and on the other hand a purely historical theology, indeed only a
bibliology as in the current faith (\emph{supranaturalismus vulgaris}). --- Martensen here
announces a conception of theology according to which it is the common mother and centre of all
sciences. It was also, as he tells in his \emph{Life}, his ideal of youth, called forth and
nourished by Steffens' theological writings: ``There must be able to be a view of the world and of
life in which everything that has significance in existence --- nature and spirit, nature and
history, poetry, art, philosophy --- fits harmoniously together into a temple of the spirit, in
which Christianity is the all-dominating and all-explaining centre.''

To raise this temple, on which he worked in theory and practice his whole life through, unshaken
by even the most violent objection, he wished, however, also to make use of philosophical aids.
That ``explication'' was to take place through the dialectical method found by Hegel. Therefore he
declares that theology and philosophy have the same method. He was convinced that when thought
immerses itself in the dogma, the latter's content will unfold itself according to its nature, let
its various moments appear in inner connection. The series of the special dogmas he then presented
in this way twelve years later in his Dogmatics (1849) --- and undeniably it was not noticed here
that there were great difficulties to overcome. (Only quite at the end, at the dogma of eternal
damnation, does he find --- what serves more to the honour of his humanity than of his logic --- a
``cross for thought.'') It is, however, not, as in Hegel, purely logical transitions (through the
demonstration and solution of contradictions), but it is the intuition, supported by historical
knowledge of dogma and guided by religious imagination, that accomplishes the work. It is no
accident when Martensen in his later days prefers to say that his Dogmatics was a reproduction of
the Church's teaching, rather than that it was an explication.

Martensen's relation to Heiberg is here quite peculiar. Heiberg did not rightly understand why
Martensen was not satisfied with what Hegel gave him. Hegel presupposes, Heiberg maintains, the
immediate constantly; it is precisely that which he wishes to transpose into the form of the
concept; but when one begins to philosophize, the immediate is the object that is philosophized
about, and it depends then only on whether one can, in one's philosophizing, exhaust it. A child
has parents; but it is to live its own life later on, as if it stood alone in the world, and not
spend its life thinking of its parents. One may well say with Martensen that the fear of God is
the beginning of wisdom; but when our wisdom is formed into a system, we must begin with doubt and
see how far it can be solved; the fear of God one is to have in one's heart, doubt in one's eye.
Since Martensen, however, wishes to keep Hegel's method, Heiberg will not reckon him (as Poul
Møller) among the deserters.\footnote{\emph{Perseus}, I, pp.~33--41. --- Heiberg's statements
concern an essay by Martensen in \emph{Maanedsskrift for Litteratur} [Monthly Journal for
Literature] that precedes his dissertation.}

What Heiberg did not see was that ``the immediate'' in Martensen does not merely mean object; it
also means authority. He begins with a dogma: the creation-dogma, and remains within it. Martensen
has with full right maintained that he was not a Hegelian when he appeared as an author. He had, on
the other hand, attached himself to the religious thinker Franz Baader, in whom then also the chief
propositions of his dissertation are already found. For Baader, autonomy (the independence of
thought and will) leads to anomy (lawlessness) --- and in the end to spiritual death. For, he says,
``every spirit lives only in and by admiration and adoration; in subordinating itself to what it
admires and adores, it goes out to glorify, spread, and present it. By turning away from the true
and only object of its admiration, adoration, and submission, it causes its own death, so far as it
is able.'' Baader expressly calls attention to the fact that Hegel confounds ``the immediate that
\emph{raises me up}, and that which \emph{I abolish} --- the life-giving breath of spirit from
above, and the earthly nourishment.''\footnote{Franz v.\ Baader: \emph{Fermenta Cognitionis}.
Berlin 1822. First Number, pp.~39; 67. --- Martensen's standpoint recalls, for the rest, not only
Baader, but also Jacobi, by whom Bishop Mynster --- admired by Martensen as a religious genius ---
was so strongly influenced. Compare above the section on Tyge Rothe, in whom we have also met the
influence of Jacobi.} That is to say precisely: Hegel (and Heiberg with him) do not distinguish
between object and authority. And philosophically they are here surely in their good right. It was
in reality not a ``higher'' philosophy that Martensen found; it was a transition from philosophy,
which has objects but no authorities, to the theology grounded not merely on authority, but on a
confusion of object and authority. With very little right did Martensen look down on
\emph{supranaturalismus vulgaris}. That he wished to keep Hegel's method was an inconsistency, and
one does not notice much of it in his Dogmatics either; he has silently let it fall, although he
likes to operate with ``higher unities,'' especially when he is in a difficulty.

In the lectures that he delivered at the university in the [eighteen-]forties, and which, partly by
his spiritedness and eloquence, partly by the great promises of a higher unity of faith and
knowledge, made them some of the most attended and admired in the history of our university, it was
now precisely his task to show that there was a higher speculation than Hegel's, a speculation to
which one got access only on the foundation of the Christian faith. He himself says in his
\emph{Life} of these lectures that when he was to lead his hearers in to his problems and his
solution, it had to happen through a survey of philosophy from Kant to Hegel. He had first to
inspire them for Hegel --- and then to lead them out beyond Hegel. He had need of the school of
thinking that he found in no one so comprehensive as in Hegel. He used the Hegelian formulas to
develop a view that was to be very different from the Hegelian. He himself admits that in many
cases he did not succeed in exorcising the spirits he had conjured forth. Many stopped when he had
led them to Hegel and would not go further. He consoles himself that many, who could later be
reckoned among the Danish Church's foremost men, faithfully followed him all the way through.

A young theologian from those days, who did not hold fast to speculative theology, though he in his
later days too embraced Martensen's personality with admiration and love --- the poet and clergyman
Johannes Fibiger --- has given a depiction of the studying youth at that time and of Martensen's
lectures.

``Theology was then regarded as a mistress as honourable as she was gracious to serve, well able to
dazzle. I too soon came high up in full flight on her fantastic wings, dangerous enough. It lay so
in the time. One must have lived in that time to be able to form a notion of its strange character
--- indeed, merely to be able to believe in the possibility of it. Under the autocracy of German
philosophy, all thoughts strove to build on the Babel-tower of fantasying; what we heard around us
was nothing less than that every great orator set his nearest life-task in building a still higher
one. The universe with all its great and small secret chambers was searched out and cleared up in
the concept, all riddles were solved, Hegel and his Berlin band of disciples had accomplished the
work. But thereby the concept of art too was unveiled and the prize laid out. It must be able to be
imitated and brought still further; one had only to build one's own system and go beyond Hegel, and
oneself become a magnate of the spirit-world. We soon breathed no other air; the only thing we
heard at the university that rang like a voice of the future was Martensen's brilliant speech, and
there was not a little in it of airy stones for that kind of building. Theology\ldots\ had then its
period of splendour. All spirits swore by theology; it was the Holy Ghost of spiritedness and had
won itself a new shape, was neither scriptural learning nor the Church's doctrine of faith nor a
priest-school --- such things were discarded stages; --- it was speculative. This mystical word,
which appeared with the ring of the ``Absolute,'' gleamed through every one of our master's
power-sentences like the flash of the sun in a fountain. Whence it came, we did not know; it was
past like a piece of magic. But when it came to some better understanding, it showed itself to be
the art of being able, with an eclectic aggregate of German systems, to bring the mysteries of
Christianity under the daylight of the concept, as surely as if one had an angel set in spirits
placed forth for analysis. Credo, ut intelligam, was the formula. Faith is ``posited,'' the angel
is set forth. Thereupon intelligence begins to let its rays fall through all its heartbeats,
nerves, and muscles. Everything was made perfectly clear in a speculative way; only it was not
quite clear how it would go if the shy Genius, unpleasantly touched by intelligence's cold
instruments, should begin to vanish.''\footnote{Johannes Fibiger: \emph{Mit Liv og Levned, som jeg
selv har forstaaet det} [My Life and Days, as I Myself Have Understood It]. (1898), p.~73 f. ---
Compare the depiction of Martensen in his later days, pp.~323--327.}

This last remark touches the decisive point. Martensen wished to philosophize on the foundation of
positive religion's exalted picture-language and let his speculations profit from the mood that the
handed-down images call forth. He did not, like Hegel, make earnest of letting the concept replace
the image. He wished ``not to lead away from the representation [the image] in order to come to the
concept, but on the contrary to seek to clarify the representation through the concept, to present a
unity of concept and intuition.'' But is it possible to reach such a ``unity''? For the concept,
every image, every intuition is only an analogy or an example; but the religious representations are
supposed to be revelations, irreplaceable by other representations. By that alone is revealed
religion different from poetry. Martensen stands then also waveringly here. He gives no higher unity
of thought and image, but gives now a thought, now an image. And as a theologian he agrees with
Grundtvig that ``with the picture-language one must not play''; consistently, when the images are
supposed to be revelations; --- but what, then, is the concept really for? In his old days he
stated that of that which lies beyond experience one can speak only in images. ``If we would not
speak of these things in images, then we would only come out into the empty, would come to some poor
concepts that would give us no light. Let us therefore hold fast to these images, which God's Word
gives us\ldots\ They are images, and yet they are more than images, for they designate the highest
realities.''\footnote{Martensen: \emph{Af mit Levnet} [From My Life], II, pp.~7; 54. --- III,
p.~240. The last-cited statement is from an Easter sermon of 1881.} Whence he knows this last, we
are not told; but it is clear that the higher unity of concept and image is given up. Martensen has,
to use Fibiger's simile, held fast the angel and laid the instrument aside.

But even this he does not always do. The year after the appearance of the last-cited statement, he
published his book on Jakob Böhme (1882), that mystic to whom he had been led by Baader, and in whom
he had always taken an interest. Occasionally he here comes to speak of Jesus' Ascension and its
compatibility with the modern conception that ``heaven'' is not a physical place in space. He
asserts that ``heaven'' means ``the inner ground of the outer world,'' and that the visible
Ascension was only a sign for the disciples that their master had left the world of space! Here he
has let the concept reinterpret the image, and that in a way that recalls the old rationalism, which
was, after all, supposed to be an overcome standpoint.\footnote{Compare more closely, on the
relation between the scientific conception of space and the biblical representations, my
\emph{Religionsfilosofi}, II A b. --- Compare also above on Heiberg's and Grundtvig's relation to
the Copernican system. In Rasmus Nielsen we shall meet the same question again. None of them was
occupied with the problem in the way Holberg was. Romanticism had on the whole a tendency to brush
it aside.} --- The profound cobbler from Görlitz, whom Martensen so enthusiastically honours in the
said book, was a freer and more consistent thinker than Martensen, and he must therefore in several
places hear that now he does not think in a Christian way, but in a heathen way. Martensen has
overlooked a statement of Böhme's, in which the latter says, against a similar objection: ``Ich
schreibe nicht heidnisch, sondern philosophisch!'' [I write not in a heathen way, but
philosophically!] The distinction between Christian and heathen ``thinking'' does not belong in
philosophy, but in the history of religion. ---

It is not the task here to follow Martensen through his theological works, in which philosophy came
to play an ever smaller role. The type of his thinking is, however, everywhere the same: the
invocation of a higher unity. As he had wished to find a higher unity of intuition and concept, so
he wished as a churchman, following Mynster, to maintain a tendency that designated a higher unity
in relation to Grundtvigianism's ``one-sided objectivity'' on the one hand and the Inner Mission's
``one-sidedly subjective tendency'' on the other. It could then rightly be called a high-church
tendency, a play on words that must not, however, be charged to Martensen's account. And finally he
wished, in his \emph{kristelige Etik} [Christian Ethics] (1871--1878), to demonstrate the union of
the Christian and the human, the great ideal that, as we have seen, stood before him right from his
earliest youth. But this unity can, he maintains, be reached only by unconditional submission to
Christianity. This does not, however, seem to be exactly a union, in any case not a higher unity,
but a subordination of the one member of the opposition under the other. Whether a harmony of this
kind is at all compatible with Christianity itself was a problem that Martensen never seriously
posed to himself.\footnote{In the weekly \emph{Nær og Fjern} [Near and Far] I wrote in the year
1878 a longer critique of Martensen's ethics. Martensen's \emph{Levnet} (III, p.~208) says, no
doubt with reference to this review, that the book was attacked ``from the purely negative side.''
I thus know from my own experience how one feels when one is ``purely negative.'' --- Of Martensen's
book on Jakob Böhme too I wrote a review at the time (\emph{Ude og Hjemme} [Abroad and at Home].
1882).} So much the more sharply was it posed by the thinker to whom we now pass over, and whose
task then also brought with it that he came to enter into the sharpest, most passionate opposition
to Martensen both as thinker and as churchman.

\begin{center}---\end{center}

% ============================================================
%  14. SØREN KIERKEGAARD   (Da.: "14. SØREN KIERKEGAARD")
% ============================================================
\dfchapter{14. Søren Kierkegaard}
\label{ch:df-kierkegaard}

\noindent Kierkegaard occupies within Danish thought a place that can be compared with
Pascal's position within French thought. Common to them is the deep feeling of the
oppositions of life, the sharp emphasis on the great problems, and the conviction that,
however great a horizon thought opens for us, our highest tasks nevertheless fall within a
narrow circle, in the end within our own personality, in the most inward reckoning of our
relation to the truth --- the truth we have courage and honesty enough to acknowledge.

His view of life forms a decisive opposition both to Grundtvig's glad security, which
revels in history, images, and song, and to Martensen's speculation, which lets faith,
presupposed as given, unfold itself in dialectical movement.

Philosophy has no right to appropriate everything in this personality and its work, which
for a large part belongs in the domains of religion and literature. But philosophical
thinking nevertheless forms the form in which he in the end sought to give an account of
his activity as a whole. His chief philosophical work, \emph{Afsluttende Efterskrift til
Filosofiske Smuler} [Concluding Postscript to the Philosophical Fragments] (1846), he
himself declared to be the midpoint of his whole production; on both sides of this central
work lie the religious and the aesthetic productivity.\footnote{\emph{Af Søren Kierkegaards
efterladte Papirer} [From Søren Kierkegaard's Posthumous Papers]. 1849, p.~69. --- Yet it
is precisely in the diaries from this year that Kierkegaard occupied himself with the
question whether he should not have the right to maintain that he had from the first had a
conscious religious aim with his whole production, the aesthetic and philosophical as well
as the religious. He came to the result that he dared not maintain it. The production in
the various directions had taken place involuntarily, and he dared in any case not ascribe
to himself the credit that it in its totality could serve a religious aim. See on this my
\emph{Mindre Arbejder} [Minor Works], Second Series, pp.~183--186.}

What he has contributed as a philosopher can be brought under three points of view. He has
given a sketch of a theory of knowledge, in which he takes up again Kant's problem, which
Hegel had let lie. He has, in his doctrine of the ``stages,'' given a comparative
philosophy of life on the basis of a deep understanding of the essence and need of
personality. And he has posed the problem of the relation between Christianity and humanity
in as definite and principled a form as no other thinker in the nineteenth century. On all
three points his thought gets its passion from the fermenting and deeply gripped
personality from which it springs. With violent speed it rises to the highest peaks, where
it is all or nothing, and the oppositions are strained so strongly that not only thought,
but the personality itself, is in the end burst asunder.

The abiding instruction he can give is predominantly of an indirect kind. So it was,
indeed, with Socrates too, one of his great models. But he has done a work of thought that
can be of great significance even for him who in essential respects chooses for himself
quite other starting points. And directly, positively, he has maintained the inner
connection of the view of life and of thinking with the personality, and thereby pointed to
a constant source of purification and renewal in the world of spirit. The concept of
personal truth set up by Poul Møller he has carried out into its consequences.

His course was not long (he was born in 1813 and died in 1855), but his passionate
productivity brought it about that he got said what, according to his personality and his
standpoint, he had to say, and death came when he had just said the last word that could be
said from his standpoint.\footnote{A fuller characterization of Kierkegaard's personality
and doctrine than I can give here I have given in my book \emph{Søren Kierkegaard som
Filosof} [Søren Kierkegaard as Philosopher] (1892).}

It is not only his doctrine, it is his very personality that has philosophical interest by
the psychological riddles it contains. In my view, which I support on the study of his
posthumous papers, melancholy [\emph{Tungsind}], nourished by a joyless childhood and a
pietistic upbringing, is the deepest current in his being. He felt himself constantly as a
clipped-wing bird that, instead of being able to use its wings freely, had to move in its
narrow circle, without being able to solve the riddle whether it was sickness or guilt that
constantly held it back. Melancholy removed him from life and made constant and close
relations to other people an impossibility for him. It gave him a different view of life
than the ordinary, easy-living people had; he saw abysses that they passed carelessly by,
and he discovered elements in Christianity that they left unheeded. But alongside melancholy
stood then the genial urge to production, the great gift of poet and thinker. It nourished
the melancholy by spinning out broodings and endless reflections, in constant circling
about itself; but it served at the same time as an outlet and a relief, in that he could
forget the inner darkness by immersing himself in production. He kept himself, as he
expresses it, alive, like Scheherazade (\emph{The Thousand and One Nights}), by constantly
composing. --- Yet still a new opposition here forced its way forward. Reflection and
imagination and their urge to spread came into strong contrast with the urge, begotten by
melancholy and deep unrest, for an absolute hold, for an anchor of distress on the troubled
sea of life --- for submission to an absolute authority, a supernatural power. Despite all
reaction against the joyless childhood and the strict upbringing, he was not able to tear
himself loose from the view of life to which they had laid the foundation; the wing was,
after all, clipped, and that especially from inner, not from outer causes. But against the
fixed authority the movements of thought dashed again and again, only to be constantly cast
back anew. His life of thought is characterized by what he states somewhere: ``The highest
potency of every passion is always to will its own downfall, and thus it is also the
understanding's highest passion to will the offence, although the offence must in one way or
another become its downfall. This, then, is thought's highest paradox: to will to discover
something that it cannot itself think.''\footnote{\emph{Filosofiske Smuler} [Philosophical
Fragments] (1884), p.~52. (\emph{Samlede Værker} [Collected Works], IV, p.~204.) --- The
expression ``offence'' [\emph{Anstød}] I take to be a reminiscence of Fichte, in whom
``Anstoss'' means the barrier that thinking constantly runs up against, where it stands over
against something merely given, something new that is not a consequence of the preceding
work of thought itself.} --- This urge to go to the boundary every critical philosopher must
have who will follow the Socratic command to distinguish between what he knows and what he
does not know. But it got in Kierkegaard a quite special force, because his interior ---
partly on account of the original melancholy, partly on account of the religious authority
rooted in by the upbringing --- presented an opposition between something that was to stand
absolutely fixed, and a never-resting striving nevertheless to draw everything into the
illumination of thought. Thereby he stood from the first otherwise toward the religious
problem than men like Sibbern, Heiberg, and Martensen. Sibbern took the Christian as
something given; he reserved to himself to test how far it really was \emph{given}, and in
the course of his long life he found that everything dogmatic in Christianity was not at all
\emph{given}. Heiberg took religion as an object that through dialectical development was
transposed into pure thought, precisely as regards the dogmatic. And for Martensen religion
was both object and authority, without its being clear when it left off being the one and
began to be the other. For Kierkegaard it was only authority, a holy, venerable authority,
mild and consoling, but with a Medusa-face toward him who would treat it as an object of
cognition, as well as toward him who would knock something off its ideal demands. And this
latter point of view stepped more and more into the foreground for him, the less the age,
and especially the contemporary Church, would acknowledge its justification. Thus he was
step by step led to raise the most violent spiritual strife that has ever been carried on in
Denmark. When he had said his last word in this strife, death closed his eyes. ---

Such was the psychological soil on which Kierkegaard's philosophizing grew forth. Naturally
he has oriented himself in the literature and philosophy of his time. His relation to Poul
Møller has already been touched on. After his theological studies he occupied himself
thoroughly with Greek philosophy and composed a spirited dissertation for the doctorate,
\emph{Om Begrebet Ironi, med særligt Hensyn til Sokrates} [On the Concept of Irony, with
Special Reference to Socrates] (1841). He had had a Hegelian period, but had, probably under
the influence of Poul Møller, gotten beyond it. Socrates became his model on account of the
close union he presented between thinking and personality. When he nevertheless here
presented Socrates as an abstract and negative ironist, it was a conception that he later
expressly corrected, just as the concept of irony too underwent a change in him during the
working-through of the doctrine of the stages (the views of life), as we shall show in its
place. Of his contemporaries' philosophers he set, after having left Hegel, the acute Adolph
Trendelenburg very high, and states in several places with warmth how much he owes him. ---
Of older ones he felt himself especially drawn to Lessing and Hamann.

\begin{center}---\end{center}

Kierkegaard's theory of knowledge he himself sums up in the two propositions: A system of
thought can be given --- but a system of existence cannot be given.

There is an intervention of thoughts that we necessarily must use when we would understand
existence. But these thoughts (the categories) perhaps themselves spring from existence, are
a kind of abbreviated expressions for our experiences! Or should they be able to come into
being in us without having anything whatever to do with real existence? Kierkegaard gives no
answer to this, but announces a closer investigation, which he did not, however, come to
undertake, since the religious authorial activity soon laid full claim to him. But in any
case he demands that those fundamental thoughts be developed in strict logical form, without
one's taking up presuppositions for whose origin one does not give an account. --- And one
must at the same time make it clear how the beginning of the series of thoughts one develops
is made. An absolute beginning it cannot be; for from nothing comes nothing. And if one
presupposes a reflection that analyses a given content, then it must be asked why this
reflection stops at a certain point in order to begin building up the system of thought from
there. Why does reflection not continue its analysing work instead of stopping and beginning
the construction? Only by a decision, thus with a jerk or by a leap, can the first stone of
the system be laid; otherwise one keeps on digging. There is thus not a necessary transition
from experience and analysis to the system of thought. Hegel could [this seems to be
Kierkegaard's thought] very well have spent his life continuing the \emph{Phänomenologie des
Geistes}, without ever having come to write his \emph{Logik als Wissenschaft}. --- Finally
it must be remembered that even if it succeeds to construct a system of thought in which each
of our fundamental concepts finds its place, there still remains the question to be answered
how this system relates to the philosopher's own personality, which, after all, cannot
wholly go up into it.

A system of thought thus leaves, even if it can be reached, various questions behind. And a
system of existence cannot be given, since existence for us, who exist in time, can never
come to lie present as concluded. The past is indeed concluded, but we, who exist and think
about the past, are not concluded; time runs on while we think, and makes its demands on us.

In \emph{Filosofiske Smuler} the question is discussed at length whether the past is really
more necessary than the future, and it is maintained that the more one makes oneself
\emph{contemporary} with the past, the more one will become clear that, during the
experiencing, it constantly stood only as possible, thus as uncertain --- as an object of
faith. Not even toward the past, which is, after all, present as concluded, can there then
be given an absolute knowledge, when one makes it rightly present to
oneself.\footnote{One sees how Kierkegaard here comes back to Hume's problem, whether the
transition from one moment to another was necessary or not. Hume too rejected the appeal to
the past (though especially from the point of view that he denied the right to infer from
past to future, while Kierkegaard seeks to become contemporary with the events of the past,
to make the past into the present). --- Kierkegaard agrees with Hume that necessity cannot
be observed. But he commits the inconsistency of thinking (in \emph{Begrebet Angst} [The
Concept of Anxiety]) that ``the leap,'' the absolute interruption of connection, could be
observed. See on this \emph{Søren Kierkegaard som Filosof}, pp.~78--81.} Lessing gets it
right: not the possession of the truth, but the constantly continued striving after truth
is the highest a person can reach. In existence, thinking and being do not cover each other.
We live forward, but understand backward. --- Only for an eternal being does existence
constitute a concluded system, not for beings who stick in the middle of time. ---

In philosophy one has too early left Kant's honest way. Hegel has not, as he and his
adherents think, refuted Kant. The difference between thinking and reality constantly
remains in force. --- Kierkegaard is one of the first who has pointed to a resumption of the
investigations of the critical philosophy. Criticism had not disappeared in the romantic
period; it maintained itself as an undercurrent through such capable researchers as Herbart,
Fries, and Beneke, not to speak of Schleiermacher and Schopenhauer. Kierkegaard has scarcely
any real relation to these thinkers. He mentions with acknowledgment Schleiermacher's
doctrine of faith (in \emph{Begrebet Angst}), and he read in his last years (see the diaries
from the [eighteen-]fifties) Schopenhauer's writings with interest; but in an epistemological
direction he is scarcely influenced by them. ---

It is not only against speculative philosophy, it is also against theology in its various
forms that Kierkegaard turns.

The intellectual interest that goes up in winning scientific understanding and demonstrating
connection between our experiences is for Kierkegaard not the highest. Proceeding from the
fact that psychic life expresses itself in a comprehending, a synthesis, he finds that this
synthesis is so much the more powerful, the greater the tension and opposition between the
elements. And greater is the tension in him who lives in longing and striving toward an
eternal goal, and who thus is to comprehend time with eternity in restless expectation, than
in the contemplative man, who works at bringing connection between various data at hand.
Therefore faith stands higher than knowledge, and therefore Kierkegaard wondered that
Martensen could sit calmly and ``arrange a Dogmatics,'' while he assumed faith as a given
matter. In science we think only with a part of our personality, but in every serious view
of life it is the whole personality that thinks, and all the soul's powers work together.
Kierkegaard's measure of value, to which we shall come back in his doctrine of the
``stages,'' is the fullness and the degree of opposition that can be connected in the unity
of the personality. He has here descended to a thought that not only since Leibniz and Kant
lies at the foundation of the theory of knowledge, but that must also make itself felt in
every energetic way of taking life. --- For my own part it was this thought that gripped me
in my youth and held me fast, so that I again and again return to it. I see especially on
this point Kierkegaard's abiding significance, however much my view of life has removed
itself from him, and however little I especially agree in the distance at which he sets the
scientific work from the rest of the personal life. ---

Purely theoretically, according to Kierkegaard, the religious problem, especially the
question of the truth of Christianity, cannot be solved. For here there is nothing given
that can be transposed into concept; here there is talk of a distant goal, of whose
attainment no objective (scientific) certainty can be had, but to which one can only relate
oneself in faith --- in the again and again repeated confirmation and assurance in which the
unrest of the tension is overcome. The object of faith exists only for the believer, not for
him who seeks universal concepts and laws.

Just as little does it avail to refer to the Bible as a historical document. For when does
the historical-critical investigation become finished? and can it --- even in the most
fortunate case --- ever become so certain that one can build the most personal decision on
it? --- Nor does it avail to refer to the Church or the confession of faith handed down in
the Church. For who guarantees its genuineness and originality? Here is a new playground for
historical research. But the decision cannot lie there. No objective guarantee whatever is
possible.

The decisive thing is constantly in the end the personal need for a hold --- for an anchor
of distress under the strong surge of temporal existence. The more the question stands
undecided for the objective, scientific (philosophical and historical) consideration, the
more it becomes clear that --- if it is possible --- it must lie in the inwardness of the
personality. Subjectivity is the truth. It is the individual's relation to that which is
believed that is the main thing: it depends on whether the individual himself is in truth,
even if that which he believes is untrue. He who prays in untruth (impersonally, without
inwardness) to the true God prays to an idol; he who prays in truth (with his whole soul's
inwardness) to an idol prays to the true God. It depends on \emph{a How}, not on a
\emph{What}. The highest truth that a being existing in time can reach is characterized by
objective uncertainty, held fast in the most passionate inwardness's appropriation. ---

Kierkegaard has here carried Poul Møller's idea of personal truth out into its utmost
consequence. He has at the same time, remarkably enough, come to state a thought similar to
the one that Heiberg, from quite other presuppositions, stated in the lines cited above.

Without entering upon a more special critique, one cannot, however, refrain from here making
two questions, the one from Kierkegaard's own standpoint, the other from a universally human
standpoint. --- How can there be talk of a revelation, a real disclosure of the secret of
existence, when the truth is not a What, but a How, and when the main thing is the
personality's inward striving, not its definite content?\footnote{Compare on this point the
interesting essay by the thorough Kierkegaard-scholar Chr.\ Schrempf: \emph{Kierkegaards
Stellung zu Bibel und Kirche} (\emph{Zeitschrift für Theologie und Kirche}, 1891).} One
understands at this point Kierkegaard's double-sided effect: some he has led deeper into the
relation to Christianity; others have precisely under his influence been led away from
Christianity, because only thereby could they preserve their personal truth and inwardness.
--- The other question is whether it should really be so that the soul's whole energy must
be used up in securing itself a defence. To this question we shall return in the discussion
of Kierkegaard's ``Stages.''

\begin{center}---\end{center}

Personality (subjectivity) appears in various forms, on various foundations. Already in the
foregoing there have shown themselves two types, the one characterized by purely intellectual
interest, the other by religious interest. In wishing to impress the significance of personal
truth, Kierkegaard depicts a whole series of types, or, as he with a less felicitous word
calls them, stages. When there is not found a special intellectual or intellectualistic stage
among them, it is surely because such a stage for him meant merely a volatilization of the
personality. He will see the individual in its relation to real life and depicts the various
ways of taking this.

The various stages or types stand in Kierkegaard sharply cut off in relation to one another.
One does not learn how they come into being, or how a transition from one to the other is
possible. In the Hegelian philosophy the transition from one stage to a higher took place
merely by raising the lower stage to full consciousness; reflection itself, the clear
consciousness of where one stood, was to be enough to realize a higher stage. Kierkegaard
maintains that the transition presupposes a leap, because the standpoints are in principle
different --- which is, after all, the same as saying that one cannot understand the
transition.

Kierkegaard's presentation of the ``stages'' is somewhat different in the poetic writings
(\emph{Enten--Eller} [Either--Or] and \emph{Stadier paa Livets Vej} [Stages on Life's Way])
and in his chief philosophical work (\emph{Uvidenskabelig Efterskrift} [Unscientific
Postscript]. The appendix to the second chapter of the second section). In the former he has
only three stages, the aesthetic, the ethical, and the religious. Here he gets six, in that
irony is inserted between the aesthetic and the ethical stage, and humour between the ethical
and the religious stage, and in that the religious stage is divided into two, heathen-human
and Christian religiosity. --- We shall now give a characterization of these six stages.

\emph{The aesthetic stage} is depicted somewhat differently in \emph{Enten--Eller} (1843) and
in \emph{Stadier paa Livets Vej} (1845). In the former the characteristic thing is a
reality-distant dwelling in melancholy and fantasying with the painting of moods, refined
reflections, and Mephistophelian thought-experiments. It is the imagined author of the first
part of \emph{Enten--Eller} who represents the aesthetic stage, not merely the figures (for
example Johannes the Seducer) that he ``in his groping existence'' conjures forth. Yet a
single one of the sections in the first part, the one called ``Vexeldrift'' [The Rotation of
Crops], is typical of this stage, also as it is depicted in the first part of \emph{Stadier
paa Livets Vej}, where ``the Aesthetes'' are presented existing in flesh and blood, not as
mere possibilities.

The characteristic thing is on the whole a shifting life of moods that shuns fixed forms and
relations, a wanton playing with all relations of life. One takes everywhere the new, fresh,
piquant, momentary, without letting oneself be held fast. One enjoys the exciting without
being entangled, the sensuous without falling, and has elasticity to break off at the right
moment. The great enemy is repetition and duty. One wants fleeting contact with people, and
one tries occupations; but one takes only the flower of everything, and avoids marriage and
friendship, citizenship and science. It is especially the erotic that is used as an example
to illuminate the fleeting touching of the circles of life that is peculiar to this
standpoint. Everything is made into an image and enjoyed as if it were only an image. It is
on account of this reflection that this stage gets the name of the aesthetic. It is not a
poetic stage; there is not depicted a poetry of life that does not itself know that it is
poetry, because it unfolds itself as an involuntary consequence of immersion in life. In his
own copy of \emph{Enten--Eller} Kierkegaard has written: ``A real being-in-love could not be
used in the first part\ldots\ What I could use was a manifoldness of erotic moods.''

\emph{The ironist} does not go up into the single moments; his task is to preserve his
connection, to maintain his self over against the various relations he is placed in, and
which are to him not means to enjoyment, but on the contrary dangers that could hinder him
from being himself. He goes then with apparent seriousness into the individual situation and
seems to be taken up with it and to ascribe value to it; but it is a tribute he pays in order
to be able to preserve himself over against externality and dispersion.

The ironist has no positive content or aim. He relates himself essentially defensively in his
striving to preserve his freedom. But thus he is an intermediate form between the Aesthete,
who relates himself enjoyingly, and the Ethicist, who has a positive task. In his dissertation
Kierkegaard had treated irony (and Socrates as its representative) as purely formal and
negative; but even there he had, in one of the appended theses, stated that as philosophy
begins with doubt, so a life that is to be able to be called worthy of a human being begins
with irony. Irony can be the incognito of ethics, its outer form, since ethically the first
condition is to have a midpoint and not let life be dissolved into single moments. There is in
irony a passion which indeed essentially expresses itself in the abstraction from all finite
relations, but which nevertheless leads to maintaining an opposition between the inner and the
outer, which is of so great significance for an ethical standpoint. ---

\emph{The ethical stage} is characterized in the poetic presentations as gathering in
opposition to the Aesthete's dispersion, and as open devotion in living-together in opposition
to aesthetic dreaming and experimenting. The ethical presupposes a decision that takes up with
repetition, the one the Aesthete feared. Thereby there comes connection into life, and the
isolation is abolished. It is most nearly Hegel's conception of the ethical as absorption in a
universal, in a life of society, that Kierkegaard here lays at the foundation, only that it is
the family, marriage, that he uses as type, not the state. Bourgeois life Kierkegaard could
see only under the point of view of philistinism. --- In the philosophical presentation the
conception has become another. The concept of the individual [\emph{den Enkelte}] plays a
chief role; from a preliminary application in the foreword to one of Kierkegaard's religious
discourses it has gradually moved forward into the foreground of his line of thought, and the
ethical stands now for him as the individual's affair, as calculated for individuality. In
one's conscience one has only to do with oneself. And all ethics is religious. Without relation
to God and belief in immortality, the ethical is only custom and use. The social and historical
have no significance, and marriage is only a hindrance. Thus the religious character of the
ethical goes into one with the individualistic. The ground of this displacement in the
conception must be sought in a line of thought that Kierkegaard had developed in the little
work \emph{Frygt og Bæven} [Fear and Trembling] (1843), where the possibility was discussed
that a religious command could conflict with ethics' universal law. The story of Abraham's
obedience toward the command to sacrifice his son is used as example. In such a conflict
between the religious and the ethical, the individual stands as the single individual, isolated
from all others; no one can understand him and his conduct. Through a violent break must the
transition from the ethical to the religious take place, when the ethical means the universal,
that which holds for all. Involuntarily, then, Kierkegaard later laid the concept of the
individual at the foundation already in ethics; the religious character of the concept has
followed with it --- and thereby, really, at the same time, the ethical is abolished as a
special stage. Now poetry and religiosity become the two only ideal possibilities: ``Between
poetry and religiosity worldly life-wisdom performs its vaudeville. Every individual who does
not live precisely poetically or religiously is stupid.''\footnote{\emph{Uvidenskabelig
Efterskrift}, p.~349. (\emph{Værker} VII, p.~397.)} I do not take it to be accidental that
Kierkegaard here speaks of the poetic and not of the aesthetic; he is thinking of something
more immediate and involuntary than the type he has characterized as ``aesthetic.''\footnote{In
\emph{Stadier paa Livets Vej} (\emph{Værker} VI, p.~158) it is said (in the section on the
Ethical): ``Woman is in her immediacy essentially aesthetic, but because she is so essentially,
therefore the transition to the religious is also close at hand. The feminine Romanticism is in
the next moment the religious.'' --- I take it that ``aesthetic'' here means poetic, and is not
taken in the same sense as when Kierkegaard speaks of ``the aesthetic stage.'' In any case one
sees that as regards woman the ethical stage does not exist at all.} ---

The Ethicist goes up in his tasks, which are to be solved by his own activity. Therefore he has
his centre in himself and seeks in a positive way to fill his place. For \emph{the humorist} his
own existence, the whole circle in which he moves, and on the whole everything finite, is
vanishing in relation to the Infinite and Eternal. He moves with his centre as an epicycle on a
larger circle, from whose midpoint he is infinitely distant. Therefore finitude becomes comic
when it will at all mean something. Even his suffering and his guilt the humorist sees as
vanishing over against the infinite horizon that opens for him. Humour encloses melancholy at,
and sympathy with, the pain in life, and its smile and jest hold a deep seriousness. The ironist
will maintain himself through the apparent seriousness toward the finite; but the humorist's
seriousness does not hinder him from a mild jest with everything finite, which despite all
busyness and importance can nevertheless never become of decisive significance. Humour is the
last stage before the religious. ---

Kierkegaard has here hinted at a view of life that corresponds to the fundamental mood that
thought's understanding of existence awakens. The individual thing is seen in its distinctiveness
and its value --- and yet as vanishing in relation to the unsurveyable existence that thought
opens for us. Oppositions can appear in their sharpness and yet fade away in relation to the
assurance of a deep harmony, a great law of unity. It is the mood in which the single mode in
Spinoza's system must be placed over against infinite nature. It is for my part the one of the
Kierkegaardian stages with which I would most nearly feel my own standpoint and my own mood
toward life akin. It is a pity that Kierkegaard, with the abilities that stood at his disposal,
has not carried out the depiction of it further. Precisely since his aim with the whole
presentation of the ``stages'' is to show ``what an enormous existence-range is possible outside
Christianity,'' that is to say, what a high and noble view of life is possible on purely human
ground, he ought to have characterized this last purely human stage more penetratingly. He lets
several of his pseudonyms\footnote{Kierkegaard has presented the ``stages'' by means of a series
of pseudonyms, because he thereby could attain ``the psychological consistency that no factually
real person, in reality's ethical limitation, dares allow himself or can wish to allow himself.''
(``A First and Last Declaration,'' after \emph{Uvidenskabelig Efterskrift}.)} speak in its name
--- but they all have direction toward the fifth and the sixth stage, whose supremacy they
acknowledge, although they do not themselves belong to them. For the same reason it is that the
earlier stages are not carried out so clearly and completely as the task demanded. It is,
besides the ``aesthetic,'' really only the religious stages that get a complete carrying-out ---
and, as we have seen, Kierkegaard gradually became inclined to push out all the intermediate
forms.

\emph{The religious stage} enters into force when one's own existence is not merely seen as
vanishing in relation to the Eternal and Infinite, but when the individual sees his real task
in making himself one with this Eternal, having his midpoint in the midpoint of the larger circle
in whose periphery he lives his life. It gets then two centres, has its life both here and yonder,
but is really to have it yonder. There is felt a deep cleavage, a painful contradiction. The pain
cannot here, as in humour, fade away in melancholy, but is constantly connected with temporal
existence on account of the opposition between time and eternity. The individual feels like a fish
on dry land. It lives in a foreign element.

``Not the bird in the cage, and not the fish on the beach, and not the sick man on the sickbed,
and not the prisoner in the narrowest prison, is so captive as he is who is captive in God; for
as God is, so the captivating representation [of God] is everywhere and in every moment.''

A standpoint like this is possible within heathendom and must indeed be distinguished from the
Christian, although many so-called Christians do not even reach up to it. In \emph{Christianity}
the opposition between time and eternity is sharpened in such a way that it can be abolished only
by the Eternal itself entering into time, living in the form of temporal existence as a finite
being --- by God becoming man. Here the truth becomes a paradox: the Deity in time! Man can on
his way through time meet the Deity itself: everything depends on whether he then knows him and
believes on him; --- but the very ability to know and believe he cannot give himself; he must
receive it from the Deity itself --- and it depends then on whether he will receive
it!\footnote{In \emph{Filosofiske Smuler} (1844) the paradox in the Christian fundamental dogma
of the God-man was presented in sharp opposition to the Socratic (human) standpoint.} --- The
whole view of life is thereby determined. For over against the one decisive point, the Deity's
revelation in time, everything else loses its significance. Only on this one point is time
touched by eternity; Christ is --- as all ideality at bottom is --- a tangent to the Earth, to
reality.\footnote{It is characteristic of Kierkegaard that both the stage he sets highest and the
one he sets lowest can be determined by the same formula: to be a tangent to the circle of life.
Like Christ, indeed all ideality (see \emph{Efterladte Papirer} [Posthumous Papers]. 1849,
pp.~30; 242), so too the aesthetic way of taking life can be described by its going off along the
tangent! (See especially ``Vexeldrift'' in the first part of \emph{Enten--Eller}.) And it is
perhaps especially characteristic of him that it is these two standpoints that he best
understands and has most respect for.} The whole rest of existence is God-forsaken, worthless.
All value is concentrated on a single point. --- From this follows again that only a
psychological, neither a speculative nor a historical, introduction to Christianity is possible.
It depends on how much has been lived --- whether the man has learned to despair of everything
else in order then to find his hold in this one point, the only one from which light falls over
existence.

This last thought Kierkegaard developed more closely in \emph{Sygdommen til Døden} [The Sickness
unto Death] (1849). He investigates in this his most profound work the various kinds of despair.
The immediate joy and happiness seem to be furthest from despair, and yet the security and calm
are illusory. The anxiety lies behind. For there is set to man the great demand to be spirit
according to the absolute, divine measure. Christianity will make man into something so
extraordinary that he cannot get it into his head. This is the ground of the offence. Weakness,
closed-up-ness, defiance are various forms in which man withdraws himself from the high demand;
all are they forms of despair, and there conceals itself in them all a will that will withdraw
itself from the demand, now by self-surrender, now by self-assertion. Spiritlessness too, which
seems to be too silly to be able to be called sin, arises nevertheless only by man's own fault.
The very consciousness of sin can take on the form of despair. The right opposite of sin is not,
as the heathen thought, virtue, but faith. And to faith one comes through despair, which teaches
one to flee to grace. Despair is therefore a sickness that it is the greatest misfortune not to
have had. But Christendom's concept of sin does not go deeper than the heathen's. One has levelled
out the great oppositions sin and faith, despair and grace. The levelling has taken place first
speculatively, then rabble-like.

It was Kierkegaard's task in the said work to show how deep-going inner experiences, how sharp
oppositions, he must know who is with truth to be able to apply the New Testament's psychology to
himself. No rendering of the work's content can give any notion of with what passion and with
what sharp-sightedness it is written. The master leads us along abysses and makes us know spiritual
possibilities that he himself dared not give out as having tried, but that, according to his
conviction, he must know at first hand who with right was to be able to call himself a Christian.

While \emph{The Sickness unto Death} impresses the psychological experiences that Christianity
presupposes, \emph{Indøvelse i Kristendom} [Practice in Christianity] (1850) lays the weight on
the meeting with ``the Deity in time.'' One relates oneself to Christ only by being contemporary
with him. From history one can learn nothing. It is constantly the abased, not the exalted, who
speaks to us. He has, after all, not yet come back, and the later generations can have no advantage
over the earlier; the task for them all is to acknowledge Christ as the abased, as the lowly man
who nevertheless said himself to be God. Constantly it is in the situation of contemporaneity that
the test is to be made. In relation to the highest there is only one time, the present. That one
has learned through dogmatics or catechism, later ages' wisdom, who that man despised by the
authorities of his time was, is only a hindrance to becoming contemporary, thus to coming into the
decisive relation to the way in which the highest has once for all been revealed on Earth. In this
respect no change, no progress is possible, and it is a sign that one has abolished Christianity
when it looks as if it were easier to become a Christian now than in the generation contemporary
with Christ. A historical Christianity is gibberish and un-Christian confusion. The sign of offence
and the object of faith has become a fairy-tale figure. One praises him for that at which one would
most be offended if one were contemporary with him, and one builds one's faith on the outcome and
on what one thinks is history's testimony. One has learned the answer and therefore does not
trouble oneself with the sum itself. --- The only rescue is now to stop the falsification and the
confusion by acknowledging one's distance from the original Christian --- to admit where one
stands, and bow to ideality's demands, even if one cannot satisfy them. ---

Step by step Kierkegaard has thus, in his comparative philosophy of life, risen up from the
``Aesthete's'' wanton play with life to the cutting seriousness in the strict demand of the highly
strained ideality. Our question must now become what is the measure according to which he orders
his stages and calls them lower or higher in relation to one another. We return herewith to the
question that seemed to be going to come to negotiation between two of our greatest poets, Baggesen
and Grundtvig, but that then stopped at the challenge.

Kierkegaard's measure is in and for itself the only one that there can be talk of. It is fetched
from the fundamental concept of spirit, soul, and personality (whatever one will call it) to which
one will constantly be forced to come back: the concept of synthesis.

Where there is consciousness, it is constantly presupposed that a manifold content is comprehended
into unity. From this it follows that a spiritual life-stage stands so much the higher, the greater
and the more heterogeneous a content is given, and the more firmly and inwardly it is worked
together. Or, to use Kierkegaard's expression, the greater the existence-range is, and the sharper
the oppositions are, when yet the unity is maintained, so much the higher a stage have we before
us.\footnote{Compare with this my determination of the concept of psychic energy (\emph{Psykologi},
Fifth ed., pp.~67; 123; 182) and Pierre Janet's corresponding concept \emph{tension psychologique}
(\emph{Les obsessions et la psychasthénie}. Paris 1903, I, pp.~265 ff.; 488 ff.).} The
psychological tension must be so much the greater, the greater the oppositions are that are to be
held together. Such oppositions are, according to Kierkegaard, the moment and the wholeness of
life, time and eternity, reality and the ideal, the world and God. But it is time with its many
moments, in opposition to eternity's constant Now, that conditions the greatest resistance that
the synthesis has to overcome. It was, after all, also the relation of time that conditioned the
constant barrier for our cognition.

But besides opposition and inwardness, Kierkegaard has still another mark of the more perfect
state, a mark that is indeed derived, but that nevertheless gradually forces its way forward to
the chief place in Kierkegaard's valuation. It was not merely the way in which a great fullness of
content is formed into unity that became the decisive thing. The oppositions and contradictions
within the psychic content were emphasized more and more, and the weight was laid on the pain that
conflicting oppositions call forth. Suffering became then gradually the real measure. Now there can
also, in the degree of suffering, especially spiritual pain, that a person can bear, lie a
testimony of his spiritual power. Suffering can be a symptom of great tasks and great fullness of
life, of a courage that in the service of great aims takes up the struggle even against a
proportionally strong resistance. But suffering is --- as little as the feeling of pleasure ---
not in and for itself an absolutely sure mark. As the feeling of pleasure can be a sign that one
has not set sufficiently great demands to oneself (it was therefore that Kant saw it as an
unfortunate symptom when one was all too glad at doing one's duty) --- so suffering can be due to
a convulsive striving after the impossible. Kierkegaard gets it right if it really is so that
there are two midpoints for our life, which each pull to their side, and if the circle in which
our life in experience moves, and where all our positive tasks lie, is not our real element, as
little as dry land is the fish's element. The fish that lies gasping on land has two midpoints,
one in the water, where it is not, one on the land, where it is. Thus, according to Kierkegaard,
man is to be placed: he lives in time, but the yonder is his home. Thereby an inner strife arises
that can never be solved as long as life lasts. Harmonization is impossible. On this point the two
last stages stand in decisive opposition to the four first.

If belief in a yonder existence, in principle different from the present life, is to be the
absolutely decisive principle for our way of taking life, so that our life here is to be
determined by what we think about life yonder, then Kierkegaard gets it right. And he has the
great merit of having called attention to the importance of becoming clear about where we stand.
Where our treasure is, there must also our hearts be, and if we have the treasure both here and
yonder, our heart must be divided. It was this pain of dividedness that Kierkegaard found so
little of when he looked about him in Christendom. One had generally, both in theory and practice,
settled things in such a way that the opposition between the two midpoints was slackened. One had
set a Both--And in place of an Either--Or. And this very Both--And had not become a higher unity.
One lived in fact ``here,'' and the prospect of ``yonder'' played no practically decisive role,
stood only as a last consolation.

Human ethics can learn much from Kierkegaard; he stands in the history of ethics as the alchemist
in the history of chemistry. During his attempt to solve an impossible task, he has cast
significant sidelights over the world of the inner, struggling psychic life. Every ethical
conception that lays weight on the personality's striving to bend life's various forces together
into unity will be able to learn from the Danish Pascal, even though, over against his
ascetic-romantic ideal, it sets the assurance that the highest human thing can be reached by the
deepening, enriching, and concentration of life, as it can be led in the world of experience.

\begin{center}---\end{center}

In his writings and especially in his diaries Kierkegaard has again and again stated that his own
faith and his own life do not reach up to the height with the strict, ideal conception of
Christianity that he had brought to bear. He felt himself here as an unhappy lover and fought
through a humiliating strife over against the thought of appearing as a Christian witness to the
truth over against his age. He became clear that the highest he could reach was self-understanding
and honesty toward that which in his eyes was the ideal. He now thought too that the Church and
its representatives must make such an admission, and his work \emph{Indøvelse i Kristendom} was a
recommendation in this direction --- a recommendation that was not followed.\footnote{There were,
however, clergymen who made the demanded admission. See \emph{Johannes Fibigers Levned} [Johannes
Fibiger's Life], pp.~285--295.} Long, then, he considered the possibility of a public, agitatorial
protest against the existing Church's falsification of the original ideal. What for a while held
him back was partly self-criticism, partly piety toward Bishop Mynster, partly, and not least,
regard for the weak and poor, for whom Christianity in its mild form was a consolation they could
not do without. But he pushed all considerations aside when Martensen, in a sermon after Mynster's
death, presented the deceased bishop as a witness to the truth, a link in the holy chain of the
truth's witnesses that stretches through the ages from the days of the Apostles. He published a
sharp protest against this application of the concept of witness to the truth in the paper
\emph{Fædrelandet} [The Fatherland] on 18 December 1854. And therewith began the great
witness-to-truth strife, which from Kierkegaard's side was carried on in newspaper articles in
\emph{Fædrelandet} from December 1854 to May 1855 and thereafter in small writings (\emph{Øjeblikket}
[The Moment], Nos.~1--9) from May to September 1855.\footnote{All Kierkegaard's contributions to the
strife are now found in \emph{Samlede Værker} XIV.}

This outer witness-to-truth strife --- which, as hinted above, had had as forerunner an inner
witness-to-truth strife in Kierkegaard's own mind --- concerned first the question of what must be
demanded of those who were to be acknowledged as Christianity's witnesses. But it led from there
naturally over to the relation between primitive Christianity (``the New Testament's Christianity'')
and the nineteenth century's Protestant Church. And thereby the strife really came to turn on the
New Testament view of life and on how far our time's Christians embrace this view of life. It is
this last point of view that gives the strife philosophical interest to this day. Only the chief
points in Kierkegaard's contribution can I here dwell on. ---

Only a single thesis have I, says Kierkegaard, while Luther had 95. It runs thus: The New
Testament's Christianity does not exist at all. He who refrains from taking part in the public
worship of God, which claims to be the New Testament's Christianity, has one great guilt less! ---
If official Christianity will at all place itself in relation to what the New Testament calls
Christianity, it has only one thing to do: to admit its distance from it.

Kierkegaard does not fight for Christianity and dares not even call himself a Christian. He fights
for honesty, and what he ventures by his appearance, he ventures for honesty.

``For what, according to the New Testament, is the becoming a Christian, to which the constant and
constant admonition against not being offended, and from which the frightful collisions (to hate
father, mother, wife, child, and so on) that the New Testament breathes? are not both because
Christianity knows very well that to become a Christian is to become, humanly speaking, unhappy for
this life, yet blessedly expecting an eternal blessedness?''

What has happened is this: Christianity has been abolished by flourishing. The religion of
sufferings has, by becoming an established Church, become the religion of life-contentment. The
relation of opposition to the world has disappeared, in that all have become Christians --- in the
sense of the meaningless. If it were right with all these Christians, then in any case the New
Testament, which constantly proceeds from the relation of opposition to the world, could no longer
be a guidance for the Christian, but would in this respect be like an outdated travel book. Since
God cannot have given us such a meaningless book, the explanation can only be sought in a forgery
on the part of men, in a rebellion --- not in defiance, but in hypocrisy. Already the Apostles were
somewhat too eager in spreading Christianity and did not take it so exactly with the individuals.
And since then one has falsified Christianity; one would not give up its consolation, but one would
still less satisfy its demand. One adapted it to one's need. Out of a couple of brighter sentences
in the funeral march that the original Christianity was, one composed for oneself a dashing gallop.
As it has gone with Christianity, so it goes with all great ideas: they appear only in the
beginning in their purity, but are falsified in the course of the historical process and often end
in complete opposition to their beginning.

An essential cause of this falsification of Christianity was that the dogmatic was pushed into the
foreground, and the ethical pushed aside. Christ as the Reconciler stepped into the place of Christ
as the Prototype; one held to his benefactions and forgot the following.

And so one got Sunday-Christians, as one has Sunday-hunters --- battalions of Christians that not
even the Devil will trouble himself to fetch. And men learned on the whole to set lower demands to
themselves. They saw that life became easier the more insignificant they made themselves, while it
becomes more and more difficult for those who see man's task in a striving to be in kinship with
the Deity. Be a triviality, and you shall see, all difficulties disappear, and this world becomes a
glorious world! ---

Shortly after having stated this last sentence (in the last number that appeared of
\emph{Øjeblikket}), Kierkegaard fell sick in the street and was brought to Frederik's Hospital,
where he died on 11 November 1855.

\begin{center}---\end{center}

The attack Kierkegaard directed against the existing Church and Christendom was carried on from a
standpoint similar to the one from which John Henry Newman some years earlier (1836--1843) had
attacked the English Church, and from which he was led over to the Catholic Church. His chief
question was: What have we ventured for Christ? And it was his conviction that ``if the present
distress of which Paul speaks [1 Cor.\ 7, where he advises against marriage] does not indicate the
constant position of the Christian Church, then the New Testament is not written for us, but must
be reworked in order to be applicable.'' This statement recalls quite Kierkegaard's comparison
between the New Testament and an outdated travel handbook; only that according to him men had
already seen to the reworking. It was Newman's assertion that the New Testament very clearly
describes the primitive-Christian life-type, but that we have read the relevant passages so often
that we have lost the ability to take them in their right meaning.\footnote{Compare Richard H.\
Hutton: \emph{Cardinal Newman}. London 1891, pp.~110--119.} --- When Kierkegaard oftener stated
that the falsification was greater in the Protestant Church than in the Catholic, there shows
itself a new agreement with Newman, whom he quite certainly did not know. Thereby it shall not be
said that Kierkegaard, if he had lived, would have gone the same way as Newman. For that the
concept of the individual played too great a role in him, and besides he had expressly impressed
that the meaning in Christianity already burst when one distinguished between two kinds of
Christians, some (monks and nuns) who fulfilled the demand, others (the laity) who compromised.

He has himself emphasized that his struggle against the existing Christendom found a support in the
way in which ``the last formation of freethinkers (Feuerbach and what belongs thereto)'' has
appeared. ``They have really taken on themselves the task of defending Christianity against the
now-living Christians\ldots\ Feuerbach attacks the Christians by showing that their life does not
correspond to Christianity's teaching\ldots\ That he is no doubt a malicious demon may well be, but
in a tactical respect he is a usable figure.''\footnote{\emph{Efterladte Papirer}. 1849, p.~463.}
--- Before Feuerbach, in the eighteenth century Reimarus, in the nineteenth Schopenhauer, had
called attention to the opposition between primitive Christianity's and modern Christendom's ethics.

Whither Kierkegaard would have turned if he had lived, it is idle to discuss. The remarkable thing
was that when he had said the last word that he could say from his standpoint, he had also used up
his life-force.\footnote{His last money too was then spent. ``The day he went out to the hospital,
he sent Giødvad 300 rixdollars, so that he could settle some smaller bills for him, and no doubt
also pay the stay at the hospital. These 300 rixdollars were the last remnant of his fortune; from
them the hospital's bill could just be paid, but scarcely even the funeral costs.'' (Letter from
Brøchner to Chr.\ Molbech, 17 Feb.\ 1856.)}

On the other hand it is of interest for the history of religion to investigate what lies at the
foundation of the opposition, maintained by Kierkegaard and others from such different standpoints,
between the New Testament and the modern Christian view of life --- an opposition that is
undoubtedly present.

Kierkegaard has not brought forward the presupposition on which the New Testament ethics rests, and
without which its negative or indifferent relation to human culture is not understood. This
presupposition is the living expectation of Christ's near return, of the Judgment, and of the end
of all things. ``God's kingdom'' means in the New Testament neither something distant nor something
eternally present, but something that is breaking forth and will soon come fully to light.
Therefore the Apostle warns against founding a family; therefore personal freedom, science, art,
and life of the state got no significance. The gaze was solely directed toward the great and
glorious that was approaching, and that the then generation, according to its master's word, was to
experience. Despite ``the sufferings of the present time,'' one was sure that the victory over
``the world'' was near. Primitive Christianity was no funeral march, but a march of victory. The
enthusiastic expectation was its leading mood.\footnote{Compare, on this whole question, my
\emph{Religionsfilosofi}, IV C.}

Remarkably enough, Kierkegaard explains away this fact. That expectation was, according to him, an
illusion that arose with psychological necessity --- in Christ and the Apostles as in the
Christians --- on account of the violent opposition and tension in which life was to be led. ``It is
to such a degree tormenting to be the true Christian that it was not to be endured if one did not
constantly expect Christ's return as impending now. The torment, the suffering begets a necessary
illusion.'' And from this the consequence can then further be drawn that when one first expects
Christ's return in many centuries, in a distant time, then one is no longer a true Christian --- does
not feel the discord.\footnote{\emph{Efterladte Papirer}. 1849, p.~247 f. --- Characteristic of
Kierkegaard is it that he calls the expectation of the Second Coming ``a subjectively true
rejoinder.'' ``Thus must not only Christ, and next, as is indeed also factual, the Apostles, but
thus must every true Christian speak.'' From the standpoint of the psychology of religion this whole
turn in Kierkegaard is of great interest.} It shall thus not be the promise and the expectation
awakened by the promise that has called forth the tension, but conversely: the tension has called
forth the expectation and --- the promise! --- Here Kierkegaard appears for once as a rationalist.
Against his explanation it conflicts that the expectation of the Messiah's coming belonged to the
whole Judaism of that time, as is seen from a series of Jewish apocalypses. The New Testament's
standpoint is in this respect only different from these contemporary Jewish apocalypses by the
belief that the Messiah had come once, so that the expectation thus aimed at a return. But
Kierkegaard was no lover of apocalyptic literature. Every painting of heavenly conditions, of
blessedness, of the new heaven and new earth that was to come, he regarded as ``aesthetics,'' which
did not belong in religion. It was the way that mattered to him, not the goal. But when now the way
is precisely determined by the goal? Primitive Christianity did not fear ``aesthetics,'' but with
passion and inwardness gave itself over to the religious imagination. Here as at many other points
there shows itself Kierkegaard's difference from men like Grundtvig and Martensen, who precisely
felt the need to paint for themselves the ``goal.'' But he had, after all, his own opinion of how it
went with them as regards the ``way.'' ---

For the historical-religious consideration it is in and for itself natural that Christianity passed
over into becoming a power of culture, in that the original expectation's energy was transposed into
a work for the development of human life, while the ideal conclusion became merely a distant final
end. The strife that then is to be carried on is how the Church has satisfied its task as an
educating and consoling authority. On this point, as we have seen, Sibbern placed himself very
critically in his old days. Kierkegaard did not at all take up the matter from this side.

\begin{center}---\end{center}

Kierkegaard stands as a tragic figure in the history of Danish thought. With his great principle of
truth as the personality's affair, as lying in subjectivity, he dashed against what continued to
stand for him with the power of unconditional authority, the dogmatic formulation of Christianity.
He did not see that the two principles were incompatible. The dogma is, after all, an answer that is
given, and given with the very highest sanction. And one can then really not blame men for not
troubling themselves with carrying out the sum. Kierkegaard's chief proposition gets its
significance rightly only when one with it steers out on a sea where the dogmas --- those lighthouses
that earlier generations have raised --- no longer show the way. Out beyond everything given seeks
the inquiring thought, when the boundary-problems seriously dawn on it. In the work that here is
constantly to be done, Kierkegaard did not come to take part, taken up as he exclusively became with
``reading through, once again, the original text of the individual, human existence-relations, the
old, familiar, and handed down from the fathers.''

But deeply has he dug in the soil of thought and spirit. He has felt the tension and the distance
between ideal and reality, both in the domain of cognition and of life, more inwardly and
passionately than anyone in his century. And however one places oneself toward his personal
standpoint --- one thing he has taught us: only through personal appropriation and working-through
of the thoughts of life (wherever we then fetch them) and through full sincerity and clarity about
where we stand can our spiritual life become sound and true. And thereby he has in truth performed
a philosopher's work.

\begin{center}---\end{center}

% ============================================================
%  15. FREDERIK DREIER   (Da.: "15. FREDERIK DREIER")
% ============================================================
\dfchapter{15. Frederik Dreier}
\label{ch:df-dreier}

\noindent While Kierkegaard in his loneliness felt himself in ever sharper opposition to
the existing Church, while Sibbern immersed himself in his fantasies of the future, and
while H. C. Ørsted worked on \emph{Aanden i Naturen}, there was a young medical student
who, in a series of small works, sketched a programme for free thinking and for the reform
of society --- a programme that at once stood in opposition to those older thinkers'
standpoints and formed a continuation of them.

It was Romanticism's ideas that the older thinkers, each in his way, proceeded from, even
though they sought to bring them into more critical and realistic form. Ørsted and Sibbern
did indeed lay weight on the knowledge of nature, but it was nevertheless the spiritual side
of existence that lay at the foundation in their interpretation of the world-riddle. Howitz
had, to be sure, brought the scientific point of view more into the foreground, but his
influence had not stretched far.\footnote{Kierkegaard divined that natural science would come
to play a more prominent role than hitherto. In the strife between man and God, man will ---
he thinks --- (when the time of speculation is over) entrench himself behind natural science.
He is amused at the new distress theology will then come into; it is Nemesis, because it
would be a science! In his own eyes the matter is not difficult: Christian ethics is and
remains the same, even if one confuses an apple with a pear, and whether the sun goes around
the Earth or the reverse. And God perhaps mixes some gibberish into the world on purpose, in
order to test men and offend the naturalists and the Society for the Dissemination of Natural
Science! (\emph{Efterladte Skrifter} [Posthumous Writings]. 1851--53, pp.~287--289.)}

Now Frederik Dreier (born 1827, died 1853) attempted to lay natural science at the
foundation --- indeed, really to make it the only science. Supported by it, he wishes first
and foremost to work for the reform of consciousness, for the liberation of thought not only
from external authorities, but also from superstition, mysticism, and abstraction, which
still reign even in naturalists and philosophers. When this has taken place, the rest will
follow of itself; if thought is freed from theology, labour will be freed from capital, the
people from the state.

There goes a great enthusiasm through the young author's writings.\footnote{The chief work in
a philosophical respect is \emph{Aandetroen og den frie Tænkning} [The Belief in Spirit and
Free Thinking]. 1852. --- Previously Dreier had anonymously published \emph{Folkenes Fremtid.
Af en Fritænker} [The Future of the Peoples. By a Freethinker] (1848) and \emph{Fremtidens
Folkeopdragelse. Af en Socialist} [The Popular Education of the Future. By a Socialist]
(1848), as well as a sharp critique of Goldschmidt and a contribution to the Clara Raphael
feud. --- A detailed biography and presentation of his theories is found in C. E. Jensen and
F. J. Borgbjerg: \emph{Socialdemokratiets Aarhundrede} [The Century of Social Democracy],
Second Volume, pp.~32--135.} Despite his theoretical materialism, he is a practical idealist.
He was a predominantly intellectual nature. The other joys that life can bestow had, in his
eyes, great progress to make before they could reach the joy of scientific thinking and
occupation. And scientific progress leads necessarily to the social. With the intellectual
joy there connects itself therefore naturally the zeal for the reform of society. Resignation
has no place in Dreier's view of life; it yields its place to a vigorous working for reform
in all directions. His highest wish is to work for harmony and beauty in life. The men of
free science and the workers are in union to carry on the struggle against all old routine.
What religion was for earlier times, the passion for reform can be in our time. All science,
all capable and noble practical activity, supports it and stands in its service. It is also
reported of him that from childhood he nourished great fellow-feeling with the suffering and
the unfortunate.

H. C. Ørsted and Sibbern belong in his eyes to the mystical and dualistic tendency in
thinking; he finds in them, however, forerunners of the assumption of the origin of species
through natural development, so that man too has arisen as a fruit of a long course of
development. Ørsted's philosophy of nature he criticizes sharply; a sketch for a special work
on this subject is found among his papers at the Royal Library. Kierkegaard he takes into his
account for the sharp opposition between theology and science, our time's sharpest opposition.
There is, precisely by the extreme opposition in reality, a certain analogy between
Kierkegaard's and Dreier's thinking. Both wish to get away from ``concept-mysticism'' to
existence. But while Kierkegaard mostly or exclusively has the existing subjects in view
(``Subjectivity is the truth''), Dreier lays the chief weight on the existing objects
(``Everything outside us is matter, and everything inside us is matter''). In both respects
there could then also, after Romanticism's overflowing developments of thought, be need for a
sober reduction.

His chief influence Dreier has gotten from his time's natural science, especially, it seems,
from Du Bois-Reymond and Carl Vogt. He admires besides the eighteenth century's
Encyclopaedists, especially Holbach and Condorcet, in whom he sees his predecessors, and he
finds on many points the way prepared for him by the radical, left wing of Hegel's disciples
(Feuerbach, Stirner), as well as by the French socialist authors. Of English authors he names
Bentham, Owen, and especially Stuart Mill, whom he admires as a logician, as ``a sharp,
consistent thinker who does not respect concept-mysticism,'' while he criticizes his political
economy.

It is, as one sees, currents that had not hitherto made themselves felt here in the country
that the young reformer has come into contact with. We shall now see how his chief thoughts
have shaped themselves during his short career. It came only to germs, but had his life not
been so early cut off, they would surely have unfolded themselves in luxuriant growth.

\begin{center}---\end{center}

It is the philosophical side of Dreier's view that I here especially hold to, although it is
his social-reformatory ideas that in the most recent time have especially reawakened attention
to him.

Dreier felt himself in a strong opposition to Romanticism's philosophy, in which he could see
only a last attempt to defend ``the belief in spirit.'' He divides all philosophy into
apologetic and critical philosophy. The apologetic schools presuppose without further ado the
inherited faith's fundamental concepts and have only the aim of adapting them together. The
critical schools aim at abolishing the oppositions (God and the world, the good and the evil,
and so on) that theology presupposes and the theologizing philosophers transform. When this
critical work is done, philosophy is no longer a peculiar science; it melts together with the
whole rest of science as its most general concept-distinguishing and concept-comprehending
method.

By this last hint Dreier's conception goes in a similar direction as Kierkegaard's idea of a
system of thought that was to comprise the fundamental concepts our cognition works with. And
it is a task that philosophy at all times must set itself, the most central of all its tasks.
What it can accomplish toward the critique of current and handed-down representations depends
in the end on the understanding of the very nature and mode of operation of our thinking, our
faculty of representation. It is, as lies in Dreier's words, the last differences and the last
unities that philosophy is to find, so far as it can take place at the stage of cognition's
development where the philosopher undertakes his work.

Dreier conceives only the theological representations and their after-effects as the object of
this critique. He has not seen that natural science's own fundamental concepts and
presuppositions must be objects of the same kind of critique, and that this critique had long
since been under way. Thereby there comes something naïve into his line of thought, naturally
called forth by his enthusiasm for his special science and its mission.

In his conception of cognition, Dreier builds on the presupposition that sensation is the
foundation of all thinking. But since by sensation he understands only the reception of
impressions --- impressions that can then leave behind residues (sense-images) that can be
repeated even when no impressions are present --- he nevertheless presupposes, as different
from this reception and its after-effects, the process by which the sense-images are connected
and compared. It looks as if the individual sense-images were independent of every connection
and comparison. But then ``thinking'' must have another foundation than ``sensation'' (in the
sense in which Dreier takes this word). In a similar way he explains motion as the result of
the interaction of things --- as if things were first there, then came into interaction, and
motion in turn were the consequence of this interaction --- and as if things' (the material
things') interaction were due to something other than motion.

Dreier complains of the concept-mysticism that arises through the inclination to make the
abstract concepts into similar, self-appearing particulars like the sense-representations
themselves. But he is himself a mystic as regards the simple sense-representations, ascribes
to them an independence they do not have. And he who hates all dualism proceeds from a
dualism between the individual representations on the one side, their connection and
comparison on the other side.

In the domain of the conception of nature, it is especially the opposition between matter and
force that he will do away with. It is in his eyes the last remnant of the opposition between
the world and God, the belief in spirit's last off-shoots. One must not regard matter and
force as independent beings, and he praises the \emph{Système de la nature} for having
cleared up well here. According to his own opinion, the abolition of that dualism must take
place by making matter into everything: ``Everything that exists outside us is matter;
everything that exists inside us, in our body or in our consciousness, is also matter, in
certain definite connections, forms, and activities.'' --- But is ``matter'' not also an
abstract concept? Does it exist apart from the special forms and activities that are what we
have given, and that our thinking is to investigate?

As it goes with the opposition between matter and force, so it is to go with the opposition
between body and soul, which has the same theological origin. The concept of soul is due to
one of the beginning thinking's usual personifications. That the soul is not material means
that it is not real; for we know nothing that is not bodily. What one meant by the word soul
was what we now call brain-activity. There is indeed a difference between consciousness and
brain, but it is the same difference as between the stomach (with the intestines) and
digestion. Carl Vogt is not right in saying that thought relates to the brain as the bile to
the liver; for thought is an activity.\footnote{Carl Vogt uses, for the rest, a more indelicate
comparison. See on his statement and on the materialist theory on the whole my \emph{Psykologi}
II, 8 c.} --- Here, then, the same difference comes forth again as before, between matter and
its activity. And besides, an organ's activity consists, after all, only in motion. Whence,
then, comes the difference between a thought and a motion? ---

The enthusiastic young medical man has no time to answer such questions. He wishes to have the
life of thought as simple and clear as possible, in order to be able to pass over to his
social plans of reform. He has not even time rightly to think over his own enthusiasm. What
one calls feeling is in his eyes something subordinate, often only unclear cognition; if one
analyses it more closely, it consists partly of thoughts, representations, and fantasy-images,
partly of organ-sensations that stem from changes in nutrition, secretion, circulation of the
blood, and so on.\footnote{As one sees, a hint in the direction of the theory later set up by
William James and Carl Lange.} He would, however, perhaps on closer consideration have found
this theory somewhat too elementary.

That he cannot ascribe any significance to religious feeling follows from what has been said
in the foregoing. It is next his conviction that it will go with poetry as it goes with
religion. Its time is over. Where poetry still reigns, there is a hole in science, just as
where religion reigns. Originally poetry corresponded quite to the peoples' intellectual
development, in that it gave such a description and explanation of reality as one could have
at a low stage of cognition. But now one must, after all, study old mythologies in order to
understand poetry! --- Imagination is indispensable for science, but it is to be serving. In
opposition to poetry, science describes reality as exactly and with as great connection as
possible; poetry uses images and representations that belong to a long-vanished time. The
significance of pictorial art is to be illustration, landscape painting for example for
geography. A similar significance the psychological and social novel will be able to have in
its domains.\footnote{In a manuscript with the heading \emph{Poesien} [Poetry], which is found
among Dreier's papers at the Royal Library, this line of thought occurs.}

It is the same with the national feeling. The Danish people is the sum of all Danes. It does
not disappear because each individual Dane speaks German, any more than the individual
disappears because he gets another language. To what end, then, all the violent national
disputes? Especially since the adoption of the foreign language will be only provisional,
since a universally human language will be the necessary consequence of the universally human
science. And already now there is not that community within a people that there is between
individual classes in the people and the corresponding classes in other peoples.

In place of the feelings whose time is over, a new order of society will develop new feelings.
Now it is competition that dominates everything. But it is a matter of fostering love of labour
itself and of making possible the recognition that men's interests are essentially the same,
and that the happiness of the one is bound up with that of the other. This will be possible by
a socialist order of society, where capital is common property. Then the essential thing in
Christianity will be able to be satisfied: the great protest against all injustice, the
commandment of love, and the demand not to do to others what one does not wish they should do
to us. The first Christians' communism was a presupposition of their brotherhood. It was, now,
a less fortunate form of organization; but the modern social reform will have a similar effect.

The really value-producing classes in society are the workers, the artists, and the men of
science. It is their affair to combat the old illusions and bring forth a new order of things.
The reforms will not proceed from the state, but from the individuals, who must join together
and then gradually, through the extension of the municipal and political franchise, attain
complete self-government. It is then a matter of hindering private capital-formation, which can
take place by the exclusion of middleman-trade, in that storehouses under public administration
are set up, where producers can deliver their goods, and the consumers receive them. A mutual
credit-association orders the exchange and fixes wages and prices of goods.\footnote{See
Dreier's periodical \emph{Samfundets Reform} [The Reform of Society]. 1853. (Especially the
article ``Kapital og Arbejde'' [Capital and Labour].) --- On the great significance that
cooperative societies later acquired, compare my \emph{Etik}, 3rd ed., p.~382 f., where their
significance for the solution of the social question is shown.} --- Whatever one thinks of the
feasibility of such an order, the idea nevertheless marks a progress in relation to Sibbern's
naïve reference to the fact that everyone can take what he wills, since no one has reason to
take more than he needs. Dreier also deviates from Sibbern in that he lets the economic order
be the condition for the brotherly disposition, while this in Sibbern is the presupposition of
the former (as was surely also the case with the primitive-Christian communism, although Dreier
seems to think the opposite). Dreier recalls here Karl Marx, whose first writings he has known,
and according to whom all order of society and all spiritual development is determined by the
economic production and its ordering at any given time. There is, however, here a constant
interaction between mode of thought and order of society; they constantly influence each other,
and it was a speculative after-effect in Marx when he wished to construct all cultural
development from the economic point of view.\footnote{Compare, on Marx, my work \emph{Den nyere
Filosofis Historie} [The History of Modern Philosophy], Second edition, II, p.~282, and my
\emph{Etik}, Third edition, pp.~387--394.}

Dreier's pedagogical and linguistic investigations I do not enter upon more closely. Even
without taking account of them, one will have gotten the impression of how many ideas occupied
the young man's interest. What fate they would have gotten in his mind and outside him, if his
career had not been cut off, is not easy to say. Now he came merely to stand as a herald that
Romanticism's time was over, and that a sober stand toward reality was to be taken, both in
practice and in theory.

\begin{center}---\end{center}

% ============================================================
%  16. RASMUS NIELSEN   (Da.: "16. RASMUS NIELSEN")
% ============================================================
\dfchapter{16. Rasmus Nielsen}
\label{ch:df-nielsen}

\noindent Dreier's appearance was a little episode that was not much noticed, though it might
well have given occasion for reflection in the philosophizing. There came in his small
writings, as in Kierkegaard, the demand for a philosophy that carried through an
investigation of science's own fundamental concepts and fundamental presuppositions. But
speculation still exercised its power, and the investigation that was demanded both from a
religious and a scientific point of view was for the time being postponed.

In order to follow philosophy's continued fate in Denmark, we must now go back again to 1840,
the point of time when both Martensen and Kierkegaard began their authorial activity, the
Hegelian period.

Like Martensen, Rasmus Nielsen (born 1809, died 1884)\footnote{Biography and presentation in
P. A. Rosenberg: \emph{Rasmus Nielsen. Nordens Filosof} [Rasmus Nielsen: The Philosopher of
the North]. 1903. --- I believe that my old teacher, whose influence I always remember with
gratitude, would have been better served by a more critical presentation than the panegyric
Mr.\ Rosenberg has here given. I have in any case left Rasmus Nielsen's school with a
stronger feeling of the sting of the problems than is traced in the cited book, where almost
a problem is ``solved'' on every second page.} had set himself the task of applying the
Hegelian philosophy for the benefit of theology, yet in such a way that the said philosophy
at the same time underwent a necessary change. Nielsen took his degree of licentiate in
theology in 1840 with a dissertation on \emph{den spekulative Metodes Anvendelse paa den
hellige Historie} [The Application of the Speculative Method to Sacred History]. The
objection that is here directed against Hegel is roughly the same as Martensen had put
forward some years before. It is the difference between thought and reality that is
impressed: not man's thought, only the divine thought, which is at the same time creative, is
expressed in reality.\footnote{\emph{De speculativa historiæ sacræ tractandæ methodo
commentatio}. 1840. --- The fundamental thought is the same as in Martensen's dissertation
some years before: faith as the foundation for speculation, in the confidence that in faith
there is contained a hidden dialectic (p.~43 f).} Martensen had expressed this difference as
an opposition between knowledge and power, and this opposition came to play a great role in
Nielsen's thinking through all its various periods. The concept of power Hegel had introduced
in his Logic in the section in which he makes the transition from the concept of substance as
mere inherence or ground to the concept of causality, and Martensen had maintained that one
could not make this transition without the presupposition of the creation-dogma, since only
the divine thought can be cause of a reality different from itself.

The work that was to be done was now divided between the two men, who nevertheless for the
time being worked much together, Martensen on his Dogmatics and Nielsen on a philosophical
system that was to let Hegel undergo the same fate that Hegel himself had prepared for so many
earlier philosophers, namely to be made into an element in a ``higher unity.'' For Nielsen's
lively imagination, a faculty that was strongly developed in him, the new system already
shaped itself in its fundamental thought, and he published an invitation to subscribe to a
speculative work, which was, however, not completed, as well as to a moral philosophy that did
not appear. Kierkegaard mocked at the many announcements of ``the System'' that was to solve
all riddles, even those that Hegel had not been able to master.

Nielsen's most important work from this period is his \emph{Propædeutisk Logik} [Propaedeutic
Logic] (1845). He develops the already hinted-at critique of Hegel, whom he reproaches with a
formalism that constantly lets the relation between thought and reality stand uncleared. ``The
formal comprehending lacks omnipotence!'' Hegel's dialectic, the constant mutual illumination
to which the concepts are subjected, is acknowledged and admired. But by the unclear
identification of concept and reality Hegel becomes a logical mystic. --- The
concept-mysticism that Dreier complained so much of, he may perhaps have become aware of
through Nielsen's critique of Hegel. ---

It was, according to this line of thought, not difficult to make the transition from philosophy
to theology; it was, after all, really already made. The concept of God was already introduced
as last or first cause (omnipotence), and it does at the same time service as absolute subject,
whereby the difficulty is to be solved that there should exist objects that were not there for
any subject, thus were not objectified. The two concepts power and objectification go into one.
The confusion of a logical (epistemological) and a metaphysical line of thought, on which
German speculation rests, thus constantly makes itself felt. It was this that Kierkegaard
protested against by his distinction between a logical system and a system of existence, and
that, for the rest, Fries had already criticized in Fichte and Schelling.

The relation between philosophy and theology, as it presented itself for Nielsen in this first
period, is expressed in the following propositions: ``The scientific medium in which
philosophy's theoretical and religion's practical idea mutually explain and ethicize each
other is theology. For theology seeks not merely a life-cognition for truth's sake, but also a
truth-cognition for life's sake, and has therefore with right been regarded as the essential
centre of all sciences, as the proper and true central science.'' Theology is here thus both
starting point and conclusion.

It was the harmony between theology and philosophy that reigned at our university in the
[eighteen-]forties during Martensen's and Nielsen's collaboration, disturbed only by
Kierkegaard's irony and criticism and by some radical Hegelians' (Brøchner's and Beck's)
objections. Sibbern at the same time made in silence the transition to his later standpoint,
which he stated in the programmes for 1846 and the following years.

Then there came a sudden end to the collaboration and the harmony.

\begin{center}---\end{center}

It became clear, namely, to Nielsen that in Kierkegaard's writings it had been so energetically
pointed out, the personal, existential side of Christianity, that the possibility of
transposing the Christian dogmas into the form of speculation must fall away. ``The System''
was burst, a speculative theology became an impossibility. He entered into connection with
Kierkegaard, became his eager adherent. And as one had earlier in Copenhagen seen Sibbern
walking with Poul Møller, and Poul Møller with Kierkegaard, so one could now (every Thursday)
see Kierkegaard walking with Rasmus Nielsen. All these walks have had their significance for
Danish intellectual life.

When Martensen's Dogmatics appeared, Nielsen, supported by Kierkegaard's \emph{Uvidenskabelig
Efterskrift}, directed a strong attack on it, and a whole theological-philosophical feud spun
itself out. Nielsen published besides, in the following years, several popular writings, of
which especially must be named \emph{Evangelietroen og den moderne Bevidsthed} [The Gospel
Faith and the Modern Consciousness] (1849).

For two years (1848--1850) the personal relation subsisted between Kierkegaard and the disciple
he had so unexpectedly gotten, and who had broken with friends and colleagues in order to
attach himself to the lonely thinker. It was not, as one has said, out of wounded vanity that
Kierkegaard again drew back, in that he was supposed to have taken Nielsen ill that in his new
writings he had not sufficiently given his new master the credit for his views. To be sure,
Kierkegaard found that Nielsen adorned himself somewhat with borrowed feathers. But his
conception of the relation he stated in the following words (at the close of the diaries from
1849): ``R. Nielsen's merit is that he has become attentive (not a little, when one is an older
man). His fault is that he then nevertheless makes that which is no science into science\ldots\
And next it is his fault that, instead of limiting himself to simply confessing something about
himself, he wants at once to use what he has learned to attack another.'' --- What Kierkegaard
missed in Nielsen was the inwardness in the personal life that first and foremost works through
what is received and seeks to test it practically. And that which here was the received was
precisely a practical conception, a personal attitude. But as formerly Hegel's influence, so
now Kierkegaard's influence spurred first and foremost the gifted man's imagination and zeal
for speculation. He was a lively and awakening speaker, who knew how to present the oppositions
of thought and life in great images and strong strokes.

In his Hegelian period he was bound constantly to have to find the ``higher unity,'' whether he
could now find it within philosophy or took theology to his aid. But now, when not only
philosophy and theology had become irreconcilable oppositions, but on the whole it looked bad
for all ``higher unity,'' the ability and the desire to work by means of oppositions must come
to play so much the greater a role. And he exercised a not slight effect through the antitheses
he often presented with striking force and with piquant sharpness.

\begin{center}---\end{center}

But the new standpoint nevertheless also became the introduction to a great period of work. The
acknowledgment that the relation between ``life-cognition'' and ``truth-cognition'' was not so
simple as he had thought in his first period must lead him to subject the relation between
thought and reality on the whole to a new testing. There came now a series of years of work in
which he immersed himself in the study of mathematics and natural science in order to get a
solid foundation for the illumination of the said relation. In this period he surely exercised
his greatest effect, because he was himself a seeker.

Yet he did not take up the question from the ground up. He did not become an epistemologist,
but held fast, despite everything, to Hegel's fundamental thought, though he now gave it a new
modification. There reigns indeed, in the various introductory works that he published under
the name of Propaedeutics in the [eighteen-]fifties and the beginning of the sixties, still a
certain critical caution. As regards form, the Propaedeutics of 1857 is his most finished work;
here reigns the limitation that he was otherwise no master of. But the stronger stamp of
independent reality that the individual facts were now, according to the programme, to have was
nevertheless not attained. Either he takes up the material from the other sciences quite
unworked, or else it is used merely as examples of the Hegelian dialectic. Mathematical
functions, physical and chemical processes, organic forms of development, everything appeared as
dialectical movements --- only that, as said, the higher unities became more and more invisible.
The connection in which natural phenomena constantly stand with one another, so that for example
two substances can by combination produce a third, the dialectician interprets without further
ado as a logical connection --- to the advantage of the demonstration of ``the objective
concept's'' validity. And throughout, that which is essentially method or usable analogy, like
the infinitesimal calculus and the atomic theory, is made into absolute points of view. The
unity in existence was under various forms made intuitable by the analogy with human concepts'
relation to their content-elements (moments), without its here, however, any more than in Hegel,
being admitted that it was (only) analogy that reigned. The very justification and the
definitive significance of such an analogy was not discussed.

After the conclusion of his preliminary studies, Nielsen began the publication of a great work
whose name was \emph{Grundideernes Logik} [The Logic of the Fundamental Ideas]. Only two parts
appeared (1864--66) --- and they made up only two-thirds of the intended work's first ninth!
Again here it showed itself that Nielsen's imagination hurried ahead of his thought. He saw in
the spirit that, in order to describe the mutual relation between the subjective and the
objective, or, as he constantly calls it, knowledge and power, there was needed an infinite
analysis, since every subject presupposes an object, and every object again a subject. When one
will not conclude in a speculative and theological way, it is a single combat that can never
come to an end --- and it was therefore no accident that the work was not concluded.
Scientifically seen, one must here be content with demonstrating the constant relation between
subject and object as a law: every quality that one ascribes to the subject can itself be made
into an object and must have its ground in certain objective relations; and every object, or
every objective relation, is conceived in a way determined by the nature of the conceiving
subject --- and thus one can keep on, by virtue of the law for the relation between subject and
object.

Nielsen did not now conceive the matter in this way. His line of thought was that since every
object must be objectified, i.e. presupposes a subject, and since the human subject cannot
conceive (objectify) everything, there must, if the reality of objects is to be maintained, be
an absolute (``ontological'') subject, for which the absolute reality exists. He overlooks that
the game must begin again here; even a God is bound to the law of the relation between subject
and object. The attempt, through \emph{Grundideernes Logik}, to found an abstract theism has
therefore not succeeded. We cannot come further than to determine and describe a subject in
relation to which certain phenomena (objects) hold, just as the astronomer must determine a
point (on the Earth, on the sun, or wherever) from which the positions and motions of the
heavenly bodies can be described as if they were absolute.

There are found many fine analyses in Nielsen's unfinished work. And what has special interest
in the history of Danish philosophy is the strong weight that is laid on the concrete and
individual, which does not go up into any concept. Nielsen here even uses the expression ``the
antilogical'' to designate this thing insoluble in universal concepts. It is for him not, as
for the strict Hegelians, an indifferent residue. It is for him the striking expression for a
power that we cannot get around. Yet by the expression ``the antilogical'' the opposition to
cognition is not aptly characterized. It is the inexhaustible, the constantly incommensurable
in every concrete phenomenon that we here stand over against; but ``antilogical'' is not the
right expression, since there is constantly, through new observations and new comparisons, the
possibility of new determinations. ---

As the tireless worker Nielsen was, he constantly sought to find new forms for his ideas.
Besides a whole series of popular writings, must be named his \emph{Religionsfilosofi}
[Philosophy of Religion] (1869), his \emph{Naturfilosofi} [Philosophy of Nature] (1873), and
finally \emph{Filosofiske Grundproblemer} [Fundamental Problems of Philosophy] (in the
university's Festschrift 1879); this last is the clearest and most condensed presentation
Nielsen has given of his conception, as it shaped itself in his later years. He here ascribes
greater independence to the problem of cognition than he has done in his earlier works.

I will here merely dwell on the way in which the religious problem presented itself in Nielsen's
philosophy after the alliance with Kierkegaard.

\begin{center}---\end{center}

When I in the year 1861 came to the university, Nielsen stood in his full force, full of zeal
and desire for activity, and the course I heard with him in my first student year got great
significance for me. He lectured essentially on Kierkegaard, in that he went through the
``Stages'' according to a little selection of chief passages in Kierkegaard's writings. Here
again, indeed, the plan was so great that it was not carried out, but with great eloquence he
nevertheless led his young hearers on the way to an understanding of the great figure who had
recently left the field of battle, and whose appearance was still felt in its after-effects. To
this presentation there attached itself then a little outline of logic. I got here a living
impression both of the most personal and of the most abstract that thinking must come to stand
over against, the two end-points whose mutual relation it is a matter of determining. Nielsen
was then still a seeker, and precisely therefore he pointed out to great horizons. Since I now
had my distress with theology, in which I could rightly see neither the personal nor the
logical come fully into their right, the dualistic standpoint concerning faith and knowledge
that Nielsen in that period maintained was to me a reassurance. I took part some years later in
the dispute about faith and knowledge that Nielsen's statements awakened, with a little (highly
immature) writing. My development led me away from this standpoint,\footnote{I permit myself to
refer on this to the little essay I have called \emph{Forord og Efterskrift til min
Religionsfilosofi} [Preface and Postscript to My Philosophy of Religion], which is found in
\emph{Mindre Arbejder} [Minor Works], Second Series. (Especially pp.~47--51.)} but what I owe
Nielsen is the interest he awakened for the problem, and that has held itself in me, even if my
interest in it has become predominantly psychological.

Psychology was precisely something one could not learn with Nielsen. Either he spoke in
rhetorical and edifying turns or in the language of abstract logic. I seek therein the ground
that he, without at first noticing it himself, gradually glided away from the strict
Kierkegaardian conception of the problem. As Brøchner has called attention to, there was from
the first a peculiar difference between Kierkegaard's and Nielsen's grounding of their common
standpoint. For Kierkegaard, as we have seen, the main thing was that the existing individual
lives in the middle of time, without any sure decision of the most burning questions, and among
oppositions that cannot be reconciled; for Nielsen the problem had its origin from the relation
of opposition between the individual and the universal, which, after all, admits of so many
nuances and transitional forms. This difference is again connected with the fact that
Kierkegaard stood nearer the critical philosophy, and indeed also demanded its resumption, while
Nielsen's speculative dialectic could indeed for a while be damped, but soon after unfolded
itself again, so far as, after the stoppage that had taken place, it was possible.

The Kierkegaardian paradox, the boundary for thinking that the dogma's absurd content sets, when
it is nevertheless to be acknowledged, was gradually mildened for Nielsen. In his popular
writings, in his Philosophy of Religion, and in his essay on the fundamental problems, he sets
``the mystery'' in all indefiniteness in place of the paradox, and the improbable in place of
the impossible. And a still more significant turn was made by the way he placed himself toward
the miracle. The chief weight --- such is his line of thought --- is, religiously seen, not to
be laid on all the particulars in the outer event; the decisive thing is its spiritual
significance; the reality that is here in question is ``spirit-reality,'' ``reality for the
spirit.'' In his last writing he explains this expression thus: ``Spirit-reality means that here
there is really a content of faith, an unambiguous revelation of the divine will; reality for
the spirit means that the outer must not be held fast in so sensuously immediate an externality
that there can here be talk of a breach of the laws of nature.'' Here Nielsen stands right by the
concept of myth. When the miracle does not need to appear ``so externally'' as real events
otherwise are wont to, and when its whole significance depends on what it means spiritually seen,
then we have precisely what one calls a religious myth. The miracle is a symbol, a sign --- ``an
outer sign of an essential inner.'' The weight lies then, for example, not on what really
happened when Jesus ``ascended to heaven,'' but on what the Apostles experienced in their
interior when this event took place. It became a similar explanation of the story of the
Ascension as we have already met with in Martensen --- and which, for the rest, Nielsen had
already hinted at in his dissertation. There is, it was said there, always an element of illusion
in the phenomena of revelation (\emph{in phænomenis revelationis semper momentum illusionis
insit, necesse est}), because the revelation must be adapted to the human conception. It then
makes no difference whether the Ascension takes place in the outer order of things (\emph{in
externo rerum ordine}), or whether it was only for the Apostles' thought (\emph{apostolis
meditantibus}) that the phenomenon was produced.\footnote{\emph{De speculativo methodo}, p.~148.}
It is a hair-fine difference that is here between the speculative theology and the mythical
conception. It is thus so that God has always been a Copernican, or rather a Brunonian, but has
had to adapt himself to the human conception. Since now the philosopher too is a Brunonian, he
must also adapt himself; what for him stands as a clear concept is for the positive religious
consciousness a vision, an image. This was, after all, precisely the fundamental thought in
Hegel's philosophy of religion; to philosophy's highest concepts there correspond in the
believers concrete intuitions or images that express the same spiritual content in the form of
representation or imagination.

When Nielsen had thus involuntarily returned to his youth's conception, it had its ground in the
fact that he came to stand over against a question that Kierkegaard had always pushed aside, and
that one also in general goes around, namely the interpretation that dogmas and miracles must
undergo when they are really to have personal and spiritual significance. It is the problem of
the religious picture-language that had already occupied Grundtvig and Martensen. One may well
issue a prohibition against ``playing with the picture-language,'' but the question does not,
after all, let itself be dismissed, what real value the picture-language has for the spiritual
life, and why precisely this definite picture-language is used, and not another.\footnote{Compare
more closely, on the relation between thought and image, especially with regard to the religious
problem, my \emph{Religionsfilosofi} II C.}

It was no accident that Nielsen in his later years felt much sympathy with the Grundtvigian
tendency, which precisely laid great weight on the picture-language and its interpretation, but
at the same time laid little weight on scientific dogmatics. Here he could hold fast his doctrine
of faith and knowledge as absolutely heterogeneous principles and his protest against scientific
theology, while yet a free spiritual interpretation had to be permitted. The sympathy was,
however, not quite mutual, since one on the Grundtvigian side was somewhat doubtful of the
``spirit-reality,'' which not without right recalled mythical symbolism, as soon as one was to
think anything definite by it. Nielsen ``played with the picture-language'' in a suspicious way.

But the real motive for Nielsen's approach to Grundtvigianism lay, however, surely in the
ethical domain. For Kierkegaard, as we have seen, Christian ethics was the main thing. When he
declared that the New Testament's Christianity did not exist, he thought first and foremost of
the New Testament ethics. Over natural science's advance he consoled himself, after all, with
the thought that Christian ethics must remain the same, whether one in natural science came to
the one or the other result. It was the opposition between the religious and the human that he
wished to maintain. Precisely this opposition Nielsen sought to get beyond, like
Grundtvigianism. And he separates himself definitively from Kierkegaard by, in his book \emph{Om
Aandsdannelse} [On Spiritual Formation] (1877), distinguishing between life and form of
existence. ``Life and existence are not synonymous words. The same life can express itself in
different existences\ldots\ The eternal life in Christianity's essence does not stop at a single
age's stereotyped scheme of existence\ldots\ The apostle-existence ceased when the
breakthrough-period was concluded.'' The religious life can, Nielsen thinks, be the same as in
the Apostles' days, in that one can live in the same faith and hope; but one cannot throw away
everything that the long work of culture has procured, and in one's conduct of life without
further ado imitate the way of taking life that was peculiar to the great epoch when Christianity
stepped into the world. Kierkegaard was right against his ecclesiastical opponents, because these
themselves solemnly acknowledged that the New Testament scheme of existence still held. He was
also right that the ideals are to be proclaimed. But another question is how they are to be
realized and shape themselves into existence. He did not understand present-day existence and its
spiritual conditions. It is a matter of living along in modern culture and letting the great
Christian ideals penetrate it. --- This was now precisely the conception against which Kierkegaard
in ``The Moments'' had directed the most violent attack, emptied his whole quiver of mockery and
indignation. And it was the conception that Nielsen's earlier comrade and later enemy, the year
after, in the second part of his Ethics, developed far more richly than Nielsen was able to. More
and more it must then become clear that there was need for an investigation of ethics'
presuppositions and results. There was, after all, the possibility that present-day existence,
and human existence on the whole, out of its own forces and life-conditions, could win a
foundation for an ethical conduct of life. This possibility Kierkegaard had pointed to in his
stages (the second, third, and fourth), but let it fall in his zeal to reach the ``scheme of
existence'' that it was his mission to set up in its original clarity, over against all sorts of
adaptations that would not be, that they were adaptations. That Nielsen let this possibility lie
was a chief ground that his successors had to strike out on quite other ways than his.

\begin{center}---\end{center}

% ============================================================
%  17. HANS BRØCHNER   (Da.: "17. HANS BRØCHNER")
% ============================================================
\dfchapter{17. Hans Brøchner}
\label{ch:df-brochner}

\noindent According to the plan for this work, it is to conclude with the philosophers who
have been my teachers. So far I believe it will be possible to carry through a purely
historical conception; my own personal standpoint naturally makes itself felt in the whole
presentation, both in that which is nearer in time and in that which is more distant, but need
not step forth more in the former than in the latter.

When I came to the university, Sibbern, despite his stiffening in strange outward forms,
nevertheless stood, by what I knew about him, as so venerable and interesting that I could not
think otherwise than that I must hear him. It was, however, only far later that his worth and
significance dawned on me, so that from an essential side I wished and strove to continue his
work. For the time being it was Rasmus Nielsen whose enthusiastic reference to Kierkegaard and
whose awakening eloquence got the greatest influence on me. But I nevertheless in my student
years constantly heard Brøchner's lectures on the history of philosophy, and also took part in
exercises with him. Later he was to become the one of my teachers with whom I came into the
closest personal relation. Of no one I have known does it hold as of him, that he had and
sought his strength in his limitation. He was a man of will from first to last, however much he
moved in the abstract world of thought (and even made this more abstract than it needed to be).
In a letter to his friend Chr. Molbech (24 Dec.\ 1862) he writes: ``I honestly endeavour to
come to know my strength and my limitation, in order to be able to set my goal in such a way
that I can approach it; for one is not to \emph{reach} a self-set goal. And should I be
mistaken about the goal and the ways, one thing I nevertheless keep, in which one cannot be
mistaken: the having willed, and the ideality that lies therein.'' These words appeared on the
occasion of a contemplated philosophical work, but they are characteristic of him both as a
human being and as a thinker. He felt his limitation as the means that stood at his disposal
when he was to give his contribution to the spiritual development and thereby develop his own
personality. There was a mood of resignation over him, nourished by sickness and sorrow, as
well as by the consciousness of how little he reached of what stood before his thought, and how
imperfect an expression he had been able to give to what he had reached. But his powerful will
hindered the resignation from getting the timbre of bitterness or coldness. As the knight of
thought that he was, he did not lay down his arms as long as it was still possible for him to
fight.\footnote{In Brøchner's letters to Chr. Molbech, and in the way in which his friend in his
letters comes to depict him, one has a beautiful picture of him, better than any biography could
give. In the introduction to the edition of these letters I have given a short characterization
of my teacher and friend. (\emph{Hans Brøchner og Chr. K. F. Molbech. En Brevvexling} [Hans
Brøchner and Chr. K. F. Molbech: A Correspondence]. 1902.)}

Brøchner (born 1820, died 1875) was in his youth, like Nielsen and Kierkegaard, strongly
influenced by Hegel, and this influence never quite disappeared, not even in his language, which
preserved its abstract character, without, as so often in Hegel, being interwoven with striking
images and mighty pathos. He left Hegel's school in another direction than the other Danes (Poul
Møller, Martensen, Kierkegaard, and Nielsen) who had belonged to the school and then made their
own ways. He granted Strauss and Feuerbach right in their critique of Hegel's attempt to solve
the problem of faith and knowledge. Strauss's historical-critical presentation of the Christian
doctrine of faith he translated into Danish. The way by which he sought further was for an
essential part determined by Kierkegaard's personal influence. It is of interest both for the
understanding of Brøchner and of Kierkegaard to see the following statement of Brøchner's a few
weeks after Kierkegaard's death: ``On me his death has made a painful, but not any depressing
impression. He has been much to me, both by his writings and by the personal relation in which I
have stood to him. There is no one whose personality has to such a degree worked awakeningly and
spurringly on me, and the friendly benevolence he always showed me often gave me courage when I
was about to lose it. It is with joy that I remember him, in the long series of years I have
known him, almost 20 years, how the mild and loving in his personality constantly more got the
preponderance over the strong ironic and polemical element that was by nature in him, and how his
thought constantly became richer, surer, and clearer, so that often a word of his could work
calmingly and reconcilingly, and clear up a confusion that one could not with one's own forces
get the better of. I shall indeed come to miss him. But when I think of how completely he has
gotten accomplished that which was his life's task --- how rich and full this life has been in
its short duration, and how much of him remains --- I cannot with any oppressive feeling think of
his death; on the contrary, it appears to me beautiful and happy.'' (Letter from Brøchner to
Molbech, 2 Dec.\ 1855.)\footnote{Among Brøchner's papers there was found a little atmospheric
essay: \emph{Erindringer om Søren Kierkegaard} [Recollections of Søren Kierkegaard], which I
published in the periodical \emph{Det nittende Aarhundrede} [The Nineteenth Century] (1876).} ---
In order to understand this relation between Kierkegaard and Brøchner despite their so highly
different views, one must remember that for Kierkegaard personal truth and sincerity were the
first and the decisive thing. In his younger friend he found the striving after truth that
Lessing had declared to be the highest, and for Kierkegaard the dogmatic formulation of the found
truth got constantly less and less significance. It has been a purely Socratic influence that was
here exercised by the elder on the younger. But it has of course also been of great significance
that they both stood in a polemical relation to theology and the existing Church. As we have seen
in an earlier place, Kierkegaard maintained that ``the last formation of freethinkers (Feuerbach
and what belongs thereto)'' at bottom defended Christianity against the modern Christians. --- The
relation between Kierkegaard and Brøchner has been of an essentially different kind than the
relation between Kierkegaard and Nielsen.

It was only late in his life that Brøchner proceeded to a presentation of his philosophical view
in its totality. While his view of life was shaped practically and in silence under his life's
dispensations, he immersed himself in the study of the history of philosophy, which he really
founded here at home. In a thorough monograph on \emph{Spinoza} (1857) he gave a presentation of
the philosophical standpoint with which he felt himself most nearly akin. In his work \emph{Om
Udviklingsgangen i Filosofiens Historie} [On the Course of Development in the History of
Philosophy] (1869) he gave a strongly condensed presentation of the fundamental thoughts that
have reigned in philosophy, of their relation to other sides of cultural development, and of the
historical connection that, consciously or unconsciously, makes itself felt among them. This
historical study found its conclusion in a handbook of \emph{Filosofiens Historie} [The History
of Philosophy] (1872--74), the only Danish presentation that comprises both Greek and modern
philosophy.

His most important writings are the two books \emph{Problemet om Tro og Viden} [The Problem of
Faith and Knowledge] (1868), which gives a critique of Martensen, Kierkegaard, and Nielsen, and
\emph{Det Religiøse i dets Enhed med det Humane} [The Religious in Its Unity with the Human]
(1869), which develops his own religious-philosophical view.

\begin{center}---\end{center}

Brøchner's starting point is psychological, and is akin to the one that, as measure, lies at the
foundation in Kierkegaard's doctrine of the stages: the human spirit as energy and totality. The
view of life that is to be able to subsist must be such a one that makes possible an unfolding of
human nature, so that the various elements in it can work harmoniously together. In the concept
of self there lies both energy (activity) and wholeness. In self-consciousness both make
themselves known in the urge to form and harmonize impulses and drives, as well as to order all
thoughts in inner connection and wholeness --- the urge that is man's essential urge. Man
distinguishes himself from the animal by this urge to go out beyond the individual impulses and
the individual sensations and reach the universally valid, the universal.

Brøchner uses the expression essence-totality to designate his chief point of view. Such a
totality lies, according to his conception, behind all psychic expressions of life, even if it
does not directly and positively come to consciousness. Hereby his psychology gets a dogmatic
character. The essence-totality is not grounded on analysis of experience, but is set up in
advance as indubitable. Over against the thinkers with whom he most nearly negotiated, he could
also very well assume presuppositions about the human spirit's essence as totality as given,
since they each in his way acknowledged it; but it would have served for clarity and more
definite decision if a closer grounding had been given. And apart from the lack of sufficient
experiential grounding, there is also this doubtfulness in Brøchner's principle, that he too much
uses it to explain individual psychological phenomena. Thus he teaches that the drive ``with
inner necessity'' develops from being the will's rudimentary form to becoming a true will. Now
experience shows, however, that not all drives develop thus, and that when it happens, it is
because other psychological phenomena, especially other drives, appear, and it is through
interactions and conflicts with them that the development in the hinted-at direction can take
place. A whole quantity of psychological relations --- contrast-effects, associations, and
displacements --- can hereby make themselves felt. For Brøchner's speculative psychology, on the
other hand, that ``inner necessity'' springs from the ``essence-totality,'' in that under the
presupposition of this it becomes self-contradictory to stop at a single impulse in its
isolation. --- In a corresponding way the new element that by the decision comes to the preceding
unfinished cognition (deliberation) is to be a supplementation stemming from that essence-totality
as ``immediate source.''\footnote{The mutual relation that here shows itself between the conscious
and the unconscious Brøchner has discussed more closely in an essay in \emph{Nyt dansk
Maanedsskrift} [New Danish Monthly] (1871).} In all practical sense the same makes itself known,
for example when one acts in opposition to one-sided theories. In sympathetic impulses there
expresses itself an urge for wholeness that has not yet come to full ethical clarity. The true
acting is that in which man develops and confirms himself through the devotion to the object of
the action. And this is only possible when this object is not something accidental, finite, and
perishable; it is a contradiction to nourish an infinite interest for such a thing. With the
finitude in the aim all egoism falls away. In a similar way, cognition is first completed when its
content for it itself no longer stands merely as given and received, but is formed into a unity
that corresponds to the unity that makes itself known in self-consciousness. As a self-energy
man's essence makes itself known both where it works in the form of cognition and where it works
in the form of will. Of these two forms cognition must constantly keep the primacy, since the
special activities of spirit are constantly accompanied by consciousness and only thereby are
activities of spirit. For the rest, there is a constant mutual relation between cognition and
willing: there is a passion of cognition, and there is a willing of the cognized universally
valid. Where in theory or in practice one stops at something finite and limited with an
inclination to make it into an absolute, the inner essence-totality will sooner or later lead out
beyond the limitation. Even in the evil, that is, the willing that uses all its energy to go up
into something finite, there is an untenable opposition between the will's form and its content.
The development and exercise of force here applied is not lost, but steps, through repentance,
into the good's service. Under the transformation of force here, as in the physical domain, the
sum of force is constant. And by means of the inner contradiction that is in the evil will and
expresses itself in restless unrest and striving, either as tumbling about in enjoyment or defiant
struggle against the good, the evil must more and more hollow itself out, so that a change of
direction becomes the only way out.

The ``essence-totality'' or man's idea strives under all these forms to unfold itself, however
hard the struggle may become. It is Brøchner's fundamental thought, on essential points a
deepening --- taking place under Kierkegaard's influence --- of Hegel's teaching, by which this
was led back to the psychological and ethical points of view from which it for a large part
originally sprang. It is by means of this fundamental thought that he, although he does not
believe in the intervention of supernatural forces, holds fast the belief in ethical development
and progress.

\begin{center}---\end{center}

But when Brøchner assumes ``an inner necessity'' in the development and assumes an infinite
fullness of possibilities that can be realized through this development, he is clear that he
cannot herein support himself on strict demonstration. If one sets knowledge and faith in
opposition to each other, he is in no doubt that he will choose knowledge, because it is by
cognition, by thought, that man's essence alone can unfold itself in such a way that it deserves
the name of spirit. But only if one by faith understands authority-faith or belief in
supernatural forces and revelations does that opposition exist for Brøchner.

At no stage of development can a finite being fully realize the infinite content, the
essence-totality, that is its innermost ground. Often it must anticipate what can first later be
demonstrated, and often it must stop at propositions that are set up without their presuppositions
being able to be brought forward; these presuppositions lie hidden in its innermost essence.
Experience is constantly incomplete, and therefore cognition can only be concluded by a postulate,
a faith. But since this faith is knowledge's self-supplementation by virtue of the same principle
according to which there exists a knowledge at all, there becomes here no necessary strife between
faith and knowledge. The principle according to which there exists a knowledge is a principle of
unity and totality. Faith aims at the point that such a principle has validity in the whole of
existence: ``The unity-principle of the existence not exhausted by experience is grasped not as
that whose certainty experience secures, but as that which is reached by a necessary postulate
proceeding from the individual consciousness, which without it would lose its hold in itself,
become uncertain of its own unity, be confused in itself.'' It is a rational faith, in that it
preserves the agreement with the rational cognition. There is a devotion in faith; but neither in
cognition is this element lacking: ``In a true cognition there is a purification of thought, a
giving-up of individual entanglement, opinion, and prejudice --- in a word, of thought's
selfishness; there is in it a self-devotion to the truth's absolute power.'' In faith this element
only steps forth more strongly.

It is by this way that one reaches the concept of the Absolute, of ``that which rests in itself
and has certainty in itself,''\footnote{It is Spinoza's definition of the absolute substance that
here lies at the foundation: ``By substance I understand that which is in itself and is conceived
through itself, that is, whose concept does not presuppose the concept of any other thing from
which it should be formed.'' Spinoza: \emph{Ethica}, I, Def.~3.} man's highest and last thought.
This thought we approach by analysis of existence, as experience shows it to us, in that
interaction and causal connection comprise everything we know, and presuppose an energy that stirs
in everything, an infinite totality within which the finite totalities that experience displays
find their explanation. Such an analysis leads only indirectly toward the last presupposition.
Directly we approach it when we proceed from the highest forms of existence, from life, which is
inexplicable from the merely material, but especially from the psychic and the spiritual. These
phenomena make it necessary to assume an absolute subjectivity (spirit), of which the
subjectivities appearing in experience are only limited forms. --- Brøchner prefers the more
abstract expression subjectivity (spirit) to the more concrete expression personality, since this
expression can be used only of a finite being that stands in relation to other beings. He is a
pantheist, not a theist. ``When we then do not limit the concept of pantheism to the imperfect
and one-sided forms of it, then we must acknowledge the pantheistic conception as essential and
necessary for philosophy\ldots\ No philosophy can think the Infinite as being separated from and
outside the Finite\ldots\ All philosophy seeks the unity and connection of cognition and is only
philosophy in so far as it pursues this\ldots\ Every higher thought to which it rises is precisely
the higher only thereby, that it comprises the lower.'' --- As little as other Hegelians will
Brøchner admit that in his concluding thoughts he supports himself on an analogy. Analogy is in
reality the source of all metaphysical and religious representations.\footnote{Compare my
\emph{Religionsfilosofi} II C and my \emph{Mindre Arbejder} [Minor Works], Second Series,
pp.~41--46.} It really lay near to Brøchner to make this admission, as soon as he separated
himself from Hegel by the relative difference between knowledge and faith, a difference that Hegel
would not acknowledge in this form. But the admission would have brought with it the acknowledgment
of a greater co-operation of other psychic elements than the intellectual in the arising and
development of those representations.

Over against the positive religions, it is Brøchner's most essential objection that they stop at
finite representations. The Divine appears in the form of finite personalities, and therefore with
sensuous qualities; it is limited in time and space, and it works (in the miracle) from without on
nature, without law and inner connection. Therewith the egoistic in the positive religions is
connected; man relates himself in them as a finite being to another, an I to a Thou. Real
self-devotion without egoism is only possible over against the Infinite, and only such a
self-devotion is one with the true self-confirmation, in that new force is released through this
relation, in which ``the sources of the being spring forth from the deep of the Infinite.'' ---
Constantly there lies at the foundation in Brøchner the presupposition that there is an infinite
content, a fund of force, that shall and can be unfolded or released, and that the most inward
self-recollection is, as regards man, the means to get a share in it. His line of thought here
gets the character of a mysticism, and he would surely not himself quite deny this kinship. In the
inward self-recollection we discover that we are members of an infinite connection and discover a
law that comprises both ourselves and the whole rest of existence. ---

Over against Feuerbach, Brøchner maintains that in all religions, especially the higher ones,
there are found rational elements, and that one must not merely fix one's gaze on the individual
religions each for itself, but look at the course of development in the history of religion, from
the lowest forms of nature-religion to the highest of the spirit-religions. It is a fact that it
has been possible for the human spirit to work its way up from lower to higher forms of religion.
And this has not happened by the properly religious disappearing, but precisely by its stepping
forth more --- when one by religion understands man's absolute relation to the Absolute as to a
principle that grounds and comprehends the human as well as the whole rest of existence. It is a
relation of unity, not a relation of opposition. The religious expresses itself in every positive
form of spiritual life. --- In such a religion, the religion of humanity, the essential thing in
the positive forms of religion is preserved, while the historical forms are dissolved. It is
reality's true and eternal religion that is present, though often unclearly and unconsciously, in
our time's human consciousness. It is not metaphysics that is set in religion's place; the
abstract concepts are only to unravel what stirs in the humane-religious reality of life.

Brøchner himself calls attention to the fact that the religious, according to his determination,
corresponds both to Kierkegaard's ethical stage and to the religious stage that is not yet the
Christian. He could here have taken the intermediate stage, ``humour,'' along too. These three
stages go for Brøchner into one; in particular he will not set the religious in opposition to the
ethical, since it is for him precisely the soul in this. --- It is connected with his strongly
concentrated presentation that he does not find occasion to hold fast Kierkegaard's distinctions,
or himself to set up others. Here too a psychological foundation is lacking, and the attempt to
give such a one would surely have led to bringing forward difficulties and questions that now
could not rightly come forth.

By his abstract language Brøchner has cut himself off from becoming accessible to larger circles;
but he has the merit that he in his way has maintained the connection in the personal life, which
was volatilized by speculative philosophy and burst by Kierkegaard. And in opposition to every
more or less materialist attempt to build the human view of life exclusively on what experience
teaches about material nature, he has maintained thought's reality as that whose calling it is to
interpret all forms of existence, the material as well as the spiritual. ---

In a lecture that Brøchner delivered in the year that was to become his year of death, he stated
that he was an adherent of ``the idealist philosophy that Germany has the honour of having brought
forth.'' We have seen how this philosophical tendency in him himself, as well as in the other
thinkers who had been subjected to a similar influence, broke with the Danish interest in
psychology, ethics, and criticism. If in the philosophical generation that followed after Brøchner
and his contemporaries these interests have come forth more strongly than before, that is due, to
be sure, in part to the influence from English and French thought that made itself felt here at
home in the last part of the nineteenth century; but it may also be conceived as a continuation of
the specially Danish tendency in the predecessors --- right from Holberg's time.

\begin{center}---\end{center}

\end{document}
